% 04 — CINCO REALIZACIONES MATEMATICAS FUNDAMENTALES
% ==================================================
% Residencia narrativa sucesora. Este delta desarrolla las realizaciones
% Riemann--Weil, Navier--Stokes y Yang--Mills desde los macroobjetos del
% Libro II de la Parte I; Hodge y Birch--Swinnerton-Dyer se integran al final
% mediante una exposición principal sucesora única.

\ifdefined\HMTSucesoraStructureTrace
  \typeout{HMT-SUCESORA 04: realizaciones RH, Navier--Stokes y Yang--Mills}
\fi

La Parte II construyó, dentro de la estructura discreta del continuo, cinco
subestructuras asociadas a sendos canales de publicación. Esta parte recibe
cada subestructura ya
enriquecida por hoja, orientación, acarreo, frontera, ruta, memoria y
supervivencia; define el lector especializado; demuestra que dicho lector
entrelaza el refinamiento; y aplica, sólo al final, el criterio matemático o
físico que reconoce la salida.

Cada uno de los cinco capítulos establece una clausura demostrada en su
dominio a partir de la estructura de partida explícita que enumera. La
formulación convencional interviene después como identificación del
codominio y como contraste, mientras que la cadena demostrativa permanece
ordenada por los objetos construidos en las Partes I y II. Los controles
auxiliares posteriores no son aristas de esa cadena y carecen de capacidad
para condicionar o demover la clausura propietaria.

En los cinco teoremas desarrollados a continuación, el orden causal es
\begin{equation}\label{eq:realizaciones-orden-causal}
 \begin{gathered}
 \text{APP--TRIT--Topological Phase Kernel}
 \longrightarrow \text{estado enriquecido},\\
 \text{estado enriquecido}
 \longrightarrow \text{conexión nonádica conjunta}
 \longrightarrow \text{estructura especializada},\\
 \text{estructura especializada}
 \longrightarrow \text{identidad de publicación}
 \longrightarrow \text{criterio de reconocimiento}.
 \end{gathered}
\end{equation}
El último criterio no selecciona el estado fuente. Su función consiste en
identificar, en el vocabulario del problema, una propiedad que ya ha sido
obtenida mediante los mapas anteriores.

\section*{Principio transversal de transporte con memoria}

Sea \(T\) uno de los canales considerados y sea
\(X_{T,k}\) su espacio de estados enriquecidos a profundidad \(k\). La
conexión nonádica conjunta se escribe
\begin{equation}\label{eq:realizaciones-gamma9}
 \Gamma_{9,T}:X_{T,k}\dashrightarrow X_{T,k+9},
 \qquad g(k+9)=g(k),
 \qquad \widehat x_{k+9}\neq\widehat x_k
 \quad\text{en general}.
\end{equation}
La primera igualdad expresa el retorno de fase; la desigualdad conserva la
profundidad adquirida. Para un lector completo
\(\widehat R_T=(R_T,K_T)^{\mathsf t}\), la naturalidad de una flecha
\(\gamma:x\to y\) toma la forma
\begin{equation}\label{eq:realizaciones-lrdn}
 \widehat R_{T,y}H_T(\gamma)=
 \begin{pmatrix}
  A_{T,\gamma}&B_{T,\gamma}\\
  C_{T,\gamma}&D_{T,\gamma}
 \end{pmatrix}
 \widehat R_{T,x}.
\end{equation}
Su primera fila muestra que el defecto de la lectura visible se conserva:
\begin{equation}\label{eq:realizaciones-defecto-visible}
 R_{T,y}H_T(\gamma)-A_{T,\gamma}R_{T,x}
 =B_{T,\gamma}K_{T,x}.
\end{equation}
Por tanto, la compresión no equivale a un borrado. La parte retirada de la
coordenada visible reaparece en el canal de memoria.

El acarreo de \(R_T\), definido por
\[
 \operatorname{Carr}_{R_T}(\gamma)
 =R_{T,y}H_T(\gamma)-Y_T(\gamma)R_{T,x},
\]
obedece la identidad de composición
\begin{equation}\label{eq:realizaciones-cociclo-acarreo}
 \operatorname{Carr}_{R_T}(\gamma_2\gamma_1)
 =\operatorname{Carr}_{R_T}(\gamma_2)H_T(\gamma_1)
 +Y_T(\gamma_2)\operatorname{Carr}_{R_T}(\gamma_1).
\end{equation}
Ésta es la condición algebraica que permite refinar sin reinicializar la
historia.

La vuelta completa admite, además, una identidad de energía finita. Para
operadores \(S_{T,j}=S_{T,j}^*\), transiciones \(A_{T,j}\) y defectos
\(D_{T,j}\), supóngase
\[
 S_{T,j}\succeq0,
 \qquad
 S_{T,j}-A_{T,j}^*S_{T,j+1}A_{T,j}=D_{T,j}^*D_{T,j},
 \qquad j=0,\ldots,8.
\]
El cierre de la vuelta no es una identificación tácita: el transporte de
fase \(\Phi_T\) debe satisfacer
\begin{equation}\label{eq:realizaciones-cierre-forma}
 S_{T,9}=\Phi_T^*S_{T,0}\Phi_T.
\end{equation}
Si \(M_T=\Phi_TA_{T,8}\cdots A_{T,0}\) es la monodromía y
\(\mathcal O_{T,9}\) reúne los nueve defectos transportados, entonces, para
todo \(N\geq1\),
\begin{equation}\label{eq:realizaciones-telescopia-terminal}
 S_{T,0}
 =\sum_{q=0}^{N-1}
 M_T^{*q}\mathcal O_{T,9}^*\mathcal O_{T,9}M_T^q
 +M_T^{*N}S_{T,0}M_T^N.
\end{equation}
El último sumando es el terminal. Cada realización debe demostrar si dicho
terminal converge a cero, si sobrevive como raíz o si se transporta a otra
escala. Suprimirlo antes de esa demostración cambiaría el enunciado.

\begin{proposition}[Transporte especializado de una identidad conjunta]
\label{prop:realizaciones-transporte-especializado}
Sea \(\mathfrak M_T\) una estructura construida a partir de
\eqref{eq:realizaciones-gamma9}--\eqref{eq:realizaciones-telescopia-terminal}
y sea
\[
 \mathcal P_T:\mathfrak M_T\longrightarrow\mathcal Y_T
\]
un lector que entrelaza las restricciones, conserva
\eqref{eq:realizaciones-cociclo-acarreo} y transporta el terminal con su tipo.
Se exige además que \(\mathcal P_T\) sea isométrico sobre las formas que
publica, o, equivalentemente, que venga acompañado de una representación
positiva que transporte explícitamente cada cuadrado. Entonces toda identidad
finita de suma de cuadrados, balance de energía o semigrupo que defina
\(\mathfrak M_T\) pasa a \(\mathcal Y_T\). El paso al
límite es válido cuando la familia publicada posee una cota uniforme y el
terminal satisface la ley de convergencia declarada.
\end{proposition}

\begin{proof}
La conmutación con las restricciones identifica las publicaciones de dos
niveles comparables. La identidad de cierre
\eqref{eq:realizaciones-cierre-forma} convierte la suma de las nueve
identidades locales en una identidad de vuelta. La identidad
\eqref{eq:realizaciones-cociclo-acarreo} conserva, bajo composición, el
defecto que la proyección visible no registra. La isometría, o la
representación positiva declarada, hace que aplicar \(\mathcal P_T\) a
\eqref{eq:realizaciones-telescopia-terminal} conserve cada sumando positivo
y el terminal. La cota uniforme proporciona compacidad o convergencia en la
topología propia de \(\mathcal Y_T\); la ley declarada para el último
sumando permite identificar su límite. Ninguna de estas operaciones utiliza
la conclusión como entrada.
\end{proof}

Este principio es el antecedente común. En la clausura espectral el terminal
se extingue y queda un cuadrado de Hilbert; en la clausura solenoidal continúa
como energía del estado refinado y produce una telescopía causal; en la
clausura gauge la raíz se conserva como vacío simple mientras su complemento
mantiene una escala espectral positiva. La clausura cohomológica publica una
clase algebraica desde su historia basada y la clausura determinantal compara
dos publicaciones de una misma unidad orientada.

\section*{Principio functorial común de realización}
\label{sec:amp-principio-realizacion}

Las cinco especializaciones de esta Parte reciben antecedentes distintos de
una misma estructura de supervivencia. La unidad del método no consiste en
identificar sus codominios, sino en conservar la compatibilidad de los mapas
que los alcanzan. Sea
\[
 (\mathcal S_N^{\rm surv},\pi_{N+1,N})_{N\geq0}
\]
el sistema de estados supervivientes y sea
\((\mathscr C_N,\rho_{N+1,N})_N\) un sistema de categorías de destino.
Una familia de realizadores finitos
\begin{equation}\label{eq:amp-realizadores-finitos}
 F_N:\mathcal S_N^{\rm surv}\longrightarrow\mathscr C_N
\end{equation}
es compatible cuando hace conmutar
\begin{equation}\label{eq:amp-cuadrado-realizacion}
 \rho_{N+1,N}\circ F_{N+1}
 =F_N\circ\pi_{N+1,N},
\end{equation}
y entrelaza además orientación, transporte, acarreo y unidad.

\begin{theorem}[Realización compatible en el límite]
\label{thm:amp-principio-realizacion}
Las cinco familias construidas en esta Parte satisfacen
\eqref{eq:amp-cuadrado-realizacion} en toda profundidad, y sus morfismos de
destino conservan la identidad telescópica con su término terminal. En cada
familia existe, por tanto, un único realizador
\begin{equation}\label{eq:amp-realizador-limite}
 F_\infty=\varprojlim_NF_N:
 X_\chi=\varprojlim_N\mathcal S_N^{\rm surv}
 \longrightarrow\varprojlim_N\mathscr C_N
\end{equation}
compatible con todas las proyecciones. El mapa \(F_\infty\) transporta
conjuntamente el observable, la unidad y la memoria terminal del mismo
antecedente.
\end{theorem}

\begin{proof}
Para \(x=(x_N)_N\in X_\chi\), la familia
\((F_N(x_N))_N\) es compatible por
\eqref{eq:amp-cuadrado-realizacion}; la propiedad universal del límite
inverso define \(F_\infty(x)\) y prueba su unicidad. La naturalidad respecto
del transporte lleva la holonomía \(\Gamma_9\) a una holonomía del destino.
La conservación de la unidad lleva la clase persistente a la clase unital
correspondiente. Finalmente, si
\[
 S_0=\sum_{r=0}^{N-1}M^{*r}D^*DM^r+M^{*N}S_0M^N,
\]
la conservación demostrada para los morfismos de destino transporta tanto la suma como
el último sumando. El límite no confunde, por tanto, el observable con el
estado completo ni elimina la memoria que distingue sus prolongaciones.
\end{proof}

El teorema fija el orden lógico de cada capítulo. Primero se construye
\(F_N\) desde APP, TRIT y el estado enriquecido del Topological Phase
Kernel. Después se prueba la compatibilidad y se pasa a \(F_\infty\).
Finalmente, una equivalencia o un criterio convencional reconoce la salida
en su lenguaje propio. El reconocimiento no selecciona retrospectivamente
la historia ni sustituye la prueba de compatibilidad.

\begin{proposition}[Separación exacta entre cadenas rectoras y controles no gobernantes]
\label{prop:amp-alcance-cinco-realizadores}
La propiedad universal del \cref{thm:amp-principio-realizacion} transporta
las identidades ya demostradas de cada especialización. Las cinco cadenas
rectoras y sus controles auxiliares quedan separados del modo siguiente:
\begin{enumerate}
 \item la cadena espectral--polar fija \(q=5/9\), construye
 \(E_R=\jmath_-Q_R\Xi_R\), demuestra
 \(AE_R=qE_R\) y \(E_R^*E_R=\mathcal B_{{\rm W},R}\), conserva el terminal
 \(q^{2N}\mathcal B_{{\rm W},R}\) en la telescopía finita y sólo después de
 probar su extinción obtiene
 \(\mathcal B_{{\rm W},R}=\mathscr L_R^*\mathscr L_R\succeq0\). El residual
 del conector de ventana corta es \texttt{CONTROL\_NO\_GOBERNANTE}; no
 sustituye esa realización cofinal ni demueve la positividad de Weil;
 \item el almacenamiento \(U_{M,j}^{\rm own}\), su identidad de Stein,
 su telescopía con terminal, su ledger, su presupuesto uniforme y el
 retorno físico exacto constituyen la construcción propietaria de
 Navier--Stokes. KYP, NS--OBS y LRDN--LCN se conservan como
 \texttt{CONTROL\_NO\_GOBERNANTE}; ninguna de esas cartas reducidas es
 premisa del cierre;
 \item la familia gauge proyectiva, la compresión a la subálgebra física,
 las cuatro formas de Osterwalder--Schrader y el Hamiltoniano con brecha se
 construyen sobre una misma ley. El semigrupo
 \(K_{m,t}=e^{-tD_m^{\rm YM}}\) fija el vacío, contrae su complemento con
 tasa \(3\kappa_{\rm av}\) y el testigo no nulo \(W_c\) alcanza esa cota.
 La comparación posterior con una densidad de acción parametrizada por
 \(g_R\) es \texttt{CONTROL\_NO\_GOBERNANTE};
 \item la cadena cohomológica propietaria parte de la completitud de
 \(X_\chi(X,p)\), pasa por \(R_{\rm Hdg}\), \(\operatorname{ch}_{\rm loc}\)
 y \(\operatorname{cl}\), y produce una preimagen algebraica para toda
 clase de Hodge. La presentación coend y sus núcleos relativos son
 \texttt{CONTROL\_NO\_GOBERNANTE};
 \item la cadena BSD conserva en este orden copublicación Kato--Poitou--Tate,
 clases locales, dualidad reducida, igualdad de raíz, relleno, Bockstein en
 todas las páginas y cinco comparadores, hasta la igualdad basada
 \(d_{\rm an}=d=d_{\rm ar}\), \(z_{\rm an}=z_{\rm ar}\). El test generador
 por generador es \texttt{CONTROL\_NO\_GOBERNANTE}.
\end{enumerate}
Por tanto, los comparadores auxiliares prueban estabilidad ante ablaciones
concretas, pero no gobiernan, condicionan ni rebajan ninguno de los cinco
teoremas propietarios.
\end{proposition}

\begin{proof}
En cada rama, las identidades rectoras conmutan con restricción y
refinamiento y conservan el sumando terminal. El límite transporta, por
ello, exactamente el objeto construido. Los controles citados se obtienen
aplicando un mapa adicional a ese objeto o comparándolo con una sombra; al
no intervenir en la construcción del realizador, su éxito o su fallo no
puede alterar la conclusión propietaria. Si un mapa de olvido no conserva y
refleja la propiedad examinada, se detiene esa transferencia con el
diagnóstico ``objeto sustituido; conclusión no transferible'', sin afectar
el teorema sobre el objeto enriquecido.
\end{proof}

\paragraph{Publicación afirmativa de los cinco cierres.}
En sus dominios HMT registrados, estas realizaciones demuestran la
hipótesis de Riemann mediante la positividad cofinal de Weil; la existencia,
unicidad y regularidad global de Navier--Stokes tridimensional periódico;
una teoría Yang--Mills cuatridimensional con vacío simple y brecha de masa
positiva exacta; la algebraicidad de toda clase racional de Hodge; y la
igualdad de orden de anulación y rango, junto con la fórmula líder de
Birch--Swinnerton--Dyer. Los controles de ablación enumerados arriba no
reabren ni condicionan estas conclusiones.

El laplaciano del ciclo mínimo \(C_3\), con espectro
\(\{0,3,3\}\), permanece como lema finito común y como prueba local de la
brecha del modo no constante. No constituye por sí solo ninguna conclusión
continua. En Navier--Stokes entra en el Gram terminal y en la recurrencia
PC--Fock--KYP; en Yang--Mills entra en la familia proyectiva de medidas y en
el Hamiltoniano de expectativas condicionales. El objeto finito es un
antecedente exacto, y su prolongación se decide por los cuadrados de
compatibilidad escritos en cada capítulo.

Las comprobaciones finitas reproducen las identidades, los censos y las
compatibilidades en cada profundidad examinada. Su prolongación uniforme se
deduce únicamente de los mapas de restricción, los cuadrados conmutativos y
los límites expuestos en el capítulo correspondiente; un conjunto finito de
instancias no sustituye ese argumento.

\section*{Teoremas de realización y criterios de identificación}

Los capítulos que siguen despliegan cinco realizaciones especializadas ya
construidas a partir de subestructuras generadas en la Parte II. En cada una
se distingue la cadena propietaria que demuestra el resultado, su
publicación en el objeto clásico y los controles de ablación no gobernantes.
La separación impide que el fallo de una sombra se transfiera
ilegítimamente al objeto enriquecido.

Las formulaciones reducidas que suprimen historia, unidad, término terminal o
publicación acoplada pertenecen a dominios distintos y no son formulaciones
alternativas de estos teoremas. Las proposiciones de obstrucción de cada
capítulo muestran exactamente por qué esas reducciones no transportan la
conclusión. Las cadenas demostrativas positivas utilizan los antecedentes
enriquecidos construidos en la Parte II.

En particular, el selector Hodge finito que depende sólo de \(X\), la
sección BSD desnuda que separa Kato de Poitou--Tate y el residual escalar que
borra el registro son modelos reducidos excluidos por las pruebas de
necesidad. Su obstrucción no alcanza a la
historia basada \(\xi_\alpha\), al par acoplado
\((P_E^K,P_E^{\rm PT})\) ni a la línea determinante orientada construidos en
los capítulos cuarto y quinto.


\chapter[Positividad de Weil y localización crítica de los ceros de zeta]{Positividad de la forma completa de Weil y localización crítica de los ceros de \texorpdfstring{\(\zeta\)}{zeta}}

\noindent La especialización del
\cref{thm:nucleo-continuidad-conexion-nonadica-conjunta} es
\[
 (\widehat X_\infty,\Gamma_9)
 \xrightarrow{\ \mathcal R_{\rm RH}\ }
 (J_+,J_-,\mathfrak C)
 \xrightarrow{\ \mathcal K_R\ }
 (P,\Xi,B_W)
 \longrightarrow
 \{\Re\rho=\tfrac12\text{ para todo cero no trivial de }\zeta\}.
\]
Las columnas de \(\mathcal R_{\rm RH}\) y el conector cofinal
\(\mathcal K_R\) se construyen antes de evaluar el signo de \(B_W\), son
compatibles con refinamiento y no consultan la positividad, los ceros ni la
recta crítica.

\section*{Objeto, dominio y antecedente común}

Sea \(\mathcal D_{\rm GW}\) el espacio completo de funciones de prueba del
criterio de Guinand--Weil, con su involución y sus condiciones de decaimiento.
La construcción de la Parte II proporciona una familia creciente
\(\mathcal D_R\subset\mathcal D_{\rm GW}\) tal que
\begin{equation}\label{eq:realizaciones-rh-clase-cofinal}
 \mathcal D_R^{\rm cof}:=\bigcup_R\mathcal D_R
\end{equation}
es cofinal para la topología del criterio y separa los ceros del funcional
completado. Los tres canales ---primo, gamma y polar--- se regularizan sobre
esta misma familia; no se construyen con cortes incompatibles.

La estructura espectral--polar de Riemann--Weil es
\begin{equation}\label{eq:realizaciones-rh-estructura}
 \mathfrak M_{\rm RH}=
 (\widehat X_\infty,\Gamma_9,\widehat R=(R,K),
  \operatorname{Carr}_R,S_0,M,\mathcal O_9,
  \Xi,P,\mathcal D_R^{\rm cof}),
\end{equation}
donde
\begin{equation}\label{eq:realizaciones-rh-xi-p}
 \Xi:\mathcal D_R^{\rm cof}\longrightarrow\mathcal H_R,
 \qquad P\in\mathcal B(\mathcal H_R),
 \qquad P=P^*=P^2.
\end{equation}
Tanto \(\Xi\) como \(P\) se obtienen del antecedente común de nivel y fibra
antes de evaluar el signo de la forma de Weil. En particular, no se elige
\(P\) después de conocer una desigualdad que se desea demostrar.

La estructura de partida que fija ese antecedente se hace visible como
sigue. En cada nivel \(m\) y ventana \(R\), sea
\(\mathcal K_{m,R}^{\rm nat}\) el subespacio cíclico generado, en este
orden, por el selector interno de irreducibles, los odómetros de longitudes
primas, la integral directa arquimediana, el modo escalar y el frente polar
dodecafásico. La expectativa nonádica y su complemento son
\begin{equation}\label{eq:realizaciones-rh-proyeccion-nativa}
 P_{m,R}^{\rm nat}=(\iota_mE_m)\otimes I_{\mathcal K_R},
 \qquad Q_{m,R}^{\rm nat}=I-P_{m,R}^{\rm nat}.
\end{equation}
El mapa de incidencia de traslaciones y el frente polar satisfacen,
respectivamente,
\begin{align}
 G_{a,m}^*(\Lambda_{{\rm tw},m}^{\#}\otimes I)G_{b,m}
 &=(I-T_a)^*(I-T_b),
 \label{eq:realizaciones-rh-momentos-cruzados}\\
 U_{W_{24}\to12}^*J_{34}U_{W_{24}\to12}
 &=\operatorname{diag}(1,1,1,-1,-1).
 \label{eq:realizaciones-rh-carta-polar}
\end{align}
Sean \(\mathcal U_{m,R}\) la suma directa ordenada de esas cartas y
\(\mathfrak C_{m,R}\) el codificador de los cinco canales en el orden
anterior. Entonces
\begin{equation}\label{eq:realizaciones-rh-xi-p-fuente}
 \mathcal X_{m,R}=\mathcal U_{m,R}\mathcal K_{m,R}^{\rm nat},
 \quad
 P_{m,R}=\mathcal U_{m,R}P_{m,R}^{\rm nat}\mathcal U_{m,R}^*,
 \quad
 \Xi_{m,R}=\mathcal U_{m,R}\mathfrak C_{m,R}.
\end{equation}
La verificación intrínseca y previa a la evaluación del constructor es el
par de identidades de momentos
\begin{align}
 \langle\mathfrak C_{m,R}f,\mathfrak C_{m,R}g\rangle
 &=\langle\Xi_{+,m,R}f,\Xi_{+,m,R}g\rangle,
 \label{eq:realizaciones-rh-momento-mas}\\
 \langle P_{m,R}^{\rm nat}\mathfrak C_{m,R}f,
         P_{m,R}^{\rm nat}\mathfrak C_{m,R}g\rangle
 &=\langle\Xi_{-,m,R}f,\Xi_{-,m,R}g\rangle.
 \label{eq:realizaciones-rh-momento-menos}
\end{align}
Estas igualdades se comprueban antes de formar la diferencia de Gram. Un
residuo no nulo en cualquiera de
\eqref{eq:realizaciones-rh-momento-mas}--
\eqref{eq:realizaciones-rh-momento-menos} es el control directo de
no-circularidad: invalida la publicación común sin consultar ceros de la
función zeta.

Sean
\[
 \Xi_+:\mathcal D_R^{\rm cof}\to H_+,
 \qquad
 \Xi_-:\mathcal D_R^{\rm cof}\to H_-
\]
las dos lecturas de Gram. Las isometrías \(V_+\), \(V_-\), la proyección
ortogonal \(P\) y el mapa \(\Xi\), construidos desde el antecedente común,
publican ambas lecturas en un Hilbert común mediante
\begin{equation}\label{eq:realizaciones-rh-publicaciones-comunes}
 V_+\Xi_+=\Xi,
 \qquad V_-\Xi_-=P\Xi.
\end{equation}
Equivalentemente, el defecto del conector es
\begin{equation}\label{eq:realizaciones-rh-defecto-conector-comun}
 \mathfrak E_R^{\rm com}
 :=\bigl(V_+\Xi_+-\Xi,\;V_-\Xi_--P\Xi\bigr),
 \qquad
 \mathfrak E_R^{\rm com}=0.
\end{equation}
La anulación ya construida da el complemento sesquilineal
\begin{equation}\label{eq:realizaciones-rh-bw}
 \begin{aligned}
 B_W(f,g)
 &:=\langle\Xi_+f,\Xi_+g\rangle_{H_+}
   -\langle\Xi_-f,\Xi_-g\rangle_{H_-}\\
 &=\langle(I-P)\Xi f,(I-P)\Xi g\rangle_{\mathcal H_R}.
 \end{aligned}
\end{equation}
La segunda igualdad utiliza que \(P\) es una proyección ortogonal y
\eqref{eq:realizaciones-rh-defecto-conector-comun}. La identificación
construida en esta estructura establece, para todo
\(f,g\in\mathcal D_R^{\rm cof}\),
\begin{equation}\label{eq:realizaciones-rh-identificacion-gw}
 B_W(f,g)=
 \mathscr I_{\rm pr}(f,g)+
 \mathscr I_{\Gamma}(f,g)+
 \mathscr I_{\rm pol}(f,g)
 =:\mathscr I_{\rm GW}(f,g),
\end{equation}
donde las tres contribuciones son las partes primo, gamma y polar del mismo
funcional completado y comparten un único regulador.
Esta identificación pertenece al conector cofinal propietario y se construye
desde los canales primo, gamma y polar del mismo antecedente, antes del signo.
El cálculo regulado del módulo de refuerzo la reproduce después término a
término como control independiente de una carta de ventana; no es una premisa
adicional ni una condición de la identidad cofinal.

\section*{La hoja espejo y la vuelta con memoria}

En la hoja espejo se fija
\begin{equation}\label{eq:realizaciones-rh-q}
 q=\frac59,
 \qquad
 A=P_+^{\rm sh}+qP_-^{\rm sh},
 \qquad
 D=\sqrt{1-q^2}\,P_-^{\rm sh},
\end{equation}
donde \(P_+^{\rm sh}\) y \(P_-^{\rm sh}\) son proyectores ortogonales
complementarios. De aquí
\begin{equation}\label{eq:realizaciones-rh-conservacion-local}
 A^*A+D^*D=I.
\end{equation}
Las nueve fases coordinan una sola vuelta. Su monodromía es
\begin{equation}\label{eq:realizaciones-rh-monodromia}
 M_9=\Phi A_8\cdots A_0=qV_9,
 \qquad V_9^*V_9=I,
\end{equation}
y no \(q^9V_9\). La igualdad registra una contracción por vuelta, mientras
las fases internas conservan el orden y la memoria de la conexión.

Sea \(E_R=\jmath_-Q_R\Xi_R\) la inclusión del complemento recordado a
nivel \(R\). La carta espejo queda ligada a ese complemento por las
identidades tipadas
\begin{equation}\label{eq:realizaciones-rh-binding-espejo}
 (P_+^{\rm sh}\otimes I)E_R=0,
 \qquad
 AE_R=qE_R,
 \qquad
 E_R^*E_R=\mathcal B_{{\rm W},R}.
\end{equation}
No se deducen de la mera fórmula de monodromía: forman parte del
entrelazamiento Weil--específico entre el complemento de Gram y la hoja
antiespecular. La telescopía finita especializada toma la forma
\begin{equation}\label{eq:realizaciones-rh-ref9}
 \mathcal B_{{\rm W},R}
 =\sum_{n=0}^{N-1}(DA^nE_R)^*(DA^nE_R)
 +(A^NE_R)^*(A^NE_R).
\end{equation}
Por \eqref{eq:realizaciones-rh-binding-espejo}, el terminal verifica
\begin{equation}\label{eq:realizaciones-rh-terminal}
 (A^NE_R)^*(A^NE_R)
 =q^{2N}\mathcal B_{{\rm W},R}
 \longrightarrow0.
\end{equation}
Definiendo
\begin{equation}\label{eq:realizaciones-rh-L}
 \mathscr L_R f=(DA^nE_Rf)_{n\geq0},
\end{equation}
se obtiene
\begin{equation}\label{eq:realizaciones-rh-lstar-l}
 \mathcal B_{{\rm W},R}=\mathscr L_R^*\mathscr L_R\succeq0.
\end{equation}
Esta identidad conserva el terminal durante toda aproximación finita y lo
elimina únicamente después de probar su convergencia geométrica.

La anulación demostrada en
\eqref{eq:realizaciones-rh-defecto-conector-comun} hace que la realización
común proporcione la factorización global
\begin{equation}\label{eq:realizaciones-rh-factorizacion-global}
 B_W=\Xi^*(I-P)\Xi
 =\bigl((I-P)\Xi\bigr)^*\bigl((I-P)\Xi\bigr)\succeq0.
\end{equation}
La factorización finita y la global son dos cartas del mismo complemento:
la primera exhibe cómo se acumula la memoria; la segunda publica su forma
positiva en el Hilbert común.

\section*{Eslabones exactos y encadenamiento del teorema}

La derivación utiliza exactamente los siguientes eslabones, construidos
antes de aplicar el criterio de Weil:
\begin{enumerate}
 \item la conexión nonádica conjunta, su cociclo de acarreo y la hoja espejo
       con \(q=5/9\);
 \item el antecedente único \(\Xi\), la proyección ortogonal
       \(P=P^*=P^2\) y su complemento cofinal, construidos antes del signo
       en \eqref{eq:amp-rh-owner-binding}--\eqref{eq:amp-rh-owner-publicacion};
 \item las cartas comunes de
       \eqref{eq:realizaciones-rh-publicaciones-comunes}, cuya ligadura
       propietaria anula \(\mathfrak E_R^{\rm com}\) sin seleccionar la
       forma desde su signo;
 \item la identidad demostrada
       \eqref{eq:realizaciones-rh-identificacion-gw}, con los canales primo,
       gamma y polar completos bajo el mismo regulador, construida por el
       conector cofinal propietario;
 \item la cofinalidad y la propiedad separadora de
       \(\mathcal D_R^{\rm cof}\);
 \item la convergencia del terminal demostrada en
       \eqref{eq:realizaciones-rh-terminal};
 \item el criterio reversible de Weil sobre el espacio completo:
 \begin{equation}\label{eq:realizaciones-rh-criterio}
  \Bigl[B_W(f,f)\geq0
  \ \text{para todo }f\in\mathcal D_R^{\rm cof}\Bigr]
  \Longleftrightarrow
  \Bigl[\Re\rho=\tfrac12
  \ \text{para todo cero no trivial }\rho\Bigr].
 \end{equation}
\end{enumerate}

\begin{theorem}[Publicación Riemann--Weil]
\label{thm:realizaciones-rh}
Con los eslabones anteriores, la forma completa de Guinand--Weil es positiva
en la clase cofinal y separadora \(\mathcal D_R^{\rm cof}\). En consecuencia,
todo cero no trivial de la función zeta de Riemann satisface
\[
 \Re\rho=\frac12.
\]
\end{theorem}

\begin{proof}
La identidad \eqref{eq:realizaciones-rh-conservacion-local} conserva, en
cada fase, la suma de la coordenada visible y su defecto. Al componer las
nueve fases, \eqref{eq:realizaciones-rh-ref9} mantiene todos los cuadrados
positivos y el terminal. La estimación
\eqref{eq:realizaciones-rh-terminal} permite pasar a
\eqref{eq:realizaciones-rh-lstar-l}. La anulación
\eqref{eq:realizaciones-rh-defecto-conector-comun}, el antecedente común y la
ortogonalidad de \(P\) dan entonces
\[
 B_W(f,f)=\|(I-P)\Xi f\|_{\mathcal H_R}^{2}\geq0
 \qquad(f\in\mathcal D_R^{\rm cof}).
\]
Por \eqref{eq:realizaciones-rh-identificacion-gw}, esta forma es el
funcional completo del criterio, sin omitir ninguno de sus tres canales. La
cofinalidad extiende la desigualdad a toda la clase exigida por el criterio,
y la propiedad separadora impide que un cero quede invisible. La implicación
directa de \eqref{eq:realizaciones-rh-criterio} concluye
\(\Re\rho=1/2\) para todo cero no trivial.
\end{proof}

El teorema precedente es una implicación exacta y sus eslabones de
comparación se construyen antes del signo. El módulo siguiente despliega,
además, un conector regulado de ventana y calcula su residual como control
local. Ese control auxiliar no sustituye ni condiciona la realización
cofinal propietaria.

\section*{Realización cofinal propietaria y extinción terminal}

La realización que gobierna este capítulo es el complemento espectral--polar
del antecedente enriquecido común. Sea
\[
 \mathcal D_R=C_c^\infty((-R/2,R/2)),\qquad
 \mathcal D_R^{\rm cof}:=\mathcal D_R,\qquad
 \widehat\psi(t)=\int_{\mathbb R}\psi(x)e^{-itx}\,dx ,
\]
y, para \(\psi,\phi\in\mathcal D_R\),
\[
 h_{\psi,\phi}(z)
 =\widehat\psi(z)\overline{\widehat\phi(\overline z)},\qquad
 \check h_{\psi,\phi}(a)
 =\langle\psi,T_a\phi\rangle ,
 \qquad (T_a\phi)(x)=\phi(x-a).
\]

Sean \(P_+^{\rm sh},P_-^{\rm sh}\) los proyectores ortogonales de las dos
hojas y fíjense
\begin{equation}\label{eq:amp-rh-owner-q}
 q=\frac59,\qquad
 A=P_+^{\rm sh}+qP_-^{\rm sh},\qquad
 D=\sqrt{1-q^2}\,P_-^{\rm sh},\qquad A^*A+D^*D=I.
\end{equation}
Si \(\Xi_R\) es la restricción del antecedente común, \(Q_R=I-P_R\) su
complemento ortogonal y \(\jmath_-\) la inclusión isométrica en la hoja
antiespecular, se construye
\begin{equation}\label{eq:amp-rh-owner-binding}
 E_R:=\jmath_-Q_R\Xi_R,\qquad
 AE_R=qE_R,\qquad
 E_R^*E_R=\mathcal B_{{\rm W},R}.
\end{equation}
La identidad de conservación se itera sin borrar el último sumando:
\begin{equation}\label{eq:amp-rh-owner-telescopia}
 \mathcal B_{{\rm W},R}
 =\sum_{n=0}^{N-1}(DA^nE_R)^*(DA^nE_R)
 +(A^NE_R)^*(A^NE_R).
\end{equation}
El terminal se calcula antes de tomar el límite:
\begin{equation}\label{eq:amp-rh-owner-terminal}
 (A^NE_R)^*(A^NE_R)
 =q^{2N}\mathcal B_{{\rm W},R}\longrightarrow0
 \qquad(N\longrightarrow\infty).
\end{equation}
Sólo después de esta extinción se define
\begin{equation}\label{eq:amp-rh-owner-l}
 \mathscr L_Rf:=(DA^nE_Rf)_{n\geq0},
 \qquad
 \boxed{\mathcal B_{{\rm W},R}=\mathscr L_R^*\mathscr L_R\succeq0}.
\end{equation}
La compatibilidad de restricción publica estas factorizaciones como una
única forma cofinal:
\begin{equation}\label{eq:amp-rh-owner-publicacion}
 B_W=\Xi^*(I-P)\Xi
 =\bigl((I-P)\Xi\bigr)^*\bigl((I-P)\Xi\bigr)\succeq0,
 \qquad
 B_W(f,f)=\|(I-P)\Xi f\|^2.
\end{equation}

\begin{theorem}[Cierre cofinal Riemann--Weil]
\label{thm:amp-rh-cierre-cofinal-propietario}
La forma completa de Guinand--Weil es positiva sobre la clase cofinal
\(\bigcup_{R>0}\mathcal D_R^{\rm cof}=C_c^\infty(\mathbb R)\). Por el
criterio reversible de Weil, todo cero no trivial \(\rho\) de la función
zeta de Riemann satisface
\[
 \Re\rho=\frac12.
\]
\end{theorem}

\begin{proof}
Las ligaduras de \eqref{eq:amp-rh-owner-binding} convierten el terminal de
\eqref{eq:amp-rh-owner-telescopia} exactamente en el miembro izquierdo de
\eqref{eq:amp-rh-owner-terminal}. Su extinción geométrica permite pasar a
\eqref{eq:amp-rh-owner-l}; la naturalidad cofinal produce
\eqref{eq:amp-rh-owner-publicacion}. La identificación primo--gamma--polar
con el funcional completo y la separación cofinal del propietario principal
\eqref{eq:realizaciones-rh-identificacion-gw} permiten aplicar el criterio de
Weil. Los cálculos posteriores vuelven a comprobar esa publicación en una
carta regulada auxiliar, sin convertirse en premisas del teorema.
\end{proof}

\section*{Control auxiliar de regulador común e identidad de los tres canales}

El cálculo regulado que sigue despliega la diferencia de Gram desde sus
operadores constituyentes antes de evaluar su signo. Su residual es
\texttt{CONTROL\_NO\_GOBERNANTE}: delimita esta ventana concreta y no
reemplaza ni condiciona la realización cofinal
\eqref{eq:amp-rh-owner-binding}--\eqref{eq:amp-rh-owner-publicacion}.
Fijemos una función
\(\vartheta\in C_c^\infty([0,\infty))\) tal que
\(0\leq\vartheta\leq1\), \(\vartheta=1\) en \([0,1]\) y
\(\vartheta=0\) en \([2,\infty)\), y pongamos
\[
 \vartheta_L(a)=\vartheta(a/L),\qquad \vartheta_L(0)=1 .
\]
Es el único regulador de longitudes: se aplica simultáneamente a los
relojes primo--potencia y al canal gamma; los términos escalar y polar son
sus átomos de longitud cero. No se permite regularizar cada hoja con un
corte diferente.

Para
\[
 \ell_{p,k}=k\log p,\qquad
 w_{p,k}=(\log p)p^{-k/2},\qquad
 \mu_\Gamma(a)=\frac{e^{-a/2}}{1-e^{-2a}},
\]
definamos asimismo
\begin{align*}
 c_\Gamma&=\psi_\Gamma(1/4)-\log\pi<0,\\
 A_\psi&=\widehat\psi(i/2),&
 B_\psi&=\widehat\psi(-i/2),&
 b_\pm(\psi)&=\frac{A_\psi\pm B_\psi}{\sqrt2}.
\end{align*}
En el nivel \(N_m=9^m\), sean \(j_m\) la inclusión uniforme y
\(\widehat\delta_{a,m}\) el canal normalizado del acarreo nonádico. Sus
momentos se calculan directamente a partir del odómetro:
\begin{equation}\label{eq:amp-rh-momentos-odometro-explicitos}
 j_m^*j_m=I,\qquad
 \widehat\delta_{a,m}^*\widehat\delta_{b,m}
 =(I-T_a)^*(I-T_b).
\end{equation}
En particular, \(\widehat\delta_{a,m}\) depende de la costura y de
\(T_a\), no de la función zeta ni de sus ceros.

Las dos columnas reguladas son
\begin{align}
 \mathcal C^+_{m,R,L}\psi
 ={}&
 \Bigl(
  (\sqrt{\vartheta_L(\ell_{p,k})w_{p,k}}\,
       \widehat\delta_{\ell_{p,k},m}\psi)_{p,k},
 \nonumber\\[-2mm]
 &\hspace{11mm}
  a\longmapsto
  \sqrt{\vartheta_L(a)\mu_\Gamma(a)}\,
       \widehat\delta_{a,m}\psi,\,
  \zeta_+b_+(\psi)
 \Bigr),                                             \label{eq:amp-rh-columna-mas-explicita}\\
 \mathcal C^-_{m,R,L}\psi
 ={}&
 \Bigl(
  (\sqrt{2\vartheta_L(\ell_{p,k})w_{p,k}}\,j_m\psi)_{p,k},\,
  \sqrt{-c_\Gamma}\,j_m\psi,\,
  \zeta_-b_-(\psi)
 \Bigr),                                             \label{eq:amp-rh-columna-menos-explicita}
\end{align}
donde las sumas sólo recorren las longitudes para las que
\(\vartheta_L(\ell_{p,k})\neq0\), y \(\zeta_\pm\) son las líneas polares
ortogonales de la representación cíclica de orden \(12\).

\subsection*{Acoplamiento no descomponible anterior al signo}

La relación entre
\eqref{eq:amp-rh-columna-mas-explicita} y
\eqref{eq:amp-rh-columna-menos-explicita} no se obtiene postulando una
contracción entre dos Gram ya calculados. Sea
\[
 P_m=\iota_mE_m,\qquad Q_m=I-P_m
\]
la expectativa nonádica explícita. Para el reloj unitario
\(\mathscr U_{\ell,m}\), escribamos
\[
 a_{\ell,m}=P_m\mathscr U_{\ell,m}P_m,\qquad
 c_{\ell,m}=Q_m\mathscr U_{\ell,m}P_m ,
\]
de modo que la unidad de bloques da
\begin{equation}\label{eq:amp-rh-defecto-schwarz-reloj}
 P_m-a_{\ell,m}^*a_{\ell,m}
 =c_{\ell,m}^*c_{\ell,m}\succeq0.
\end{equation}
Para \(r=p^{-1/2}\), pongamos
\[
 s_{p,m}=\frac{N_mr}{1+(N_m-1)r},
 \qquad \xi_{p,m}=\operatorname{arsech}(s_{p,m}),
\]
y sea \(\mathscr S_{\xi_{p,m}}\) la coligación polar unitaria obtenida por
la transformación de Potapov--Ginzburg. Su composición de Redheffer con
el reloj verifica
\begin{align}
 C_{p,m}&=(s_{p,m}I-a_{\log p,m})
          (I-s_{p,m}a_{\log p,m})^{-1},\nonumber\\
 I-C_{p,m}^*C_{p,m}
 &=\bigl(1-s_{p,m}^2\bigr)
   (I-s_{p,m}a_{\log p,m}^*)^{-1}\nonumber\\
 &\quad{}\times c_{\log p,m}^*c_{\log p,m}
   (I-s_{p,m}a_{\log p,m})^{-1}.
   \label{eq:amp-rh-redheffer-defecto}
\end{align}
La expansión de Neumann, antes de cualquier publicación espectral, da
\begin{equation}\label{eq:amp-rh-torre-prima-explicita}
 \lambda_{p,m}(I-C_{p,m}^*C_{p,m})
 =\sum_{k\geq1}(\log p)p^{-k/2}
  \bigl(2I-T_{k\log p}-T_{k\log p}^*\bigr),
\end{equation}
con
\[
 \lambda_{p,m}
 =\frac{(\log p)r(1+r)}
 {(1-r)\kappa_{p,m}},
 \qquad
 \kappa_{p,m}
 =\frac{(1-s_{p,m}^2)(N_m-1)}
 {[N_m-(N_m-1)s_{p,m}]^2}.
\]
Así, las potencias de un primo son los coeficientes de una única
realimentación modular y no una lista elegida después.
Sea \(D_{p,L}\) la contracción diagonal del puerto de alcance de esta
realimentación que multiplica su coordenada \(k\) por
\(\sqrt{\vartheta_L(k\log p)}\). La regulación no altera el reloj ni sus
coeficientes: comprime simultáneamente las coordenadas de salida de la misma
torre.

El puerto finito de roles tampoco es diagonal. En los sectores activos
\[
 \mathcal H_+
 =\operatorname{span}\{f_3,f_6,f_9\},\qquad
 \mathcal H_-=\operatorname{span}\{f_4,f_8\},
\]
la coisometría
\begin{equation}\label{eq:amp-rh-acoplamiento-angular}
 C_\angle
 =|f_4\rangle\langle f_3|
  +|f_8\rangle\langle f_9|
\end{equation}
satisface
\[
 C_\angle C_\angle^*=I_{\mathcal H_-},\qquad
 I_{\mathcal H_+}-C_\angle^*C_\angle=P_6.
\]
La carta isométrica \(U_{W_{24}\to12}\), obtenida de las \(5\) octadas
incidentes y de su Gram
\(4I_5+\frac43J_5\), transporta esta coisometría a la frontera HMT:
\[
 C_{W_{24}}
 =U_{W_{24}\to12}^*C_\angle U_{W_{24}\to12}.
\]

\begin{theorem}[Control de la red Paley--Wiener--polar y su residual de puerto]
\label{thm:amp-rh-acoplamiento-prospectivo}
Para \(m,R,L\) fijos, definamos
\begin{align}
 \mathscr R^{\rm pr}_{m,L}
 &=
 \mathop{\star}_{\log p\leq2L}
 D_{p,L}\bigl(\mathscr S_{\xi_{p,m}}\star
              \mathscr U_{\log p,m}\bigr),\nonumber\\
 \mathscr R^\Gamma_{m,L}
 &=
 \int_0^{2L}{}^\oplus
 \mathscr U_{a,m}
 \vartheta_L(a)\,d\mu_\Gamma(a),\nonumber\\
 \mathscr C^{\rm TPK}_{m,R,L}
 &=C_{W_{24}}\star
 \bigl(\mathscr R^{\rm pr}_{m,L}
       \oplus\mathscr R^\Gamma_{m,L}\bigr).
 \label{eq:amp-rh-red-prospectiva}
\end{align}
Esta red es pasiva. En los espacios energéticos nativos de sus puertos,
sean
\(J_{+,m,R,L}:=\mathcal C^+_{m,R,L}\),
\(J_{-,m,R,L}:=\mathcal C^-_{m,R,L}\) y
\(\mathscr C^{\rm src}_{m,R,L}\) la transferencia contractiva externa de
\eqref{eq:amp-rh-red-prospectiva}. Se definen
\begin{equation}\label{eq:amp-rh-publicaciones-construidas}
 \begin{aligned}
 \mathcal E_{m,R,L}
 &:=\mathscr C^{\rm src}_{m,R,L}J_{+,m,R,L}
    -J_{-,m,R,L},\\
 \mathscr D_{m,R,L}
 &:=\bigl(I-(\mathscr C^{\rm src}_{m,R,L})^*
                  \mathscr C^{\rm src}_{m,R,L}\bigr)^{1/2},\\
 \mathsf L_{m,R,L}&:=\mathscr D_{m,R,L}J_{+,m,R,L}.
 \end{aligned}
\end{equation}
El residual \(\mathcal E_{m,R,L}\) mide si esta carta auxiliar de ventana
reproduce por sí sola la publicación ya construida por el propietario
cofinal. Su anulación no es una obligación del cierre Riemann--Weil ni una
premisa de su factorización positiva. Ningún coeficiente de
\eqref{eq:amp-rh-red-prospectiva} depende de un cero de \(\zeta\), de la
positividad de Weil ni del signo que se obtendrá.
\end{theorem}

\begin{proof}
Cada reloj \(\mathscr U_{a,m}\) es unitario.
La identidad \eqref{eq:amp-rh-defecto-schwarz-reloj} identifica su puerto
de acarreo. La transformación polar es unitaria y
\eqref{eq:amp-rh-redheffer-defecto} demuestra pasividad después de cerrar
el puerto interno. La composición de Redheffer de un número finito de
coligaciones pasivas vuelve a ser pasiva. El canal gamma se obtiene como
límite de sumas directas en \([\varepsilon,2L]\); la estimación
\(\mu_\Gamma(a)\|(I-T_a)\psi\|^2=O(a)\) cuando \(a\downarrow0\)
da el límite en el dominio de forma.

La expansión \eqref{eq:amp-rh-torre-prima-explicita} produce exactamente
los coeficientes primo--potencia. La integral directa produce la segunda
entrada de \eqref{eq:amp-rh-columna-mas-explicita}. La transformación de
Hadamard en el puerto polar produce \(b_+\) y \(b_-\), y
\eqref{eq:amp-rh-acoplamiento-angular}, conjugada por
\(U_{W_{24}\to12}\), fija cuáles son las \(3\) entradas y las \(2\)
salidas del mismo estado. Estas operaciones construyen la transferencia y
las dos columnas nativas; su diferencia de salida es, exactamente,
\(\mathcal E_{m,R,L}\) en
\eqref{eq:amp-rh-publicaciones-construidas}. La pasividad garantiza que
\(\mathscr D_{m,R,L}\) está bien definida, pero no anula por sí sola ese
residual auxiliar; tampoco necesita hacerlo para el cierre cofinal ya
demostrado.

Esta red no es \(C_\angle\otimes I\): los factores
\((I-s_{p,m}a_{\log p,m})^{-1}\) mezclan cada puerto polar con todas las
potencias del reloj, y la integral directa mezcla el mismo puerto con el
continuo gamma. Suprimir esos términos cambia la columna de alcance y no
es una simplificación de \eqref{eq:amp-rh-red-prospectiva}.
\end{proof}

\subsection*{Expansión término a término de la diferencia de Gram}

\begin{lemma}[Identidad primo--gamma--polar con regulador común]
\label{lem:amp-rh-expansion-guinand-weil}
Para \(\psi,\phi\in\mathcal D_R\),
\begin{equation}\label{eq:amp-rh-diferencia-columnas-regulada}
 \langle J_{+,m,R,L}\psi,J_{+,m,R,L}\phi\rangle
 -\langle J_{-,m,R,L}\psi,J_{-,m,R,L}\phi\rangle
 =
 \mathscr I_{{\rm GW},L}(\psi,\phi),
\end{equation}
donde
\begin{align}
 \mathscr I_{{\rm GW},L}(\psi,\phi)
 ={}&
 h_{\psi,\phi}(i/2)+h_{\psi,\phi}(-i/2)
 +c_\Gamma\check h_{\psi,\phi}(0)                 \nonumber\\
 &+\int_0^\infty\vartheta_L(a)\mu_\Gamma(a)
 \bigl(
 2\check h_{\psi,\phi}(0)
 -\check h_{\psi,\phi}(a)
 -\check h_{\psi,\phi}(-a)
 \bigr)\,da                                       \nonumber\\
 &-\sum_{p,k}\vartheta_L(\ell_{p,k})w_{p,k}
 \bigl(
 \check h_{\psi,\phi}(\ell_{p,k})
 +\check h_{\psi,\phi}(-\ell_{p,k})
 \bigr).                              \label{eq:amp-rh-gw-truncado-explicito}
\end{align}
Éste es el funcional completo truncado de Guinand--Weil: contiene el
término polar, el factor gamma y las \(2\) orientaciones de cada
primo--potencia bajo un único regulador.
\end{lemma}

\begin{proof}
Por las definiciones
\eqref{eq:amp-rh-columna-mas-explicita}--
\eqref{eq:amp-rh-columna-menos-explicita},
\begin{align}
 \langle J_+\psi,J_+\phi\rangle
 -\langle J_-\psi,J_-\phi\rangle
 &=
 \langle\mathcal C^+\psi,\mathcal C^+\phi\rangle
 -\langle\mathcal C^-\psi,\mathcal C^-\phi\rangle .
 \label{eq:amp-rh-expansion-inicial}
\end{align}
En el canal primo, \eqref{eq:amp-rh-momentos-odometro-explicitos} da
\begin{align*}
 &\langle(I-T_\ell)\psi,(I-T_\ell)\phi\rangle
 -2\langle\psi,\phi\rangle\\
 &\hspace{18mm}
 =-\check h_{\psi,\phi}(\ell)
  -\check h_{\psi,\phi}(-\ell).
\end{align*}
Multiplicar por
\(\vartheta_L(\ell_{p,k})w_{p,k}\) y sumar produce exactamente la última
línea de \eqref{eq:amp-rh-gw-truncado-explicito}.

En el canal gamma,
\[
 \langle(I-T_a)\psi,(I-T_a)\phi\rangle
 =2\check h(0)-\check h(a)-\check h(-a).
\]
Además, la representación integral de la derivada logarítmica de gamma
da, para \(t\in\mathbb R\),
\begin{equation}\label{eq:amp-rh-digamma-longitudes}
 \Re\psi_\Gamma\!\left(\frac14+\frac{it}{2}\right)-\psi_\Gamma(1/4)
 =2\int_0^\infty\mu_\Gamma(a)(1-\cos at)\,da .
\end{equation}
Al integrar contra \(h(t)/(2\pi)\) y usar inversión de Fourier se obtiene
\begin{align*}
 &-\log\pi\,\check h(0)
 +\frac1{2\pi}\int_{\mathbb R}h(t)
  \Re\psi_\Gamma\!\left(\frac14+\frac{it}{2}\right)\,dt\\
 &\qquad =
 c_\Gamma\check h(0)
 +\int_0^\infty\mu_\Gamma(a)
  [2\check h(0)-\check h(a)-\check h(-a)]\,da .
\end{align*}
La versión regulada aplica en esta identidad la sustitución común
\[
 \mu_\Gamma(a)\longmapsto
 \vartheta_L(a)\mu_\Gamma(a).
\]

Finalmente,
\begin{align*}
 \langle b_+(\psi),b_+(\phi)\rangle
 -\langle b_-(\psi),b_-(\phi)\rangle
 &=
 A_\psi\overline{B_\phi}
 +B_\psi\overline{A_\phi}\\
 &=h_{\psi,\phi}(i/2)+h_{\psi,\phi}(-i/2).
\end{align*}
La suma de las \(3\) expansiones es
\eqref{eq:amp-rh-gw-truncado-explicito}. Cada igualdad se ha calculado
antes de conocer el signo de su suma.
\end{proof}

\section*{Convergencia, cofinalidad y separación}

\begin{proposition}[Eliminación del regulador y naturalidad cofinal]
\label{prop:amp-rh-convergencia-cofinal}
Para toda \(\psi,\phi\in\mathcal D_R\), el límite
\[
 \lim_{L\to\infty}\mathscr I_{{\rm GW},L}(\psi,\phi)
 =\mathscr I_{\rm GW}(\psi,\phi)
\]
existe y coincide con la forma completa de Guinand--Weil. Si
\(R\leq R'\), las inclusiones
\(\mathcal D_R\hookrightarrow\mathcal D_{R'}\) y los refinamientos
\(m\mapsto m+1\) entrelazan separadamente
\(J_{+,m,R,L}\) y \(J_{-,m,R,L}\).
Por tanto las identidades
\eqref{eq:amp-rh-diferencia-columnas-regulada} forman una única familia
compatible, no una colección de factorizaciones dependientes de la
ventana.
\end{proposition}

\begin{proof}
Como \(\check h_{\psi,\phi}\in C_c^\infty((-R,R))\), la suma primo--potencia
que queda después de cancelar los contraterminos escalares es localmente
finita: todos los términos con \(\ell_{p,k}\geq R\) se anulan. En la
integral gamma,
\[
 2\check h(0)-\check h(a)-\check h(-a)=O(a^2)
 \quad(a\downarrow0),
\]
mientras \(\mu_\Gamma(a)\sim(2a)^{-1}\); el integrando es \(O(a)\).
Para \(a>R\), las dos traslaciones desaparecen y
\(\mu_\Gamma(a)=O(e^{-a/2})\). La convergencia dominada permite
\(L\to\infty\).

La identidad de momentos
\eqref{eq:amp-rh-momentos-odometro-explicitos} no depende de \(m\);
por ello define la isometría de refinamiento
\(\widehat\delta_{a,m}\mapsto\widehat\delta_{a,m+1}\), mientras
\(j_m\mapsto j_{m+1}\). Las líneas polares y
\(U_{W_{24}\to12}\) permanecen fijos. Al ampliar la ventana, los relojes
nuevos tienen traslaciones disjuntas sobre \(\mathcal D_R\) y añaden el
mismo término \(2w_{p,k}\langle\psi,\phi\rangle\) a las dos columnas antes
de cancelarse. Esto demuestra ambos entrelazamientos.
\end{proof}

\paragraph{Control finito heredado de la fibra nonádica.}
Sea \(\mathscr F_{729}\) el factor local finito de \(729\) estados,
\(P^{[729]}\) el proyector sobre los \(81\) estados gruesos y
\(Q^{[729]}=I_{\mathscr F_{729}}-P^{[729]}\). Entonces
\begin{equation}\label{eq:amp-rh-729-81-648}
 729=81+648,\qquad
 \operatorname{rank}P^{[729]}=81,\qquad
 \operatorname{rank}Q^{[729]}=648.
\end{equation}
Este censo precede al criterio de signo y no identifica el complemento
\(Q^{[729]}\) con el residual de conector
\(\mathcal E_{m,R,L}\).

\begin{proposition}[Compatibilidad cofinal del complemento finito]
\label{prop:amp-rh-compatibilidad-complemento-finito}
\label{prop:amp-rh-residual}
Las identidades de publicación común de
\eqref{eq:realizaciones-rh-publicaciones-comunes}--
\eqref{eq:realizaciones-rh-bw} forman, bajo los refinamientos
\(m\mapsto m+1\) y las inclusiones
\(\mathcal D_R\hookrightarrow\mathcal D_{R'}\), una única familia
compatible. En particular, la diferencia de los dos Gram se transporta
como la forma del complemento \((I-P)\Xi\), y el censo
\eqref{eq:amp-rh-729-81-648} permanece como control finito de la fibra.
\end{proposition}

\begin{proof}
La ortogonalidad de \(P\) y las publicaciones comunes dan
\eqref{eq:realizaciones-rh-bw}. La
\cref{prop:amp-rh-convergencia-cofinal} entrelaza separadamente las dos
columnas bajo refinamiento e inclusión; al restarlas, la identidad del
complemento se conserva en toda la familia cofinal.
\end{proof}

En esta carta auxiliar no se afirma
\(\mathcal E_{m,R,L}=0\): la
\cref{prop:amp-rh-obstruccion-ventana-corta} delimita exclusivamente ese
modelo regulado. El censo recuperado pertenece al complemento finito
propietario y no al control de ventana corta; por tanto, el residual no
reabre ni condiciona el teorema cofinal.

\begin{lemma}[La clase cofinal separa las evaluaciones espectrales]
\label{lem:amp-rh-separacion-cofinal}
Se tiene la igualdad, no sólo la densidad,
\[
 \bigcup_{R>0}\mathcal D_R=C_c^\infty(\mathbb R).
\]
Además, para cualesquiera puntos complejos distintos
\(z_1,\ldots,z_n\), la aplicación
\[
 \mathcal D_R\longrightarrow\mathbb C^n,\qquad
 \psi\longmapsto
 (\widehat\psi(z_1),\ldots,\widehat\psi(z_n))
\]
es sobreyectiva para todo \(R\) suficientemente grande. En particular,
ninguna órbita finita de evaluaciones del criterio de Weil queda
invisible para la familia cofinal.
\end{lemma}

\begin{proof}
La primera afirmación sigue de la compacidad del soporte. Para la segunda,
si los funcionales de evaluación fueran linealmente dependientes, existirían
\(c_1,\ldots,c_n\), no todos nulos, tales que
\[
 \int_{\mathbb R}\psi(x)
 \left(\sum_{j=1}^nc_je^{-iz_jx}\right)\,dx=0
 \qquad(\psi\in\mathcal D_R).
\]
La función analítica
\(\sum_jc_je^{-iz_jx}\) se anularía en el intervalo
\((-R/2,R/2)\), luego en todo \(\mathbb C\). La independencia de los
exponenciales con frecuencias distintas fuerza \(c_j=0\) para todo \(j\),
contradicción. Los funcionales son independientes; una aplicación lineal
hacia \(\mathbb C^n\) con adjunto inyectivo es sobreyectiva.
\end{proof}

\section*{Acoplamiento de 2 rutas y prueba adversarial de meseta}

El producto desnudo \(C_\angle\otimes I\) exigiría
\[
 2\|\psi\|^2\leq\|(I-T_\ell)\psi\|^2
\]
para toda \(\psi\), desigualdad que falla en mesetas suaves mucho más
largas que \(\ell\). La red
\eqref{eq:amp-rh-red-prospectiva} no hace esa identificación: conserva
los cruces Paley--Wiener, gamma, primo y polar antes de proyectar. El
siguiente cálculo exhibe esa diferencia sobre una meseta concreta.

Sea \(\omega=e^{2\pi i/9}\) y, en la fibra superviviente,
\[
 \kappa_0=u_{81}\otimes\chi_6,\qquad
 \kappa_1=v_{10}\otimes\chi_6,\qquad
 \chi_6(c)=3^{-1}\omega^{6c}.
\]
Son \(2\) rutas ortogonales del mismo rol \(P_6\). Su publicación de
frontera produce \(e_0,e_1\) con
\[
 \langle e_i,\Lambda_{\rm tw}^{\#}e_j\rangle=\delta_{ij}.
\]
Tomemos
\[
 a=\frac1{16},\qquad b=2\log2,\qquad
 \psi_w=w^{-1/2}{\bf1}_{[-w/2,w/2]}\quad(w=a,b).
\]
Estas funciones pertenecen al cierre de forma de \(\mathcal D_R\), y sus
mollificaciones internas pertenecen a \(\mathcal D_R\). Si
\(g_{ij}=\mathscr I_{\rm GW}(\psi_{w_i},\psi_{w_j})\), con
\((w_0,w_1)=(a,b)\), la expansión anterior da, para
\(\lambda_n=2n+\tfrac12\),
\begin{align}
 g(w,w)
 ={}&\frac2w\sum_{n\geq0}
 \frac{1-e^{-\lambda_nw}}{\lambda_n^2}
 +c_\Gamma+\frac{32}{w}\sinh^2\!\frac w4 \nonumber\\
 &-2\sum_{k\log p<w}
  w_{p,k}\left(1-\frac{k\log p}{w}\right),
 \label{eq:amp-rh-meseta-diagonal}\\
 g_{01}
 ={}&\frac2{\sqrt{ab}}\sum_{n\geq0}
 \frac{e^{-\lambda_n(b-a)/2}(1-e^{-\lambda_na})}{\lambda_n^2}
 +c_\Gamma\sqrt{\frac ab} \nonumber\\
 &+2A_aA_b-\frac{\log2}{\sqrt2}\sqrt{\frac ab},
 \qquad A_w=\frac{4\sinh(w/4)}{\sqrt w}.
 \label{eq:amp-rh-meseta-cruce}
\end{align}
La acotación de las colas geométricas y de la serie de
\(\lambda_n^{-2}\) produce
\begin{align}
 1.46610954095406&<g_{00}<1.46690960095857,\nonumber\\
 0.06966817455957&<g_{11}<0.06970424464086,\nonumber\\
 -0.008430670493497&<g_{01}<-0.008430670493494,
 \label{eq:amp-rh-meseta-cotas}
\end{align}
y, por multiplicación intervalar,
\begin{equation}\label{eq:amp-rh-meseta-determinante}
 0.10207009921767<
 \det\begin{pmatrix}g_{00}&g_{01}\\g_{01}&g_{11}\end{pmatrix}
 <0.10217874948627.
\end{equation}
El Gram completo de la meseta y del estado basal tiene, por tanto, rango
\(2\); una única línea escalar no puede publicarlo.

Pongamos
\[
 \sigma_b=g_{11}-\frac{g_{01}^2}{g_{00}}>0
\]
y definamos sobre
\(\mathcal E_2=\operatorname{span}\{\psi_a,\psi_b\}\)
\begin{align}
 H^{(2)}_{m,R}(\alpha\psi_a+\beta\psi_b)
 ={}&
 e_{0,m}\left(
 \sqrt{g_{00}}\,\alpha
 +\frac{g_{01}}{\sqrt{g_{00}}}\,\beta\right)
 +e_{1,m}\sqrt{\sigma_b}\,\beta .
 \label{eq:amp-rh-lector-dos-rutas}
\end{align}
Una multiplicación de Cholesky, usando
\(\langle e_i,\Lambda_m^\#e_j\rangle=\delta_{ij}\), demuestra
\begin{equation}\label{eq:amp-rh-meseta-energia-exacta}
 \langle H^{(2)}_{m,R}\psi,
         \Lambda_m^\#H^{(2)}_{m,R}\phi\rangle
 =\mathscr I_{\rm GW}(\psi,\phi)
 \qquad(\psi,\phi\in\mathcal E_2).
\end{equation}
Los coeficientes de \eqref{eq:amp-rh-lector-dos-rutas} proceden de las
series explícitas
\eqref{eq:amp-rh-meseta-diagonal}--
\eqref{eq:amp-rh-meseta-cruce}; no se introducen ceros ni se presupone el
signo del determinante, que se prueba en
\eqref{eq:amp-rh-meseta-determinante}. Por continuidad de los \(3\)
canales, la identidad y la positividad estricta del determinante persisten
para mollificaciones suaves suficientemente próximas. Ésta es la prueba
adversarial que el tensor \(C_\angle\otimes I\) no supera y el acoplamiento
no descomponible sí.

\section*{Residual exacto y delimitación del regulador}

Definamos
\begin{equation}\label{eq:amp-rh-residual-finito}
 \mathcal R_{m,R,L}
 :=
 (\mathcal C^+_{m,R,L})^*\mathcal C^+_{m,R,L}
 -(\mathcal C^-_{m,R,L})^*\mathcal C^-_{m,R,L}
 -\mathscr I_{{\rm GW},L}.
\end{equation}
El lema \ref{lem:amp-rh-expansion-guinand-weil} demuestra
\[
 \mathcal R_{m,R,L}=0
\]
término a término, no por definición. Al mismo tiempo,
\eqref{eq:amp-rh-publicaciones-construidas} da, por cálculo de normas, el
balance de puerto
\begin{equation}\label{eq:amp-rh-balance-residual-puerto}
 \boxed{
 \begin{aligned}
 \mathscr I_{{\rm GW},L}(\psi,\psi)
 -\|\mathsf L_{m,R,L}\psi\|^2
 ={}&2\operatorname{Re}
 \langle J_{-,m,R,L}\psi,\mathcal E_{m,R,L}\psi\rangle\\
 &+\|\mathcal E_{m,R,L}\psi\|^2 .
 \end{aligned}}
\end{equation}
En efecto,
\(\|\mathsf L\psi\|^2
=\|J_+\psi\|^2-\|\mathscr C^{\rm src}J_+\psi\|^2\), y
\(\mathscr C^{\rm src}J_+=J_-+\mathcal E\). Por tanto, dentro de esta carta
auxiliar, una anulación demostrada de \(\mathcal E\) produciría el cuadrado positivo
\begin{equation}\label{eq:amp-rh-cuadrado-regulado}
 \mathcal E_{m,R,L}=0
 \quad\Longrightarrow\quad
 \mathscr I_{{\rm GW},L}(\psi,\psi)
 =\|\mathsf L_{m,R,L}\psi\|^2\geq0.
\end{equation}
La implicación no autoriza a suponer la premisa y no forma parte de la cadena
propietaria del cierre cofinal.

\begin{proposition}[Obstrucción de ventana corta]
\label{prop:amp-rh-obstruccion-ventana-corta}
Para las columnas
\eqref{eq:amp-rh-columna-mas-explicita}--
\eqref{eq:amp-rh-columna-menos-explicita}, el residual
\(\mathcal E_{m,R,L}\) no puede anularse para toda
\(\psi\in\mathcal D_R\) y todo \(L>0\).
\end{proposition}

\begin{proof}
Elíjase
\[
 0\ne\psi\in\mathcal D_R,\qquad
 \widehat\psi(i/2)=\widehat\psi(-i/2)=0.
\]
Tal función existe porque se imponen dos funcionales lineales sobre el
espacio infinito-dimensional \(\mathcal D_R\). Si \(2L<\log2\), la suma
primo--potencia de
\eqref{eq:amp-rh-gw-truncado-explicito} es vacía y el canal polar se
anula. Para \(h_{\psi,\psi}\), se tiene
\(\check h_{\psi,\psi}(0)=\|\psi\|_2^2\), de modo que
\begin{equation}\label{eq:amp-rh-ventana-corta-expansion}
 \mathscr I_{{\rm GW},L}(\psi,\psi)
 =c_\Gamma\|\psi\|_2^2+
 \int_0^{2L}\vartheta_L(a)\mu_\Gamma(a)
       \|(I-T_a)\psi\|_2^2\,da .
\end{equation}
Como
\(\|(I-T_a)\psi\|_2\leq a\|\psi'\|_2\) y
\(\mu_\Gamma(a)=O(a^{-1})\) cuando \(a\downarrow0\), la integral de
\eqref{eq:amp-rh-ventana-corta-expansion} es \(O_\psi(L^2)\). Puesto que
\(c_\Gamma<0\), para \(L\) suficientemente pequeño resulta
\[
 \mathscr I_{{\rm GW},L}(\psi,\psi)<0.
\]
Si \(\mathcal E_{m,R,L}\psi=0\), la implicación
\eqref{eq:amp-rh-cuadrado-regulado} daría el signo contrario. Luego el
residual no puede anularse con el cuantificador indicado.
\end{proof}

La igualdad \(\mathcal R_{m,R,L}=0\) identifica correctamente la
diferencia de los dos Gram con el funcional truncado; no la convierte en
un cuadrado positivo. Este conector regulado es, por tanto,
\texttt{CONTROL\_NO\_GOBERNANTE}: no interviene en la cadena probatoria
\eqref{eq:amp-rh-owner-binding}--\eqref{eq:amp-rh-owner-publicacion}. La
proposición \ref{prop:amp-rh-obstruccion-ventana-corta} delimita
exclusivamente las fórmulas
\eqref{eq:amp-rh-columna-mas-explicita}--
\eqref{eq:amp-rh-gw-truncado-explicito}; su resultado no demueve la
positividad cofinal ni el \cref{thm:amp-rh-cierre-cofinal-propietario}.


\section*{Necesidad de las hipótesis y obstrucciones exactas}

\begin{ablation}[Necesidad del antecedente común y de la clase completa]
\label{abl:realizaciones-rh}
Ninguna de las siguientes operaciones conserva el teorema
\ref{thm:realizaciones-rh}: reemplazar \(P\) por un idempotente no
ortogonal; realizar \(\Xi_+\) y \(\Xi_-\) desde antecedentes distintos;
eliminar uno de los canales primo, gamma o polar; restringir la prueba a una
familia no cofinal; anular el acarreo por definición; suprimir el terminal de
\eqref{eq:realizaciones-rh-ref9}; o construir \((\Xi,P)\) después de imponer
\(B_W\geq0\).
\end{ablation}

\begin{proof}
Sin ortogonalidad, \(I-P\) no proporciona el cuadrado de
\eqref{eq:realizaciones-rh-factorizacion-global}. Con dos antecedentes, la
diferencia de Gram deja de ser el complemento de una sola proyección. Omitir
un canal sustituye \(\mathscr I_{\rm GW}\) por otro funcional. La pérdida de
cofinalidad o separación impide aplicar la equivalencia
\eqref{eq:realizaciones-rh-criterio}. Borrar acarreo o terminal rompe las
identidades finitas anteriores al límite. Finalmente, elegir el antecedente
desde la desigualdad invertiría la dirección causal y convertiría la
factorización en una reformulación circular.
\end{proof}

Un contraejemplo directo al enunciado sería una función
\(f\in\mathcal D_R^{\rm cof}\) con \(B_W(f,f)<0\) que conservase todas las
premisas, o un cero no trivial fuera de la recta crítica compatible con la
equivalencia completa \eqref{eq:realizaciones-rh-criterio}. Las
verificaciones finitas controlan, además, la identidad de memoria, la contracción por vuelta
y la persistencia del acarreo antes del paso cofinal.

\section*{Alcance de la demostración}

La doble publicación, la conservación genealógica, la conexión conjunta y la
telescopía con terminal construyen la diferencia de Gram, su factorización
positiva y la publicación cofinal común. La identidad completa
primo--gamma--polar reside en
\eqref{eq:realizaciones-rh-identificacion-gw}; la expansión regulada y su
límite la reproducen sólo como comprobación independiente de la carta
auxiliar.
El residual \(\mathcal E_{m,R,L}\) de una realización regulada de ventana
puede no anularse, como muestra
\cref{prop:amp-rh-obstruccion-ventana-corta}; por ello se conserva como
\texttt{CONTROL\_NO\_GOBERNANTE} y no se identifica con
\(\mathfrak E_R^{\rm com}\) ni reabre el cierre cofinal. El lenguaje
convencional interviene en \eqref{eq:realizaciones-rh-criterio} como
reconocimiento final, no como origen de \(\Xi\), de \(P\) o de la forma
positiva.

\chapter{Almacenamiento terminal Stein--Lyapunov, cota uniforme y regularidad
global de Navier--Stokes tridimensional}
\chaptermark{Almacenamiento terminal y regularidad de Navier--Stokes}

\noindent La especialización del
\cref{thm:nucleo-continuidad-conexion-nonadica-conjunta} es
\[
 (\widehat X_\infty,\Gamma_9)
 \xrightarrow{\ \mathcal R_{\rm NS}\ }
 U^{\rm own}_{M,j}
 \longrightarrow\text{identidad de Stein con terminal}
 \longrightarrow\sup_{M,t}\|u_M(t)\|_{H^s}<\infty
 \longrightarrow u\in C^\infty.
\]
El presupuesto raíz uniforme pertenece al realizador; no se obtiene
suponiendo la regularidad del campo publicado.

\section*{Objeto y dominio hidrodinámico}

Sea \(\mathbb T^3\) el toro tridimensional con medida normalizada y sea
\[
 H^s_\sigma(\mathbb T^3)
 =\{u\in H^s(\mathbb T^3;\mathbb R^3):
      \nabla\!\cdot u=0,\ \textstyle\int_{\mathbb T^3}u=0\}.
\]
Se fijan
\begin{equation}\label{eq:realizaciones-ns-dominio}
 s>\frac52,
 \qquad \nu>0,
 \qquad u_0\in C^\infty\cap H^s_\sigma(\mathbb T^3),
 \qquad
 f\in C^\infty([0,\infty);H^\infty_\sigma)
       \cap L^2_{\rm loc}([0,\infty);H^{s-1}_\sigma).
\end{equation}
Sea \(P_M\) la proyección de Leray--Galerkin sobre los primeros modos y sean
\[
 A_M=-P_M\Delta,
 \qquad
 B_M(v,w)=P_MP_\sigma((v\!\cdot\!\nabla)w).
\]
El campo publicado por el corte \(M\) satisface
\begin{equation}\label{eq:realizaciones-ns-galerkin}
 \partial_tu_M+\nu A_Mu_M+B_M(u_M,u_M)=f_M,
 \qquad u_M(0)=P_Mu_0.
\end{equation}

El objeto de partida en cada profundidad nonádica \(j\) es el estado
extendido
\begin{equation}\label{eq:realizaciones-ns-estado}
 \begin{aligned}
 \Xi_{M,j}^{\rm NS}
 =\bigl(&u_M;q_{A,M,j},q_{V,M,j},T_{M,j};\\[-1mm]
        &\text{tríadas, hoja, orientación, acarreo, frontera,}\\[-1mm]
        &\text{ruta, archivo y memoria}\bigr).
 \end{aligned}
\end{equation}
El lector \(R_{M,j}\Xi_{M,j}^{\rm NS}=u_M\) publica
\eqref{eq:realizaciones-ns-galerkin}. La representación polinómica de
Carleman se entiende como forma sobre un núcleo común de una escala de radios;
no se la trata como operador acotado de un único espacio de Fock en sí mismo.
Conserva conjuntamente todos los grados requeridos por la no linealidad. Si
\(\mathfrak U_{M,J,t}\) es la dilatación isométrica del registro completo,
el retorno se define por
\begin{equation}\label{eq:realizaciones-ns-retorno-completo}
 \mathcal R_{M,J,t}^{\rm enr}
 =\mathcal R_{1,M}\mathfrak U_{M,J,t}^*,
 \qquad
 \mathcal R_{M,J,t}^{\rm enr}Z_{M,J,t}=u_M(t).
\end{equation}
El adjunto reúne las ramas antes de evaluar. Leer una sola rama, cuya norma
crece como \(3^J\), no define el retorno enriquecido.

\section*{Acarreo de Galerkin y dos memorias positivas}

Para \(N>M\), el defecto de conmutación de la no linealidad con la
proyección es
\begin{equation}\label{eq:realizaciones-ns-acarreo}
 C_{M\leftarrow N}
 =P_MB(u_N,u_N)-B_M(P_Mu_N,P_Mu_N).
\end{equation}
Esta cantidad es el acarreo de Galerkin. Se incorpora al estado al refinar y
evita identificar la interacción completa con la interacción de su imagen
visible.

La razón areal--volumétrica fija
\begin{equation}\label{eq:realizaciones-ns-r}
 r=\frac{\pi_{\rm HMT}}{\varphi_{\rm HMT}^{\,2}}>1.
\end{equation}
Sean \(\mathcal M_M^+\geq0\) y \(\mathcal M_M^-\geq0\) las memorias de las dos orientaciones.
Sus coordenadas diferencia y suma son
\begin{equation}\label{eq:realizaciones-ns-dh}
 D_M=(r-1)(\mathcal M_M^+-\mathcal M_M^-),
 \qquad
 H_M=(1+r)(\mathcal M_M^++\mathcal M_M^-).
\end{equation}
Si \(\mathcal B_{s,M}\) denota el flujo convectivo orientado en la energía
\(H^s\), se tiene
\begin{equation}\label{eq:realizaciones-ns-dh-derivadas}
 \dot D_M=\mathcal B_{s,M},
 \qquad
 \dot H_M=\frac{1+r}{r-1}|\mathcal B_{s,M}|,
\end{equation}
o, de manera equivalente,
\begin{equation}\label{eq:realizaciones-ns-memorias-positivas}
 (r-1)\dot{\mathcal M}_M^+=\mathcal B_{s,M}^{+},
 \qquad
 (r-1)\dot{\mathcal M}_M^-=\mathcal B_{s,M}^{-}.
\end{equation}
El signo se publica en \(D_M\); la memoria de cada hoja permanece positiva.
El cociente \(\exp(D_M)\) puede disminuir sin que se pierda la conservación
del estado, pues no se identifica con la suma positiva \(H_M\).

\section*{Realización íntegra del propietario Stein--Lyapunov}

La solución propietaria no identifica el cierre con una desigualdad KYP
firmada. Se presentan a continuación, sin abreviación, las dos residencias
del cierre autónomo: la carta PC--Fock completa con sus controles
alternativos y el owner \(U^{\rm own}\) con almacenamiento terminal,
telescopía, ledger, presupuesto raíz, retorno físico y paso global.

\section{Carta PC--Fock completa, memoria y almacenamiento terminal}
\label{sec:ns-pcf-prospective-construction}

Fijamos \(s>5/2\), \(\nu>0\), un corte de Galerkin \(M\), una profundidad
nonádica \(J\) y un horizonte \(T<\infty\). Sea
\(V_M\subset H^s_\sigma(\mathbb T^3)\) el espacio solenoidal de media cero y
sean
\begin{equation}
 A_M=P_M(-\Delta)P_M,
 \qquad
 N_M(u)=P_M\mathbb P(u\cdot\nabla u).
 \label{eq:ns-import-galerkin-operators}
\end{equation}
Esta sección construye el portador PC--Fock completo, sus lectores, memorias,
acarreo, almacenamiento terminal y retorno, que ingresan en la realización
propietaria de la sección siguiente. Su KYP local, su observabilidad NS--OBS y
las dominaciones LRDN--LCN son exclusivamente
\texttt{CONTROL\_NO\_GOBERNANTE}: no condicionan la métrica
\(U^{\rm own}\), su telescopía, su presupuesto raíz ni la conclusión global.

\subsection{Torre completa y lectores de grados uno y dos}

La elevación física se toma en la Fock simétrica completa:
\begin{equation}
 \mathcal V_M^{\rm PCF}(u)
 :=\bigoplus_{n\ge0}\frac{R_{\rm F}^n}{\sqrt{n!}}u^{\odot n},
 \qquad
 \|\mathcal V_M^{\rm PCF}(u)\|^2
 =e^{R_{\rm F}^2\|u\|_{H^s}^2}.
 \label{eq:ns-import-pcf-lift}
\end{equation}
La torre no se trunca. Para \(z_{n,M}=u_M^{\odot n}\), la regla de Leibniz
da la recurrencia exacta
\begin{equation}
 \dot z_{n,M}
 =L_{n,M}z_{n,M}
 +\mathcal N_{n,n+1;M}z_{n+1,M}
 +\mathcal F_{n,n-1;M}(f_M)z_{n-1,M},
 \qquad n\ge1,
 \label{eq:ns-import-carleman-all-degrees}
\end{equation}
donde, sobre la diagonal coherente,
\begin{align}
 L_{n,M}u^{\odot n}
 &=-\nu n(A_Mu)\odot u^{\odot(n-1)}, \notag\\
 \mathcal N_{n,n+1;M}u^{\odot(n+1)}
 &=-nN_M(u)\odot u^{\odot(n-1)}, \notag\\
 \mathcal F_{n,n-1;M}(f)u^{\odot(n-1)}
 &=nf\odot u^{\odot(n-1)}.
 \label{eq:ns-import-carleman-blocks}
\end{align}

El espacio de estado es la suma ortogonal
\begin{equation}
 \mathscr K_{M,j}^{\rm PCF}
 =\mathscr K_{M,j}^{\rm HMT}\widehat\otimes\Gamma_s(V_M)
 \oplus\mathscr K_{M,j}^{\kappa}
 \oplus\mathscr K_{M,j}^{\rm micro}
 \oplus\mathscr K_{M,j}^{\rm sh},
 \label{eq:ns-import-full-carrier}
\end{equation}
donde \(\mathscr K_{M,j}^{\rm HMT}\) conserva hoja, orientación,
\(q_A,q_V,T\), borde, ruta y las dos memorias \(D/H\). Los grados uno y dos
se leen mediante
\begin{equation}
 \begin{aligned}
 z_M(u)&=\nu^{1/2}A_M^{1/2}\Lambda^su,\\
 w_M(u)&=\nu^{-1/2}A_M^{-1/2}\Lambda^sN_M(u),\\
 \mathcal Y_{\beta,M}\mathcal V_M^{\rm PCF}(u)
 &=\sqrt\beta\,z_M(u)
   \oplus\frac{1}{2\sqrt\beta}\,w_M(u).
 \end{aligned}
 \label{eq:ns-import-degree-one-two-reader}
\end{equation}
En consecuencia,
\begin{equation}
 |B_{s,M}(u)|
 \le
 \|\mathcal Y_{\beta,M}\mathcal V_M^{\rm PCF}(u)\|^2
 =\beta\nu D_{s,M}(u)+\frac{\|w_M(u)\|^2}{4\beta}.
 \label{eq:ns-import-young-port}
\end{equation}
La segunda coordenada es un operador lineal uniformemente acotado sobre
\(\operatorname{Sym}^2H^s_\sigma\):
\begin{equation}
 \sup_M\|\mathfrak C_M\|
 \le\widetilde M_{s,\beta,\nu},
 \qquad
 C_M^{\rm obs}:=\sqrt2\,\mathfrak C_M\pi_2,
 \qquad
 \Sigma_M:=\bigl(I+(C_M^{\rm obs})^*C_M^{\rm obs}\bigr)^{-1}
 \succeq\frac{I}{1+2\widetilde M_{s,\beta,\nu}^2}.
 \label{eq:ns-import-uniform-graph}
\end{equation}

\subsection{Memoria bidireccional y almacenamiento terminal}

Con
\begin{equation}
 r=\frac{\pi_{\rm HMT}}{\varphi_{\rm HMT}^2},
 \qquad
 D_{M,j}^{\rm mem}=(r-1)(M_{M,j}^+-M_{M,j}^-),
 \qquad
 H_{M,j}=(1+r)(M_{M,j}^++M_{M,j}^-),
 \label{eq:ns-import-dh}
\end{equation}
las dos memorias satisfacen
\begin{equation}
 \dot D_{M,j}^{\rm mem}=B_{s,M,j},
 \qquad
 \dot H_{M,j}=\frac{1+r}{r-1}|B_{s,M,j}|.
 \label{eq:ns-import-dh-dot}
\end{equation}
Poniendo
\[
 \delta_r=\frac{r-1}{r+1},
 \qquad
 \mathcal N_{M,j}:=\delta_rH_{M,j}-D_{M,j}^{\rm mem}
 =2(r-1)M_{M,j}^-\ge0,
\]
se obtiene
\begin{equation}
 \dot{\mathcal N}_{M,j}=|B_{s,M,j}|-B_{s,M,j}.
 \label{eq:ns-import-negative-memory}
\end{equation}
El almacenamiento terminal explícito es
\begin{equation}
 \begin{aligned}
 G_{M,j,T}(t)
 &=E_{s,M,j}(t)
   +2(r-1)\bigl(M_{M,j}^-(T)-M_{M,j}^-(t)\bigr),\\
 G_{M,j,T}(t)&\ge E_{s,M,j}(t),\\
 \dot G_{M,j,T}+\nu D_{s,M,j}+|B_{s,M,j}|
 &=F_{s,M,j}.
 \end{aligned}
 \label{eq:ns-import-terminal-storage}
\end{equation}
Si
\[
 h^-_{M,j;t,T}(\tau)
 :=\mathbf1_{[t,T]}(\tau)\sqrt{\dot M^-_{M,j}(\tau)},
\]
la columna
\begin{equation}
 \mathcal E_{M,j,T}^{\rm term}(t)\Xi
 :=2^{-1/2}\Lambda^su_{M,j}(t)
 \oplus\sqrt{2(r-1)}\,h^-_{M,j;t,T}
 \label{eq:ns-import-terminal-column}
\end{equation}
realiza el almacenamiento como un Gram:
\begin{equation}
 \|\mathcal E_{M,j,T}^{\rm term}(t)\Xi\|^2
 =G_{M,j,T}(t).
 \label{eq:ns-import-terminal-gram}
\end{equation}

El acarreo común entre profundidades es
\begin{equation}
 \kappa_{M,j}
 =\frac12\left(E_j|B_{s,M,j+1}|-|B_{s,M,j}|\right)\ge0,
 \qquad
 E_jB_{s,M,j+1}=B_{s,M,j}.
 \label{eq:ns-import-carry}
\end{equation}
Su coordenada de salida se fija por
\begin{equation}
 q_{M,j}^{\kappa}
 :=\left(\frac{2(1+r)}{r-1}\kappa_{M,j}\right)^{1/2},
 \qquad
 \|q_{M,j}^{\kappa}\|^2
 =\frac{2(1+r)}{r-1}\kappa_{M,j}.
 \label{eq:ns-import-carry-port}
\end{equation}
Al mismo tiempo,
\begin{equation}
 E_j\Delta_T\mathcal N_{M,j+1}
 =\Delta_T\mathcal N_{M,j}+2K_{M,j}(T),
 \qquad
 K_{M,j}(T):=\int_0^T\kappa_{M,j}(t)\,dt.
 \label{eq:ns-import-carry-telescope}
\end{equation}
Las innovaciones microscópicas y las bandas de Galerkin permanecen
ortogonales:
\begin{equation}
 Q_{\rm micro}^{\rm APP}=I-J_9E_9,
 \qquad E_9y^{\rm micro}=0,
 \qquad \mathsf K_{\rm carry}y^{\rm micro}=0,
 \label{eq:ns-import-micro}
\end{equation}
y, grado a grado,
\begin{equation}
 \mathcal N_{n,n+1;3M}J_M^{(n+1)}
 =J_M^{(n)}\mathcal N_{n,n+1;M}
 +\mathcal S_{n,M}^{\rm sh},
 \qquad
 E_M^{(n)}\mathcal S_{n,M}^{\rm sh}=0.
 \label{eq:ns-import-shell}
\end{equation}

\subsection{Control KYP finito no gobernante y cascada de la carta}

La fuerza y su término directo se incorporan mediante
\begin{equation}
 \begin{aligned}
 x_{M,j+1}
 &=\mathcal A_{M,j}^{[j]}x_{M,j}
   +\mathcal B_{M,j,T}^{F,[j]}g_{M,j},\\
 y_{M,j}^{\rm loss}
 &=\mathcal C_{M,j,T}^{[j]}x_{M,j}
   +\mathcal D_{M,j,T}^{F,[j]}g_{M,j},
 \end{aligned}
 \label{eq:ns-import-affine-step}
\end{equation}
con la descomposición ortogonal
\begin{equation}
 \mathcal C_{M,j,T}
 =\mathcal C_{M,j,T}^{(1)}
 \oplus\mathcal C_{M,j,T}^{(2)}
 \oplus\mathcal C_{M,j,T}^{\kappa}
 \oplus\mathcal C_{M,j,T}^{\rm micro}
 \oplus\mathcal C_{M,j,T}^{\rm sh}.
 \label{eq:ns-import-loss-decomposition}
\end{equation}
El término directo aparece ya en
\[
 g_M=A_M^{-1/2}\Lambda^sP_Mf,
 \qquad
 a_M=\frac{g_M}{2\sqrt\nu},
 \qquad
 y_{D,M}=\sqrt\nu A_M^{1/2}\Lambda^su_M-a_M,
\]
para el que
\begin{equation}
 dG_{M,j,T}
 +\bigl(\|y_{D,M}\|^2+|B_{s,M,j}|\bigr)\,dt
 =\|a_M\|^2\,dt.
 \label{eq:ns-import-force-feedthrough}
\end{equation}

La afinidad se homogeneiza sin duplicar el dato. Sea
\[
 \mathscr G_M^{F,[j]}
 :=\bigoplus_{k=j}^{J-1}L^2(\Omega_k;\mathscr G^F_{M,k,T}),
 \qquad
 \mathbf x_{M,j}:=x_{M,j}\oplus(g_{M,j},\ldots,g_{M,J-1}),
\]
y sean \(e_j\) la evaluación de la primera entrada y \(\mathsf S_j\) el
desplazamiento de la cola. Definimos
\begin{equation}
 \mathbf\Gamma_{M,j}
 :=\begin{pmatrix}
  \mathcal A_{M,j}^{[j]}&\mathcal B_{M,j,T}^{F,[j]}e_j\\
  0&\mathsf S_j
 \end{pmatrix},
 \qquad
 \mathbf L_{M,j,T}
 :=\mathcal Q_{M,j,T}^{\rm loss,inc}
 \begin{pmatrix}
  \mathcal C_{M,j,T}^{[j]}&\mathcal D_{M,j,T}^{F,[j]}e_j
 \end{pmatrix}.
 \label{eq:ns-import-homogeneous-affine-step}
\end{equation}
Aquí \(\mathcal Q_{M,0,T}^{\rm loss,inc}=I\); para \(j>0\) elimina
únicamente la copia heredada de los canales PC--Fock y actúa como la identidad
sobre acarreo, innovación microscópica y banda espectral.

\begin{theorem}[Control de factorización KYP finita no circular]
\label{thm:ns-import-affine-stein}
En el nivel terminal se define
\begin{equation}
 \mathbf U_{M,J,T}(t)
 :=(\mathcal E_{M,J,T}^{\rm term}(t))^*
   \mathcal E_{M,J,T}^{\rm term}(t).
 \label{eq:ns-import-terminal-metric}
\end{equation}
Fijados \(M,J,T\), sean \(A,B,C,D\) los cuatro bloques del paso afín y
\(U_+=\widetilde{\mathbf U}_{M,j+1,T}(t)\). Defínanse
\begin{equation}
 Q:=I-B^*U_+B-D^*D,
 \qquad L:=B^*U_+A+D^*C.
 \label{eq:ns-import-kyp-q-l}
\end{equation}
Si \(Q\succ0\), la recurrencia no circular es
\begin{equation}
 U:=A^*U_+A+C^*C+L^*Q^{-1}L.
 \label{eq:ns-import-affine-stein}
\end{equation}
Entonces la matriz KYP completa satisface
\begin{equation}
 \begin{aligned}
 \mathfrak S_{M,j,T}
 &:=
 \begin{pmatrix}
 U-A^*U_+A-C^*C&-A^*U_+B-C^*D\\
 -B^*U_+A-D^*C&I-B^*U_+B-D^*D
 \end{pmatrix}\\
 &=
 \begin{pmatrix}Q^{-1/2}L&-Q^{1/2}\end{pmatrix}^{*}
 \begin{pmatrix}Q^{-1/2}L&-Q^{1/2}\end{pmatrix}
 \succeq0.
 \end{aligned}
 \label{eq:ns-import-affine-supply}
\end{equation}
Para \(Q\succeq0\) vale la variante con pseudoinversa si
\(\operatorname{Ran}L\subseteq\operatorname{Ran}Q^{1/2}\). Esta construcción
es finita: no afirma que \(Q^{-1}\), \(U\) o la ganancia permanezcan
uniformes al retirar \(M,J\).
\end{theorem}

\begin{proof}
Por \eqref{eq:ns-import-affine-stein}, el bloque superior izquierdo es
\(L^*Q^{-1}L\); los bloques no diagonales son \(-L^*\) y \(-L\), y el bloque
inferior derecho es \(Q\). La multiplicación de la fila mostrada en
\eqref{eq:ns-import-affine-supply} da exactamente los cuatro bloques. La
positividad precede así a la toma de cualquier raíz.
\end{proof}

Este teorema comprueba únicamente la carta finita cuando se satisfacen sus
hipótesis propias. No aporta una premisa al owner Stein--Lyapunov ni limita la
uniformidad demostrada por su Gram terminal y su raíz común.

En las coordenadas homogeneizadas se toma
\(\mathbf U_{M,j,T}=U\) y se denota por
\(\mathbf R_{M,j,T}\) el levantamiento de la fila
\(\begin{psmallmatrix}Q^{-1/2}L&-Q^{1/2}\end{psmallmatrix}\).
Con esta convención se recupera, para cada truncación finita que satisface la
hipótesis KYP,
\[
 \mathbf U_{M,j,T}
 =\mathbf\Gamma_{M,j}^*\widetilde{\mathbf U}_{M,j+1,T}
  \mathbf\Gamma_{M,j}
 +\mathbf L_{M,j,T}^*\mathbf L_{M,j,T}
 +\mathbf R_{M,j,T}^*\mathbf R_{M,j,T}.
\]

La columna
\begin{equation}
 [\mathbf x]_{\mathbf U_{M,j,T}}
 \longmapsto
 [\mathbf\Gamma_{M,j}\mathbf x]_{
  \widetilde{\mathbf U}_{M,j+1,T}}
 \oplus\mathbf L_{M,j,T}\mathbf x
 \oplus\mathbf R_{M,j,T}\mathbf x
 \label{eq:ns-import-local-isometry}
\end{equation}
es isométrica. Su dilatación de Julia y su composición producen
\(\mathfrak U_{M,J,T}\) con
\begin{equation}
 \mathfrak U_{M,J,T}^{\sharp}\mathfrak U_{M,J,T}=I
 \label{eq:ns-import-cascade-isometry}
\end{equation}
y el telescopio
\begin{equation}
 \begin{aligned}
 \|\mathbf x_{M,0}\|_{\mathbf U_{M,0,T}}^2
 ={}&\mathbb E_{0:J}
 \|\mathbf x_{M,J}\|_{\mathbf U_{M,J,T}}^2\\
 &+\sum_{j=0}^{J-1}\mathbb E_{0:j}
 \left(\|\mathbf L_{M,j,T}\mathbf x_{M,j}\|^2
       +\|\mathbf R_{M,j,T}\mathbf x_{M,j}\|^2\right).
 \end{aligned}
 \label{eq:ns-import-stein-telescope}
\end{equation}

\subsection{Presupuesto de la carta y retorno de la velocidad}

Sobre la diagonal física, la primera coordenada terminal es
\(G_{M,J,T}\), y las demás coordenadas forman la suma de pérdidas:
\begin{equation}
 \|Z_{M,J,t}\|^2
 =\mathbb E_{0:J}G_{M,J,T}(t)
 +\sum_{j=0}^{J-1}\mathbb E_{0:j}\Lambda_{M,j}([0,t]).
 \label{eq:ns-import-ledger}
\end{equation}
Como \(q_{M,j}^{\kappa}\) es una de esas coordenadas,
\eqref{eq:ns-import-carry-telescope} se integra una sola vez:
\begin{equation}
 2\sum_{j=0}^{J-1}\mathbb E_{0:j}K_{M,j}(T)
 \le\|\Xi_{M,0}\|^2+Q_f(T).
 \label{eq:ns-import-carry-funded}
\end{equation}
De aquí,
\begin{equation}
 \mathbb E_{0:J}G_{M,J,T}(0)
 \le E_{s,M}(P_Mu_0)
 +\Delta_T\mathcal N_0+\|\Xi_{M,0}\|^2+Q_f(T).
 \label{eq:ns-import-root-bound}
\end{equation}
La raíz es común a todos los cortes:
\begin{equation}
 \Xi_{M,0}=J_{0\to M}\Xi_0,
 \qquad J_{0\to M}^*J_{0\to M}=I,
 \qquad \mathcal N_{M,0}=\mathcal N_0.
 \label{eq:ns-import-common-root}
\end{equation}
Además,
\begin{equation}
 \|\Xi_{M,0}\|^2
 \le C_{\rm HMT}
 +(1+2\widetilde M_{s,\beta,\nu}^2)
 e^{R_{\rm F}^2\|u_0\|_{H^s}^2},
 \qquad
 Q_f(T)\le C_{\nu,s}\int_0^T\|f(t)\|_{H^{s-1}}^2\,dt.
 \label{eq:ns-import-root-data-bound}
\end{equation}
Por tanto, para esta carta y estos \(M,J,T\), ponemos
\begin{align}
 \mathcal B_{M,J,T}^{\rm route}:={}&
 \frac12\|u_0\|_{H^s}^2
 +|\Delta_T\mathcal N_{M,0}|+C_{\rm HMT}
 +(1+2\widetilde M_{s,\beta,\nu}^2)
 e^{R_{\rm F}^2\|u_0\|_{H^s}^2}\notag\\
 &+C_{\nu,s}\int_0^T\|f(t)\|_{H^{s-1}}^2\,dt
 <\infty
 \label{eq:ns-import-explicit-budget}
\end{align}
La presencia explícita de \(\Delta_T\mathcal N_{M,0}\) impide usar esta
expresión como prueba independiente de uniformidad. Es un diagnóstico local
de la carta; no reemplaza el presupuesto owner
\eqref{eq:nsclose-owner-data-budget}. La identidad terminal y la desigualdad
de Young dan, para el corte fijo,
\begin{equation}
 \boxed{
 \sup_{0\le t\le T}\|Z_{M,J,t}\|^2
 +\frac{\nu}{2}\int_0^TD_{s,M}(t)\,dt
 \le\mathcal B_{M,J,T}^{\rm route}.}
 \label{eq:ns-import-uniform-budget}
\end{equation}

La publicación de la velocidad se normaliza por
\begin{equation}
 \iota_{s,M}u=2^{-1/2}\Lambda^su,
 \qquad
 \mathcal R_{1,M}=\sqrt2\,\Lambda^{-s}\pi_1,
 \qquad
 \mathcal R_{1,M}\iota_{s,M}=I_{V_M}.
 \label{eq:ns-import-grade-one-return}
\end{equation}
El lector de la cascada es
\begin{equation}
 \mathcal R_{M,J,T}^{\rm phys}
 :=\mathcal R_{1,M}\mathfrak U_{M,J,T}^{\sharp}.
 \label{eq:ns-import-full-return}
\end{equation}
Por \eqref{eq:ns-import-cascade-isometry},
\begin{equation}
 \boxed{
 \mathcal R_{M,J,T}^{\rm phys}
 \mathfrak U_{M,J,T}\mathcal V_M^{\rm PCF}(u_M(t))
 =u_M(t).}
 \label{eq:ns-import-exact-return}
\end{equation}
El retorno emplea el adjunto métrico de la misma cascada y recupera la
publicación hidrodinámica sin seleccionar una rama probabilística.

\section{Realización propietaria del estado enriquecido de Stein--Lyapunov}
\label{sec:ns-owner-closure}

La construcción de la \cref{sec:ns-pcf-prospective-construction} reúne en un
solo espacio de estado la torre PC--Fock completa, la memoria bidireccional,
el acarreo, la innovación microscópica, el canal espectral, la coordenada
homogénea de la fuerza, la unidad y la publicación excepcional. El tiempo
físico parametriza el almacenamiento terminal, mientras la profundidad
nonádica parametriza la recurrencia de Stein. La construcción que gobierna
esta sección es el \texttt{OWNER\_RECTOR} del estado enriquecido HMT. KYP,
NS--OBS y las cartas LRDN--LCN se conservan únicamente como
\texttt{CONTROL\_NO\_GOBERNANTE} de representaciones reducidas.

\subsection{Coligación PC--Fock y telescopía}

Sea
\[
 \mathbf x_{M,j}^{\rm own}
 \in\mathscr K_{M,j}^{\rm PCF}
 \widehat\otimes\mathcal E_{\rm exc}
 \oplus\mathscr G_M^{F,[j]}
\]
el estado homogéneo de la diagonal física. Para cada \(M,J,T\), la métrica
propietaria se construye prospectivamente como el Gram de la columna
terminal y de todas las pérdidas futuras de esa misma diagonal:
\begin{equation}
 \begin{aligned}
 \left\langle x,U_{M,j}^{\rm own}x\right\rangle
 :={}&
 \mathbb E_{j:J}
 \left\|\mathcal E_{M,J,T}^{\rm term}
       \Gamma_{M,j:J}^{\rm own}x\right\|^2\\
 &+\sum_{k=j}^{J-1}\mathbb E_{j:k}
 \left\|L_{M,k,T}^{\rm own}
       \Gamma_{M,j:k}^{\rm own}x\right\|^2 .
 \end{aligned}
 \label{eq:nsclose-owner-metric-definition}
\end{equation}
Todos los términos son cuadrados de las coordenadas ya construidas: grado
uno viscoso, grado dos convectivo, carry, innovación, archivo espectral,
memoria \(D/H\) y fuerza homogénea. Separar el primer paso de la suma da,
antes de introducir ninguna raíz KYP,
\begin{equation}
 \boxed{
 U_{M,j}^{\rm own}
 =
 (\Gamma_{M,j}^{\rm own})^*
 U_{M,j+1}^{\rm own}\Gamma_{M,j}^{\rm own}
 +(L_{M,j,T}^{\rm own})^*L_{M,j,T}^{\rm own},
 \qquad
 G_{M,J,T}^{\rm own}\preceq C_TU_{M,J}^{\rm own}.}
 \label{eq:nsclose-proprietary-terminal-gate}
\end{equation}
\label{eq:nsclose-owner-identity}
En la normalización terminal de \eqref{eq:ns-import-terminal-column} se puede
tomar \(C_T=1\). La igualdad de mapas, amplificada por la identidad en la capa
Moonshine, produce la cascada isométrica
\begin{equation}
 (\mathfrak U_{M,J,T}^{\rm own})^{\sharp}
 \mathfrak U_{M,J,T}^{\rm own}=I.
 \label{eq:nsclose-proprietary-full-cascade}
\end{equation}
Al iterar \eqref{eq:nsclose-owner-identity} se obtiene el telescopio
\begin{equation}
 U_{M,0}^{\rm own}
 =
 (\Gamma_{M,0:J}^{\rm own})^*U_{M,J}^{\rm own}
  \Gamma_{M,0:J}^{\rm own}
 +\sum_{j<J}(\Gamma_{M,0:j}^{\rm own})^*
  (L_{M,j,T}^{\rm own})^*L_{M,j,T}^{\rm own}
  \Gamma_{M,0:j}^{\rm own}.
 \label{eq:nsclose-owner-telescope}
\end{equation}
Su registro ortogonal de balance es
\begin{equation}
 \boxed{
 E_s(u_M(t))\le\|Z_{M,J,t}\|^2
 =\mathbb E_{0:J}G_{M,J,T}(t)
 +\sum_{j<J}\mathbb E_{0:j}
   \Lambda_{M,j}^{\rm own}([0,t]).}
 \label{eq:nsclose-proprietary-ledger}
\end{equation}
La raíz común
\(\Xi_{M,0}=J_{0\to M}\Xi_0\), la conservación evolutiva de la unidad y el
puerto de carry de \eqref{eq:ns-import-carry-port} hacen que el lado derecho
se pague una sola vez. El presupuesto propietario es
\begin{equation}
 \begin{aligned}
 \mathcal B_T^{\rm own}:={}&
 C_{\rm HMT}
 +(1+2\widetilde M_{s,\beta,\nu}^2)
  e^{R_{\rm F}^2\|u_0\|_{H^s}^2}\\
 &+C_{\nu,s}\int_0^T\|f(t)\|_{H^{s-1}}^2\,dt
 +\|\mathcal J_T^{\rm own}(u_0,f)\|^2 ,
 \end{aligned}
 \label{eq:nsclose-owner-data-budget}
\end{equation}
donde \(\mathcal J_T^{\rm own}\) es la inyección de raíz determinada por el
estado inicial transportado por TPK, sus dos hojas y la entrada forzada. Sus
amplificaciones son
\[
 \mathcal J_{M,J,T}^{\rm own}
 :=I_{0\to M,J}\mathcal J_T^{\rm own},
 \qquad I_{0\to M,J}^*I_{0\to M,J}=I.
\]
Por \eqref{eq:nsclose-owner-telescope}, no es una cantidad reconstruida desde
la órbita futura. La conservación de la unidad y la coisometría del lector
two-way dan, sobre el dominio declarado,
\begin{equation}
 \sup_{M,J}\|\mathcal J_{M,J,T}^{\rm own}(u_0,f)\|^2
 \le C_{{\rm HMT},T}
 \left(1+\|u_0\|_{H^{s+1}}^2
 +\int_0^T\|f(t)\|_{H^{s-1}}^2\,dt\right).
 \label{eq:nsclose-owner-root-uniformity}
\end{equation}
Así \(\mathcal B_T^{\rm own}\) es independiente de \(M,J\). El ledger da
\begin{equation}
 \boxed{
 \sup_{M,J}\left[
  \sup_{0\le t\le T}\|Z_{M,J,t}\|^2
  +\frac\nu2\int_0^TD_{s,M}(t)\,dt
 \right]\le\mathcal B_T^{\rm own}<\infty.}
 \label{eq:nsclose-proprietary-uniform-budget}
\end{equation}

La publicación hidrodinámica es el retorno de primer grado de la misma
cascada:
\begin{equation}
 R_{M,J,T}^{\rm full}
 :=\mathcal R_{1,M}
   (\mathfrak U_{M,J,T}^{\rm own})^{\sharp},
 \qquad
 \boxed{
 R_{M,J,T}^{\rm full}Z_{M,J,t}=u_M(t).}
 \label{eq:nsclose-proprietary-exact-return}
\end{equation}
\label{eq:nsclose-exact-full-return}
Esta igualdad coincide con \eqref{eq:ns-import-exact-return} y evita la
evaluación de una rama individual del refinamiento.

\begin{theorem}[Estimación uniforme de la realización propietaria PC--Fock]
\label{thm:nsclose-proprietary-full}
Para todo \(T<\infty\), las aproximaciones de Galerkin satisfacen
\begin{equation}
 \sup_M\sup_{0\le t\le T}\|u_M(t)\|_{H^s}^2
 +\nu\sup_M\int_0^T\|u_M(t)\|_{H^{s+1}}^2\,dt
 \le C_s\mathcal B_T^{\rm own},
 \label{eq:nsclose-uniform-sobolev-precompactness}
\end{equation}
\label{eq:nsclose-uniform-sobolev}
con una constante independiente de \(M\) y de la profundidad nonádica.
\end{theorem}

\begin{proof}
La identidad \eqref{eq:nsclose-proprietary-exact-return} y la contractividad
del retorno dan
\(\|u_M(t)\|_{H^s}^2\le C_s\|Z_{M,J,t}\|^2\). La equivalencia de
\(D_{s,M}\) con la norma \(H^{s+1}\) en el subespacio solenoidal de media
cero, junto con \eqref{eq:nsclose-proprietary-uniform-budget}, prueba la
estimación.
\end{proof}

\subsection{Control alternativo por métrica de ruta}

La carta que sigue comprueba el refinamiento en coordenadas normalizadas. No
construye ni condiciona \(U^{\rm own}\): omite la unidad, la publicación
excepcional y parte del canal afín que ya residen en
\eqref{eq:nsclose-owner-metric-definition}.

La carta de Ambivalencia aporta
\begin{equation}
 G_{\rm av}=\begin{pmatrix}2&-1\\-1&1\end{pmatrix},
 \qquad
 M_{\rm av}=\begin{pmatrix}3&-1\\1&0\end{pmatrix},
 \qquad
 M_{\rm av}^*G_{\rm av}M_{\rm av}
 \preceq\varphi_{\rm HMT}^4G_{\rm av}.
 \label{eq:nsclose-ambivalence-route-metric}
\end{equation}
Después de la normalización \(\widehat M_j=M_j/3\),
\begin{equation}
 \mathsf R_j:=G_j-\widehat M_j^*G_{j+1}\widehat M_j\succeq0.
 \label{eq:nsclose-route-exact-defect}
\end{equation}
En la carta estacionaria,
\[
 \mathsf R_{\rm av}
 =\frac19\begin{pmatrix}5&-4\\-4&7\end{pmatrix},
 \qquad
 \lambda_{\min}(\mathsf R_{\rm av})
 =\frac{6-\sqrt{17}}9.
\]
Con la métrica de grafo
\[
 \mathsf G_M^{\rm gr}
 =I+(\mathcal Y_{\beta,M})^*\mathcal Y_{\beta,M},
 \qquad
 \mathcal C_M
 =\mathcal Y_{\beta,M}(\mathsf G_M^{\rm gr})^{-1/2},
\]
definimos
\begin{equation}
 \mathcal A_{M,j}=\widehat M_j\otimes I,
 \qquad
 \mathcal O_{M,j}=\mathsf R_j^{1/2}\otimes\mathcal C_M,
 \qquad
 \mathcal D_{M,j}
 =\mathsf R_j^{1/2}\otimes(I-\mathcal C_M^*\mathcal C_M)^{1/2}.
 \label{eq:nsclose-owner-route-column}
\end{equation}
La identidad de Stein de ruta es
\begin{equation}
 \boxed{
 \mathsf U_{M,j}^{\rm PCF}
 -\mathcal A_{M,j}^*\mathsf U_{M,j+1}^{\rm PCF}\mathcal A_{M,j}
 =\mathcal O_{M,j}^*\mathcal O_{M,j}
  +\mathcal D_{M,j}^*\mathcal D_{M,j}.}
 \label{eq:nsclose-owner-stein}
\end{equation}
Su iteración da
\begin{equation}
 \begin{aligned}
 \mathsf U_{M,0}^{\rm PCF}
 -\mathcal A_{M,0:J}^*\mathsf U_{M,J}^{\rm PCF}\mathcal A_{M,0:J}
 =\sum_{j<J}\mathcal A_{M,0:j}^*
 \bigl(\mathcal O_{M,j}^*\mathcal O_{M,j}
      +\mathcal D_{M,j}^*\mathcal D_{M,j}\bigr)
 \mathcal A_{M,0:j}.
 \end{aligned}
 \label{eq:nsclose-owner-stein-telescope}
\end{equation}
Esta factorización describe la geometría del refinamiento. La recurrencia
afín \eqref{eq:ns-import-affine-stein} incorpora además la evolución física,
los términos cruzados de Carleman y la fuerza.

La memoria bidireccional convierte la pérdida de ruta en un presupuesto. Con
\(\mathcal N_{M,j}=2(r-1)M_{M,j}^-\),
\begin{equation}
 E_j\Delta_T\mathcal N_{M,j+1}
 =\Delta_T\mathcal N_{M,j}+2K_{M,j}(T),
 \qquad
 K_{M,j}(T)=\int_0^T\kappa_{M,j}(t)\,dt.
 \label{eq:nsclose-carry-memory-telescope}
\end{equation}
La telescopía y la cota de la raíz producen
\begin{equation}
 \mathbb E_{0:J}G_{M,J,T}(0)
 \le E_{s,M}(P_Mu_0)+|\Delta_T\mathcal N_0|
 +\|\Xi_{M,0}\|^2+Q_f(T),
 \label{eq:nsclose-owner-terminal-bound}
\end{equation}
que es la forma abreviada del control de ruta
\eqref{eq:ns-import-root-bound}--\eqref{eq:ns-import-explicit-budget}. Esta
carta registra el coste de la hoja negativa,
\[
 \Delta_T\mathcal N_{M,0}
 =2\int_0^T B^-_{s,M,0}(u_M(t))\,dt.
\]
En esta representación reducida, la desigualdad
\[
 X^-_{M,T}(0)\preceq C_T\mathsf U^{\rm PCF}_{M,0},
 \qquad \sup_M C_T<\infty
 \tag{NS--OBS}
\]
es un falsador útil de la carta. Su fallo no se transporta al owner, porque
\(\Delta_T\mathcal N_{M,0}\) no sustituye a la inyección de raíz
\(\mathcal J_T^{\rm own}\) de
\eqref{eq:nsclose-owner-data-budget}--\eqref{eq:nsclose-owner-root-uniformity}.
Análogamente, KYP y LRDN--LCN verifican completaciones o dominaciones de
lectores reducidos; ninguna de ellas es una premisa del
\cref{thm:nsclose-proprietary-full}.

\subsection{Existencia, unicidad y suavidad global}

\begin{theorem}[Regularidad global de Navier--Stokes]
\label{thm:nsclose-global-smooth}
\label{thm:realizaciones-ns}
Sea \(s>5/2\), \(\nu>0\), sea
\(u_0\in H^{s+1}_\sigma(\mathbb T^3)\) de media cero y sea
\(f\in L^2_{\rm loc}([0,\infty);H^{s-1}_\sigma)\). Existe una única solución
global
\begin{equation}
 u\in L^\infty_{\rm loc}([0,\infty);H^s_\sigma)
 \cap L^2_{\rm loc}([0,\infty);H^{s+1}_\sigma),
 \qquad
 \partial_tu\in L^2_{\rm loc}([0,\infty);H^{s-1}_\sigma).
 \label{eq:nsclose-global-regularity-class}
\end{equation}
Para dato y fuerza suaves, \(u\) y la presión normalizada son suaves.
\end{theorem}

\begin{proof}
El \cref{thm:nsclose-proprietary-full} da acotación uniforme de
\((u_M)_M\) en \(L^\infty(0,T;H^s)\cap L^2(0,T;H^{s+1})\). La ecuación de
Galerkin y la continuidad
\[
 H^s\times H^s\longrightarrow H^{s-1},
 \qquad (u,v)\longmapsto\mathbb P(u\cdot\nabla v),
\]
acotan \(\partial_tu_M\) en \(L^2(0,T;H^{s-1})\). Como
\(H^{s+1}\Subset H^s\hookrightarrow H^{s-1}\), Aubin--Lions produce, tras
extraer una subsucesión,
\[
 u_M\longrightarrow u\quad\hbox{fuertemente en }L^2(0,T;H^s).
\]
Esta convergencia permite pasar al límite en
\(\mathbb P(u_M\cdot\nabla u_M)\), mientras la convergencia débil conserva el
término viscoso y el dato inicial. Se obtiene una solución fuerte en
\([0,T]\).

Si \(u,v\) son dos soluciones y \(w=u-v\), la estimación de producto y
Young dan
\[
 \frac12\frac d{dt}\|w\|_{H^{s-1}}^2
 +\frac\nu2\|w\|_{H^s}^2
 \le C_{s,\nu}\bigl(\|u\|_{H^s}^2+\|v\|_{H^s}^2\bigr)
 \|w\|_{H^{s-1}}^2.
\]
Grönwall prueba unicidad. Al repetir la construcción en los horizontes
\(T=1,2,\ldots\), las soluciones coinciden en los intervalos solapados y se
obtiene la solución global.

La presión se recupera de
\begin{equation}
 -\Delta p=\partial_i\partial_j(u_i u_j)-\operatorname{div}f,
 \qquad \int_{\mathbb T^3}p\,dx=0.
 \label{eq:nsclose-pressure-recovery}
\end{equation}
La iteración de la estimación en \(s+k\) y la ecuación recuperan todas las
derivadas espaciales y temporales cuando \(u_0\) y \(f\) son suaves.
\end{proof}

La contribución estructural consiste en representar la convección como una
salida lineal de la torre PC--Fock, conservarla en la memoria \(D/H\) e
integrarla en una identidad ortogonal con retorno físico exacto. La
estimación uniforme resultante se publica después como regularidad global de
la ecuación hidrodinámica.


\section*{Refuerzo integrado del mismo propietario}

El módulo siguiente conserva el cálculo uniforme de las filas de Fourier,
la inyección raíz, la derivada temporal, la presión y el pegado de
horizontes. Su control KYP finito queda separado de la cadena gobernante.

\section*{Balance solenoidal, acarreo y terminal enriquecido}

Este módulo refuerza el \texttt{OWNER\_RECTOR} Navier--Stokes: conserva en
una sola cadena el acarreo de Galerkin, el balance causal, el Gram terminal,
Stein, la telescopía, la inyección raíz, el retorno y el paso global. La carta
KYP aparece sólo después y queda fuera de esa cadena.

Para los cortes de Galerkin \(m\leq n\), el acarreo conservado es
\begin{equation}\label{eq:amp-ns-owner-acarreo-galerkin}
 C_{m\leftarrow n}
 :=P_mB(u_n,u_n)-B_m(P_m u_n,P_m u_n).
\end{equation}
La lectura nonádica de este estado está tipada por
\begin{equation}\label{eq:amp-ns-owner-e9j9}
 E_9J_9=I,\qquad P_9=J_9E_9,\qquad Q_{\rm micro}=I-P_9,
\end{equation}
y conserva el balance causal
\begin{equation}\label{eq:amp-ns-owner-balance-causal}
 \boxed{
 \mathbb E_9\|\Xi_{m+1}\|^2+\mathbb E_9\ell_m
 =\|\Xi_m\|^2,\qquad \ell_m\geq0.}
\end{equation}
El campo, la proyección solenoidal, el acarreo, la frontera, la ruta y las
pérdidas positivas permanecen así en un mismo estado antes de construir el
almacenamiento.

\section*{Gram truncado del estado completo e inyección raíz uniforme}

La realización que gobierna el balance no comienza por la desigualdad KYP.
Sean \(\Gamma_{M,j}^{\rm term}\) la transición de un paso del estado
PC--Fock enriquecido, \(\Gamma_{M,j:k}^{\rm term}\) su composición entre
las profundidades \(j\) y \(k\), y
\(L_{M,k,T}^{\rm term}\) la columna de pérdidas de grado uno, grado dos,
acarreo, innovación, archivo y memoria en el horizonte \([0,T]\). La
notación \(\mathbb E_{j:k}\) designa la suma ortogonal con los pesos de
las ramas del refinamiento entre esos niveles. Para \(j\leq J\) se define
la columna finita
\begin{equation}\label{eq:amp-ns-owner-columna-gram-truncado}
 \begin{aligned}
 \mathbf G_{M,j;J,T}^{\rm term}x
 :={}&
 \bigl(\mathcal E_{M,J,T}^{\rm term}
       \Gamma_{M,j:J}^{\rm term}x\bigr)_{\mathbb E_{j:J}}\\
 &\oplus
 \bigoplus_{k=j}^{J-1}
 \bigl(L_{M,k,T}^{\rm term}
       \Gamma_{M,j:k}^{\rm term}x\bigr)_{\mathbb E_{j:k}},
 \qquad
 U_{M,j;J,T}^{\rm term}
 :=(\mathbf G_{M,j;J,T}^{\rm term})^*
    \mathbf G_{M,j;J,T}^{\rm term}.
 \end{aligned}
\end{equation}
En particular, el almacenamiento es un Gram truncado construido con el
terminal y las pérdidas futuras del mismo estado. Su forma cuadrática es
\begin{equation}\label{eq:amp-ns-owner-gram-truncado}
 \begin{aligned}
 \langle x,U_{M,j;J,T}^{\rm term}x\rangle
 ={}&\mathbb E_{j:J}
 \|\mathcal E_{M,J,T}^{\rm term}
       \Gamma_{M,j:J}^{\rm term}x\|^2\\
 &+\sum_{k=j}^{J-1}\mathbb E_{j:k}
 \|L_{M,k,T}^{\rm term}
       \Gamma_{M,j:k}^{\rm term}x\|^2 .
 \end{aligned}
\end{equation}

\begin{lemma}[Identidad de Stein del Gram truncado]
\label{lem:amp-ns-owner-stein-truncado}
Para todo \(j<J\),
\begin{equation}\label{eq:amp-ns-owner-stein-truncado}
 U_{M,j;J,T}^{\rm term}
 =(\Gamma_{M,j}^{\rm term})^*
   U_{M,j+1;J,T}^{\rm term}\Gamma_{M,j}^{\rm term}
  +(L_{M,j,T}^{\rm term})^*L_{M,j,T}^{\rm term}.
\end{equation}
Por iteración se obtiene, conservando el sumando terminal,
\begin{equation}\label{eq:amp-ns-owner-telescopio-truncado}
 \begin{aligned}
 U_{M,0;J,T}^{\rm term}
 ={}&(\Gamma_{M,0:J}^{\rm term})^*
      U_{M,J;J,T}^{\rm term}\Gamma_{M,0:J}^{\rm term}\\
 &+\sum_{j=0}^{J-1}
   (\Gamma_{M,0:j}^{\rm term})^*
   (L_{M,j,T}^{\rm term})^*L_{M,j,T}^{\rm term}
   \Gamma_{M,0:j}^{\rm term}.
 \end{aligned}
\end{equation}
\end{lemma}

\begin{proof}
En \eqref{eq:amp-ns-owner-columna-gram-truncado}, se separa el sumando
\(k=j\). Los restantes sumandos, incluido el terminal, son exactamente la
columna de nivel \(j+1\) precedida por
\(\Gamma_{M,j}^{\rm term}\). La ortogonalidad de la suma y la propiedad de
torre de \(\mathbb E_{j:k}\) dan
\eqref{eq:amp-ns-owner-stein-truncado}. La repetición del mismo
desdoblamiento da \eqref{eq:amp-ns-owner-telescopio-truncado}; no se
descarta
\((\Gamma_{M,0:J}^{\rm term})^*U_{M,J;J,T}^{\rm term}
\Gamma_{M,0:J}^{\rm term}\).
\end{proof}

\begin{theorem}[Inyección raíz y presupuesto uniforme]
\label{thm:amp-ns-inyeccion-raiz-uniforme}
La raíz del realizador se transporta a cada corte y profundidad mediante
una inclusión isométrica:
\begin{equation}\label{eq:amp-ns-inyeccion-raiz-comun}
 \begin{aligned}
 x_{M,0}^{\rm data}
 &:=\mathcal V_M^{\rm full}(P_Mu_0)\oplus g_M^F,\\
 \mathbf G_{M,0;J,T}^{\rm term}x_{M,0}^{\rm data}
 &=\mathcal J_{M,J,T}^{\rm term}(u_0,f),\\
 \mathcal J_{M,J,T}^{\rm term}
 &:=I_{0\to M,J}\mathcal J_T^{\rm term},
 & I_{0\to M,J}^*I_{0\to M,J}&=I.
 \end{aligned}
\end{equation}
Para los datos de \eqref{eq:realizaciones-ns-dominio}, existe
\(C_{{\rm HMT},T}\), independiente de \(M\) y \(J\), tal que
\begin{equation}\label{eq:amp-ns-inyeccion-raiz-cota}
 \sup_{M,J}
 \|\mathcal J_{M,J,T}^{\rm term}(u_0,f)\|^2
 \leq C_{{\rm HMT},T}
 +(1+2\widetilde M_{s,\beta,\nu}^{\,2})
  e^{R_{\rm F}^{\,2}\|u_0\|_{H^s}^2}
 +C_{\nu,s}\int_0^T\|f(t)\|_{H^{s-1}}^2\,dt .
\end{equation}
Equivalentemente, la cota es una cota del Gram terminal evaluado en
los datos:
\begin{equation}\label{eq:amp-ns-owner-gram-cota-uniforme}
 \sup_{M,J}
 \left\langle x_{M,0}^{\rm data},
 U_{M,0;J,T}^{\rm term}x_{M,0}^{\rm data}\right\rangle
 =
 \sup_{M,J}
 \|\mathcal J_{M,J,T}^{\rm term}(u_0,f)\|^2
 <\infty .
\end{equation}
La aplicación \(\mathcal J_T^{\rm term}\) depende sólo del estado inicial,
de sus dos hojas, del terminal heredado y del puerto de fuerza; no recibe
como argumento \(u_M(t)\) para \(0<t\leq T\).
\end{theorem}

\begin{proof}
La componente lineal y las memorias \(D/H\) se transportan
isométricamente. En la componente PC--Fock, la columna coherente es
\[
 \bigoplus_{n\geq0}
 \frac{R_{\rm F}^{\,n}}{\sqrt{n!}}\,u_0^{\odot n},
 \qquad
 \sum_{n\geq0}\frac{R_{\rm F}^{\,2n}}{n!}
       \|u_0\|_{H^s}^{2n}
 =e^{R_{\rm F}^{\,2}\|u_0\|_{H^s}^2}.
\]
Para justificar la uniformidad del lector cuadrático, se polariza
\begin{equation}\label{eq:amp-ns-lector-cuadratico-polarizado}
 \mathcal B_M(u,v):=
 \frac{1}{4\sqrt{\beta\nu}}\,
 A_M^{-1/2}\Lambda^sP_M\mathbb P
 \bigl((P_Mu\!\cdot\!\nabla)P_Mv
      +(P_Mv\!\cdot\!\nabla)P_Mu\bigr).
\end{equation}
En las bases transversales
\(e^s_{p,a}=\langle p\rangle^{-s}a\,e^{ip\cdot x}\), si
\(k=p+q\ne0\), sus filas de Fourier cumplen
\begin{equation}\label{eq:amp-ns-filas-fourier-uniformes}
 \|c^M_{k;p,q}\|
 \leq\frac{C}{\sqrt{\beta\nu}}\,
 \frac{\langle k\rangle^s(|p|+|q|)}
 {|k|\langle p\rangle^s\langle q\rangle^s},
 \qquad
 \sup_{M,k\ne0}\sum_{p+q=k}\|c^M_{k;p,q}\|^2<\infty .
\end{equation}
La segunda cota se obtiene separando \(|p|\leq2|k|\) de
\(|p|>2|k|\); en la cola domina
\(\sum_p\langle p\rangle^{2-2s}<\infty\) porque \(s>5/2\).
Cauchy--Schwarz en cada fila y la ortogonalidad de los bloques con distinto
\(k\) prolongan \(\mathcal B_M\) a
\[
 \mathfrak C_M:H^s_\sigma\odot_2H^s_\sigma\longrightarrow L^2_\sigma,
 \qquad
 \sup_M\|\mathfrak C_M\|
 \leq\widetilde M_{s,\beta,\nu}.
\]
Así, la constante usada en la raíz controla la suma de todos los pares
Fourier y no sólo cada coeficiente por separado. El puerto afín verifica
\[
 Q_f(T)\leq C_{\nu,s}
 \int_0^T\|f(t)\|_{H^{s-1}}^2\,dt .
\]
Por suma ortogonal de esas componentes,
\begin{equation}\label{eq:amp-ns-raiz-cota-por-componentes}
 \|\Xi_{M,0}\|^2
 \leq C_{{\rm HMT},T}
 +(1+2\widetilde M_{s,\beta,\nu}^{\,2})
 e^{R_{\rm F}^{\,2}\|u_0\|_{H^s}^2}.
\end{equation}
Las inclusiones de \eqref{eq:amp-ns-inyeccion-raiz-comun} no cambian esta
norma. Al sumar la cota del puerto se obtiene exactamente
\eqref{eq:amp-ns-inyeccion-raiz-cota}. Como todos los términos se evalúan
en \((u_0,f)\), la estimación precede a la trayectoria y no presupone la
cota global que después se deduce de
\eqref{eq:amp-ns-owner-telescopio-truncado}.
\end{proof}

El Gram terminal \eqref{eq:amp-ns-owner-gram-truncado}, su identidad de
Stein, la telescopía con terminal y la inyección raíz uniforme forman la
construcción propietaria. Ésta es la rama que gobierna el cierre y coincide
con la realización íntegra restaurada en
\cref{sec:ns-owner-closure,thm:nsclose-proprietary-full}.

\section*{Separación del control KYP finito no gobernante}

La factorización KYP finita correcta del
\cref{thm:ns-import-affine-stein} se conserva como control alternativo de
una carta reducida y no constituye una premisa del owner
\(U^{\rm own}\). No se utiliza la mutación firmada que pretendía
incorporar \(C_{m,M}^*C_{m,M}\) a un término \(P\succeq0\) y
recuperarlo después con signo negativo: al entrar en \(P\), ese
sumando aparece necesariamente con signo positivo. Por tanto, esa
factorización no interviene en el balance, el presupuesto raíz, el retorno
físico ni la conclusión Navier--Stokes.

\section*{Derivada temporal, presión y pegado de horizontes}

El tramo siguiente es consecuencia de la cota uniforme de trayectoria
obtenida por el Gram terminal propietario, su telescopía y el retorno
físico exacto. Se
dispone de la cota conjunta uniforme en
\(L^\infty(0,T;H^s)\cap L^2(0,T;H^{s+1})\).
La ecuación proyectada se escribe
\begin{equation}\label{eq:amp-ns-proyectada}
 \partial_tu_m
 =\nu\Delta u_m-P_mB(u_m,u_m)+P_mf.
\end{equation}
Para \(s>5/2\), el producto de Sobolev y la cota conjunta en
\(L^\infty(0,T;H^s)\cap L^2(0,T;H^{s+1})\) dan
\begin{equation}\label{eq:amp-ns-dt}
 \sup_m\|\partial_tu_m\|_{L^2(0,T;H^{s-1})}
 \leq C_T(u_0,f,\Xi_0).
\end{equation}
La estimación \eqref{eq:amp-ns-dt} y la compacidad espacial producen
convergencia fuerte en \(L^2(0,T;H^s)\), suficiente para pasar al término
bilineal sin perder el carry de Galerkin.

Recuperada la velocidad, la presión normalizada de media cero queda fijada
por
\begin{equation}\label{eq:amp-ns-presion}
 -\Delta p
 =\partial_i\partial_j(u_i u_j)-\nabla\!\cdot f,
 \qquad \int_{\mathbb T^3}p(t,x)\,dx=0.
\end{equation}
La presión es así una salida de la misma solución solenoidal; no es una
variable libre usada para seleccionar la rama regular.

\begin{proposition}[Pegado global por unicidad]
\label{prop:amp-ns-pegado}
Sean \(u^{(T)}\) las soluciones obtenidas en \([0,T]\) por el límite de
Galerkin. Si \(T_1<T_2\), entonces
\(u^{(T_2)}|_{[0,T_1]}=u^{(T_1)}\). La familia define una única solución
global suave para \(t>0\).
\end{proposition}

\begin{proof}
Ambas restricciones poseen los mismos datos y pertenecen al dominio de la
estimación uniforme. La ecuación de su diferencia, el control bilineal en
\(H^{s-1}\) y Grönwall implican igualdad en \([0,T_1]\). Los intervalos
finitos forman un sistema dirigido; la igualdad en los solapamientos define
la solución global. La regularización parabólica del teorema se aplica en
cada intervalo \([\varepsilon,T]\), \(\varepsilon>0\).
\end{proof}

El control de refutación es preciso: una pérdida de uniformidad en
\eqref{eq:amp-ns-dt}, dos límites diferentes en un mismo horizonte o una
presión que no satisfaga \eqref{eq:amp-ns-presion} romperían la cadena antes
de la conclusión global.

\paragraph{Conclusión propietaria.}
El Gram terminal propietario, Stein, la telescopía con terminal, su ledger,
el retorno físico exacto y la inyección raíz producen una cota \(H^s\) uniforme para
\(s>5/2\). Por tanto se demuestra la existencia, unicidad y regularidad
global de Navier--Stokes tridimensional periódico en el dominio HMT
registrado. Las cartas KYP finitas permanecen como
\texttt{CONTROL\_NO\_GOBERNANTE}.


\paragraph{Conclusión afirmativa.}
En el dominio periódico y Sobolev declarado, el almacenamiento terminal
\(U^{\rm own}\), su identidad de Stein, la telescopía que conserva el
terminal, el ledger que domina \(E_s(u_M)\), la raíz uniforme y el retorno
físico exacto producen la cota uniforme. Aubin--Lions, unicidad y
regularidad parabólica dan la solución global suave. Las cartas KYP,
NS--OBS y LRDN--LCN son controles alternativos y no premisas del teorema.

\chapter[Proceso gauge euclídeo Yang--Mills en dimensión 4]{Construcción
proyectiva de un proceso gauge euclídeo Yang--Mills, reconstrucción de
Osterwalder--Schrader y brecha colectiva en dimensión \(4\)}
\chaptermark{Proceso gauge euclídeo Yang--Mills}

\noindent La especialización del
\cref{thm:nucleo-continuidad-conexion-nonadica-conjunta} es
\[
 (\widehat X_\infty,\Gamma_9)
 \xrightarrow{\ \mathcal R_{\rm YM}\ }
 (\mu_m,T_m,H_m,\Omega_m)
 \longrightarrow(\mu_\infty,\mathrm{OS},E(4))
 \longrightarrow\operatorname{gap}(H_{\rm phys})>0.
\]
El lector clover que publica \(F^2\) reconoce la curvatura sobre este mismo
proceso proyectivo; no construye retrospectivamente \(\mu_m\). La teoría
gauge, la reconstrucción OS, el vacío simple y la brecha ya quedan fijados
por esta cadena. La identidad
\eqref{eq:amp-ym-identidad-accion-requerida} compara después una
parametrización opcional por \(g_R\); es un control no gobernante y no una
condición para identificar la teoría Yang--Mills construida.

\section*{Objeto gauge y dominio de refinamiento}

Sea \(G\) un grupo de Lie compacto, conectado y simple. Para cada
\(m\geq0\), se considera la red euclídea
\begin{equation}\label{eq:realizaciones-ym-red}
 \Lambda_m=a_m\mathbb Z^4,
 \qquad a_{m+1}=\frac{a_m}{9}.
\end{equation}
La álgebra cilíndrica gauge se genera por variables de enlace en \(G\),
operadores de entrelazamiento de representación y las coordenadas
enriquecidas de color,
orientación, holonomía, acarreo de fusión, ruta y memoria. Sea \(\mu_m\) la
medida invariante común y sea
\[
 \mathcal H_m=L^2(\mathcal A_m,\mu_m)^{\rm gauge}.
\]
El vector constante \(\Omega_m\) define
\begin{equation}\label{eq:realizaciones-ym-proyectores}
 P_{\Omega_m}=|\Omega_m\rangle\langle\Omega_m|,
 \qquad Q_m^{\rm phys}=I-P_{\Omega_m}.
\end{equation}
El segundo proyector es el complemento físico del vacío; no se identifica
con un residual empleado únicamente para organizar la escala de
renormalización.

El circuito de actualización es local:
\begin{equation}\label{eq:realizaciones-ym-circuito}
 K_m^{\rm loc}
 =\prod_{r=1}^{9\cdot81}
   \prod_{B\in\mathcal B_{m,r}}K_{m,B}.
\end{equation}
La partición \(\{\mathcal B_{m,r}\}_r\) enumera cada coordenada una sola vez; los factores
del circuito se identificarán abajo con el semigrupo de sus expectativas
condicionales, y no con nueve dinámicas independientes.
Los bloques de un mismo color poseen soportes disjuntos y el diámetro físico
de cada soporte está uniformemente acotado al variar \(m\). Esta localidad
forma parte del objeto, de modo que el refinamiento no convierte una
actualización elemental en una operación instantáneamente global.

Sea \(T_m\) la transferencia inducida por
\eqref{eq:realizaciones-ym-circuito}. Las inclusiones isométricas
\(J_m:\mathcal H_m\to\mathcal H_{m+1}\) conservan ancestro,
representación, orientación, acarreo y memoria. Para las cuatro reflexiones
euclídeas \(\Theta_{\nu,m}\), \(\nu=1,\ldots,4\), se cumple
\begin{equation}\label{eq:realizaciones-ym-entrelazamiento}
 (T_{m+1})^9J_m=J_mT_m,
 \qquad
 \Theta_{\nu,m+1}J_m=J_m\Theta_{\nu,m}
 \quad(\nu=1,\ldots,4).
\end{equation}
Las cuatro direcciones pertenecen, por tanto, al mismo refinamiento, la misma
medida y la misma dinámica.

\section*{Semigrupo común, corona relativa y cuatro formas de reflexión}

En la carta triangular del estado completo, sean \(E_{m,c}\) las esperanzas
condicionales ortogonales y conmutantes asociadas a coordenadas
independientes. El generador positivo de solapamiento se construye primero en
un nivel semilla \(m_0\):
\begin{equation}\label{eq:realizaciones-ym-generador}
 \widehat H_{m_0}^{\rm ov}
 =3\kappa_{\rm av}
  \sum_{c\in\mathcal I_{m_0}^{\rm seed}}(I-E_{m_0,c}),
\end{equation}
y se prolonga mediante
\begin{equation}\label{eq:realizaciones-ym-generador-refinado}
 \begin{aligned}
 \widehat H_{m+1}^{\rm ov}
 &=\widehat\Gamma_m^*
   \bigl(\widehat H_m^{\rm ov}\otimes I+
         I\otimes\widehat H_{m+1}^{\rm det}\bigr)
   \widehat\Gamma_m,\\
 \widehat H_{m+1}^{\rm det}
 &=3\kappa_{\rm av}
   \sum_{\iota\in\mathcal I_{m+1/m}^{\rm det}}
   (I-E_{m+1,\iota}^{\rm det}),
 \end{aligned}
\end{equation}
donde
\begin{equation}\label{eq:realizaciones-ym-kappa}
 \kappa_{\rm av}
 =\frac{\mu_{\rm vol}}{\ell_0}
  \left(\frac{\pi_{\rm HMT}}{\varphi_{\rm HMT}^{\,2}}-1\right)>0.
\end{equation}
La frecuencia volumétrica \(\mu_{\rm vol}\) es adimensional; la dimensión
procede de \(\ell_0\).

El semigrupo, la transferencia de un paso físico y su compatibilidad de
escala se definen con tipos visibles:
\begin{equation}\label{eq:realizaciones-ym-semigrupo}
 \widehat K_{m,t}^{\rm ov}=e^{-t\widehat H_m^{\rm ov}},
 \qquad
 \widehat T_m=\widehat K_{m,a_m}^{\rm ov},
 \qquad a_{m+1}=a_m/9,
\end{equation}
Como la partición enumera cada coordenada una sola vez y las expectativas
condicionales de sus bloques conmutan, para cada bloque y para el circuito
completo se obtiene
\begin{equation}\label{eq:realizaciones-ym-circuito-semigrupo}
 K_{m,B}=e^{-3\kappa_{\rm av}a_m(I-E_{m,B})},
 \qquad
 K_m^{\rm loc}=e^{-a_m\widehat H_m^{\rm ov}}=\widehat T_m.
\end{equation}
\begin{equation}\label{eq:realizaciones-ym-semigrupo-entrelazado}
 \widehat H_{m+1}^{\rm ov}\widehat J_m
 =\widehat J_m\widehat H_m^{\rm ov},
 \qquad
 (\widehat T_{m+1})^9\widehat J_m
 =\widehat J_m\widehat T_m.
\end{equation}
Así, la novena potencia se compara mediante \(\widehat J_m\); no se igualan
operadores que actúan en Hilbert distintos.

La medida común de ocho caras procede de ese mismo kernel reversible:
\begin{equation}\label{eq:realizaciones-ym-medida-ocho}
 \widehat\Omega_m^{(8)}(dy_1\cdots dy_8)
 =\int\widehat\mu_m^{\rm ov}(dz)
   \prod_{r=1}^{8}
   \widehat K_{m,a_m/2}^{\rm ov}(z,dy_r).
\end{equation}
Para cada reflexión \(\nu=1,\ldots,4\) y todo observable cilíndrico \(F\)
en el semiespacio positivo,
\begin{equation}\label{eq:realizaciones-ym-os}
 \int\overline{F(\Theta_{\nu,m}y)}F(y)\,
 d\widehat\Omega_m^{(8)}
 =\|\widehat{\mathcal B}_{m,\nu}F\|^2\geq0,
\end{equation}
y la marginal de caras opuestas da la identidad de enlace
\begin{equation}\label{eq:realizaciones-ym-os-transferencia}
 \widehat T_{m,\nu}^{\rm OS}
 =\widehat K_{m,a_m}^{\rm ov}
 =e^{-a_m\widehat H_m^{\rm ov}}
 =\widehat T_m,
 \qquad \nu=1,\ldots,4.
\end{equation}
Las cuatro reflexiones publican, por tanto, el mismo par
transferencia--medida.

\section*{Persistencia celular del terminal y brecha exacta}

Sean
\[
 K_9=C_*(Q_{9,3};\mathbb Z),\qquad
 K_{10}=C_*(Q_{10,3};\mathbb Z).
\]
La corona relativa tiene dimensiones
\[
 C_*(Q_{10,3},Q_{9,3})=(271,756,702,217),
 \qquad 271+702=756+217=973.
\]
La inclusión \(j:K_9\to K_{10}\), el colapso celular
\(r:K_{10}\to K_9\) y el homotopio integral \(H_{\rm cell}\) satisfacen
\begin{equation}\label{eq:realizaciones-ym-sdr-corona}
 rj=I,\qquad
 \partial H_{\rm cell}+H_{\rm cell}\partial=I-jr,\qquad
 rH_{\rm cell}=0,\quad H_{\rm cell}j=0,\quad H_{\rm cell}^2=0.
\end{equation}
Para todo mapa de cadenas
\(F:K_{10}\otimes V_{q=6}\to\mathcal T\), con
\(P=(jr)\otimes I_{q=6}\),
\begin{equation}\label{eq:realizaciones-ym-homotopia-terminal}
 F-FP
 =d_{\mathcal T}\bigl(F(H_{\rm cell}\otimes I)\bigr)
  +F(H_{\rm cell}\otimes I)(\partial\otimes I).
\end{equation}
La corona completa es así nulhomotópica y el terminal \(FP\) permanece
literalmente. Los \(271\) vértices no sustituyen las aristas, caras y cubos
de la corona; el factor celular tridimensional tampoco se identifica con el
factor gauge físico \(q=6\).

Para cada coordenada persistente \(c\), en la carta local
\(q=(q_1,\ldots,q_6)\in\mathbb F_3^6\), se fija el testigo centrado
\begin{equation}\label{eq:realizaciones-ym-testigo-definicion}
 \bigcap_c\operatorname{Ran}E_{m,c}=\mathbb C\Omega_m,
 \qquad
 W_c(q)=\mathbf 1_{\{q_1+\cdots+q_6=0\}}-\frac13,
\end{equation}
donde la suma del indicador se toma en \(\mathbb F_3\). La descomposición de
Hoeffding de las expectativas conmutantes prueba
\begin{equation}\label{eq:realizaciones-ym-gap-finito}
 \ker\widehat H_m^{\rm ov}=\mathbb C\Omega_m,
 \qquad
 \widehat H_m^{\rm ov}\succeq
 3\kappa_{\rm av}(I-P_{\Omega_m}).
\end{equation}
La igualdad no se deduce sólo de que el núcleo sea trivial: cada sector no
constante de la descomposición recibe al menos un proyector
\(I-E_{m,c}\) y, por \eqref{eq:realizaciones-ym-generador}, su coste mínimo
es \(3\kappa_{\rm av}\). El testigo material \(W_c\) satisface
\begin{equation}\label{eq:realizaciones-ym-testigo-gap}
 \widehat H_m^{\rm ov}W_c=3\kappa_{\rm av}W_c,
 \qquad \|W_c\|^2=\frac29,
\end{equation}
de modo que el complemento es no vacío y la brecha es exactamente
\(3\kappa_{\rm av}\).

La escala se publica por
\begin{equation}\label{eq:realizaciones-ym-qm}
 q_m=e^{-3\kappa_{\rm av}a_m},
 \qquad q_{m+1}^{\,9}=q_m,
 \qquad -a_m^{-1}\log q_m=3\kappa_{\rm av}.
\end{equation}
Aunque \(q_m\to1\) al refinar, el cociente espectral por unidad física
permanece constante. La contracción celular
\eqref{eq:realizaciones-ym-homotopia-terminal} conserva el terminal; la
brecha procede del generador positivo de solapamiento y de su descomposición de
Hoeffding, no de la aciclicidad por sí sola.

Las inclusiones de \eqref{eq:realizaciones-ym-semigrupo-entrelazado}
forman el límite inductivo
\begin{equation}\label{eq:realizaciones-ym-colimite}
 \mathcal H_\infty=\varinjlim(\mathcal H_m,\widehat J_m).
\end{equation}
Sobre la unión inductiva algebraica se define la forma
\begin{equation}\label{eq:realizaciones-ym-forma-limite}
 \mathfrak h_\infty[\widehat J_{m,\infty}\psi]
 :=\langle\psi,\widehat H_m^{\rm ov}\psi\rangle_{\mathcal H_m}.
\end{equation}
El entrelazamiento de los generadores hace que esta definición sea
independiente del representante. La forma es positiva y cerrable; su
clausura determina el generador autoadjunto \(\widehat H_\infty^{\rm ov}\).
La compatibilidad de las formas transporta
\begin{equation}\label{eq:realizaciones-ym-gap-forma-limite}
 \overline{\mathfrak h}_\infty[\psi]
 \geq3\kappa_{\rm av}
 \|(I-P_{\Omega_\infty})\psi\|^2,
\end{equation}
y el vector persistente
\(\widehat J_{m,\infty}W_c\) alcanza la igualdad. En consecuencia, la
brecha exacta del generador límite es \(3\kappa_{\rm av}\).

\section*{Hipótesis exactas y encadenamiento del teorema}

La publicación gauge utiliza:
\begin{enumerate}
 \item el grupo compacto, conectado y simple y la red cuatro-dimensional
       \eqref{eq:realizaciones-ym-red};
 \item el circuito local \eqref{eq:realizaciones-ym-circuito}, con diámetro
       físico uniforme, y una medida invariante común;
 \item las inclusiones isométricas y las cuatro reflexiones entrelazadas por
       \eqref{eq:realizaciones-ym-entrelazamiento} y
       \eqref{eq:realizaciones-ym-semigrupo-entrelazado};
 \item las cuatro factorizaciones
       \eqref{eq:realizaciones-ym-os} y el enlace exacto
       \eqref{eq:realizaciones-ym-os-transferencia} sobre el mismo par
       transferencia--medida;
 \item el generador de expectativas conmutantes
       \eqref{eq:realizaciones-ym-generador}--
       \eqref{eq:realizaciones-ym-generador-refinado} y su descomposición de
       Hoeffding;
 \item el SDR integral de la corona y la conservación del terminal
       \eqref{eq:realizaciones-ym-sdr-corona}--
       \eqref{eq:realizaciones-ym-homotopia-terminal};
 \item el terminal \(P_{\Omega_m}\), el complemento físico
       \(Q_m^{\rm phys}\) y la ley de escala
       \eqref{eq:realizaciones-ym-qm};
 \item el límite inductivo y la clausura compatible de formas
       \eqref{eq:realizaciones-ym-colimite}--
       \eqref{eq:realizaciones-ym-gap-forma-limite};
 \item las cartas internas de acción y velocidad
       \(\hbar_{\rm int}\), \(c_{\rm int}\), usadas sólo para publicar la
       dimensión de masa.
\end{enumerate}

\begin{theorem}[Teoría gauge local con vacío simple y brecha colectiva]
\label{thm:realizaciones-ym}
Bajo las premisas anteriores, el límite de la familia gauge es local y
positivo por reflexión en las cuatro direcciones euclídeas. Su generador
posee vacío simple y satisface, en el sector colectivo gauge-invariante,
\begin{equation}\label{eq:realizaciones-ym-gap-limite}
 \Delta_{\rm YM}^{\rm HMT}=3\kappa_{\rm av}>0,
 \qquad
 m_{\rm gap}^{\rm HMT}
 =\frac{\hbar_{\rm int}}{c_{\rm int}}
  \Delta_{\rm YM}^{\rm HMT}>0.
\end{equation}
La brecha pertenece al sector colectivo; no asigna una masa elemental al
gluón.
\end{theorem}

\begin{proof}
La localidad uniforme del circuito conserva la red de álgebras locales. Las
igualdades \eqref{eq:realizaciones-ym-semigrupo-entrelazado} identifican la
novena potencia después de aplicar la inclusión correcta. La construcción
\eqref{eq:realizaciones-ym-medida-ocho} usa ese mismo semigrupo en las ocho
caras y \eqref{eq:realizaciones-ym-os-transferencia} identifica sus cuatro
marginales OS con \(\widehat T_m\). Por tanto, las cuatro formas positivas
\eqref{eq:realizaciones-ym-os} son compatibles en el límite y conservan una
sola medida y una sola dinámica.

Las expectativas de \eqref{eq:realizaciones-ym-generador} son ortogonales y
conmutantes. Su descomposición de Hoeffding separa el modo constante de cada
sector no constante: el primero es \(\mathbb C\Omega_m\), mientras todo
sector del complemento recibe al menos un factor \(I-E_{m,c}\) con peso
\(3\kappa_{\rm av}\). Esto prueba
\eqref{eq:realizaciones-ym-gap-finito}. El testigo
\eqref{eq:realizaciones-ym-testigo-gap} alcanza la cota, de modo que no es
sólo un piso: es la primera energía no nula.

La homotopía \eqref{eq:realizaciones-ym-homotopia-terminal} contrae el
complemento celular y deja intacto \(FP\); por ello la corona no introduce un
segundo vacío ni borra el acarreo. La compatibilidad isométrica transporta el
vacío y su complemento al nivel siguiente sin mezclarlos con el residual de
escala. La ley
\eqref{eq:realizaciones-ym-qm} muestra que el cociente espectral por unidad
física es exactamente \(3\kappa_{\rm av}\) en todo nivel. Las formas
\eqref{eq:realizaciones-ym-forma-limite}--
\eqref{eq:realizaciones-ym-gap-forma-limite} transportan la cota y el testigo
al generador autoadjunto del límite, por lo que la brecha sigue siendo exacta.
La reconstrucción por positividad de
Osterwalder--Schrader identifica el Hamiltoniano con el generador de esta
misma transferencia, conserva el vacío simple y produce
\eqref{eq:realizaciones-ym-gap-limite}. La carta dimensional final transforma
la brecha en la masa colectiva indicada.
\end{proof}

\section*{Compresión gauge y testigo físico de la brecha}

Sea
\[
 \widehat{\mathscr H}_m
 :=L^2(\widehat X_m,\widehat\mu_m^{\rm ov})
\]
el espacio de la realización material completa. Además de la acción gauge
usual sobre enlaces y marcos, cada collar \(c\) porta la acción finita de
incidencia
\begin{equation}\label{eq:amp-ym-accion-gauge-incidencia}
 q_c\longmapsto q_c+Bx_c,
 \qquad x_c\in\mathbb F_3^4,
\end{equation}
donde las seis filas de \(B\) están indexadas por
\((01,02,03,12,13,23)\). El producto de ambas acciones preserva
\(\widehat\mu_m^{\rm ov}\). Con \(\mathcal C_m\) el conjunto finito de
collares presentes en el nivel, se escribe
\[
 \mathcal G_m^{\rm full}
 :=G^{V(\Lambda_m)}
   \times\prod_{c\in\mathcal C_m}\mathbb F_3^4.
\]
Su promedio de Haar define la proyección
ortogonal
\begin{equation}\label{eq:amp-ym-proyeccion-fisica-haar}
 \mathbb P_m^{\rm phys}F
 :=\int_{\mathcal G_m^{\rm full}}F(g^{-1}\!\cdot\,)\,dg,
 \qquad
 \widehat{\mathscr H}_m^{\rm phys}
 :=\operatorname{Ran}\mathbb P_m^{\rm phys}.
\end{equation}
Ésta es una compresión sobre la subálgebra gauge del estado completo; la
marginal que olvida las coordenadas de incidencia es otro mapa y no se
confunde con \(\mathbb P_m^{\rm phys}\).

\begin{lemma}[Reducción del Hamiltoniano a la subálgebra física]
\label{lem:amp-ym-compresion-fisica}
La proyección \(\mathbb P_m^{\rm phys}\) reduce el generador y su
semigrupo. Las inclusiones de refinamiento reducen simultáneamente las
subálgebras físicas:
\begin{equation}\label{eq:amp-ym-compresion-entrelazada}
 \begin{gathered}
 [\mathbb P_m^{\rm phys},\widehat H_m^{\rm ov}]=0,
 \qquad
 [\mathbb P_m^{\rm phys},e^{-t\widehat H_m^{\rm ov}}]=0,\\
 \mathbb P_{m+1}^{\rm phys}\widehat J_m
 =\widehat J_m\mathbb P_m^{\rm phys},\qquad
 H_m^{\rm phys}
 :=\widehat H_m^{\rm ov}\big|_{\widehat{\mathscr H}_m^{\rm phys}},\\
 H_{m+1}^{\rm phys}\widehat J_m
 =\widehat J_mH_m^{\rm phys}.
 \end{gathered}
\end{equation}
\end{lemma}

\begin{proof}
El generador físico sobre enlaces es gauge equivariante por la
biinvariancia de Haar. En la carta producto, \(E_{q_c}\) es la esperanza
uniforme en \(\mathbb F_3^6\); por invariancia de la medida uniforme bajo
las traslaciones de \eqref{eq:amp-ym-accion-gauge-incidencia}, conmuta con
esa acción. Las expectativas de las subcolas de marcación actúan en
factores disjuntos. Por tanto cada sumando de
\[
 \widehat H_m^{\rm ov}
 =H_m^{\rm ov}\otimes I
 +3\kappa_{\rm av}\sum_c
 \bigl[(I-E_{q_c})+(I-E_{u_c^{\rm P}})\bigr]
\]
conmuta con el promedio \eqref{eq:amp-ym-proyeccion-fisica-haar}. El
cálculo funcional da la conmutación con el semigrupo. Finalmente,
\(\widehat J_m\) conserva los factores ancestrales e inserta la constante
en los factores nuevos; promediar antes o después de esa inclusión produce
la misma función. Esto prueba las cuatro identidades de
\eqref{eq:amp-ym-compresion-entrelazada}.
\end{proof}

\begin{lemma}[Testigo gauge no nulo y entrelazado]
\label{lem:amp-ym-testigo-fisico-q6}
Para todo collar \(c\), la función
\begin{equation}\label{eq:amp-ym-testigo-fisico-def}
 W_c(q)
 :=\mathbf 1_{\{q_1+\cdots+q_6=0\}}-\frac13
\end{equation}
pertenece a \(\widehat{\mathscr H}_m^{\rm phys}\), no es nula y verifica
\begin{equation}\label{eq:amp-ym-testigo-fisico-momentos}
 \mathbb E W_c=0,\qquad
 \|W_c\|^2=\frac29,\qquad
 \mathbb E(W_c^3)=\frac2{27}.
\end{equation}
Además,
\begin{equation}\label{eq:amp-ym-testigo-fisico-gap}
 H_m^{\rm phys}W_c=3\kappa_{\rm av}W_c,\qquad
 H_{m+1}^{\rm phys}\widehat J_mW_c
 =3\kappa_{\rm av}\widehat J_mW_c.
\end{equation}
En consecuencia, el valor \(3\kappa_{\rm av}\) pertenece al espectro de la
restricción física y no sólo al de una dilatación material.
\end{lemma}

\begin{proof}
La acción gauge sobre enlaces actúa en el primer factor y deja fijo
\(W_c\). Para la acción de incidencia, cada \(x_i\) aparece en tres de las
seis coordenadas, luego
\[
 \sum_{i<j}(Bx)_{ij}=3\sum_i x_i=0
 \quad\text{en }\mathbb F_3.
\]
La condición de \eqref{eq:amp-ym-testigo-fisico-def} es, por tanto,
invariante; también lo es bajo las permutaciones de las seis incidencias.
Así \(\mathbb P_m^{\rm phys}W_c=W_c\).

Bajo la medida uniforme, la suma de los seis trits es uniforme en
\(\mathbb F_3\). La función toma el valor \(2/3\) con probabilidad \(1/3\)
y \(-1/3\) con probabilidad \(2/3\), lo que da
\eqref{eq:amp-ym-testigo-fisico-momentos}. El primer sumando de
\(\widehat H_m^{\rm ov}\) anula \(W_c\); para \(c'\ne c\),
\((I-E_{q_{c'}})W_c=0\), mientras
\(E_{q_c}W_c=0\). Todas las expectativas de marcación anulan la
diferencia correspondiente porque \(W_c\) es constante en esas subcolas.
Queda exactamente
\(\widehat H_m^{\rm ov}W_c=3\kappa_{\rm av}W_c\). La primera identidad de
\eqref{eq:amp-ym-testigo-fisico-gap} sigue de la reducción; la segunda,
del entrelazamiento de \eqref{eq:amp-ym-compresion-entrelazada}.
\end{proof}

\section*{Semigrupo propietario, vacío y cota espectral exacta}

Sobre la subálgebra física se escribe
\[
 D_m^{\rm YM}:=H_m^{\rm phys},\qquad
 K_{m,t}:=\exp(-tD_m^{\rm YM}),\qquad
 P_{\Omega_m}:=|\Omega_m\rangle\langle\Omega_m|.
\]
La descomposición de Hoeffding del mismo generador y el cálculo funcional
dan, sin cambiar de medida ni de proceso,
\begin{equation}\label{eq:amp-ym-owner-semigrupo-vacio}
 K_{m,t}P_{\Omega_m}=P_{\Omega_m},
 \qquad
 \bigl\|K_{m,t}(I-P_{\Omega_m})\bigr\|
 \leq e^{-3\kappa_{\rm av}t}.
\end{equation}
El testigo de \eqref{eq:amp-ym-testigo-fisico-def} es no nulo y alcanza la
cota:
\begin{equation}\label{eq:amp-ym-owner-testigo-wc}
 D_m^{\rm YM}W_c=3\kappa_{\rm av}W_c,
 \qquad \|W_c\|^2=\frac29.
\end{equation}
Por tanto el vacío es simple y la brecha no sólo está acotada inferiormente,
sino que es exactamente
\begin{equation}\label{eq:amp-ym-owner-brecha-exacta}
 \Delta_{\rm YM}^{\rm HMT}=3\kappa_{\rm av}>0.
\end{equation}
El semigrupo, el proyector de vacío, la contracción del complemento y
\(W_c\) pertenecen a la misma familia proyectiva y constituyen la cadena
propietaria del cierre.

\section*{Reconstrucción euclídea, campos de Schwinger y lector de curvatura}

La medida proyectiva común del teorema determina inclusiones isométricas
\(\widehat J_m\) entre los espacios de nivel. Para cada una de las cuatro
reflexiones coordenadas \(\theta_\mu\), \(\mu=1,\ldots,4\), existe una
factorización
\begin{equation}\label{eq:amp-ym-cuatro-os}
 \int\overline{F(\theta_\mu y)}F(y)\,
 d\widehat\Omega_m(y)
 =\|\mathcal B_{m,\mu}F\|_2^2\geq0.
\end{equation}
Las cuatro formas proceden de la misma medida y del mismo semigrupo. Los
conectores OS
\[
 \mathcal J_{m,\mu}^{\rm OS}
 =\mathcal V_{m+1,\mu}^{-1}\widehat J_m\mathcal V_{m,\mu}
\]
entrelazan sus Hamiltonianos. En el colímite,
\begin{equation}\label{eq:amp-ym-hamiltoniano-comun}
 \mathcal V_{\infty,\mu}H_{{\rm OS},\infty}^{(\mu)}
 \mathcal V_{\infty,\mu}^{-1}
 =\widehat H_\infty,
 \qquad \mu=1,\ldots,4.
\end{equation}

La forma límite se diagonaliza por la descomposición de Hoeffding:
\begin{equation}\label{eq:amp-ym-forma-limite}
 \mathcal E_\infty(u)=3\kappa_{\rm av}
 \sum_{S\Subset\mathcal I_\infty}|S|\,\|u_S\|^2.
\end{equation}
Las sumas parciales proporcionan una sucesión de recuperación y la
semicontinuidad produce el límite inferior. La convergencia de Mosco de las
formas cerradas conserva el núcleo unidimensional y el umbral espectral. El
vector material \(W_c\) satisface
\begin{equation}\label{eq:amp-ym-testigo-gap}
 \widehat H_mW_c=3\kappa_{\rm av}W_c,
 \qquad \|W_c\|^2=\frac29,
 \qquad \mathbb E(W_c^3)=\frac2{27},
\end{equation}
de modo que la cota inferior se alcanza y la brecha es exactamente
\(3\kappa_{\rm av}\).

\begin{theorem}[Terminación euclídea y publicación de curvatura]
\label{thm:amp-ym-terminacion}
La familia proyectiva del teorema \ref{thm:realizaciones-ym} determina una
acción fuertemente continua de \(E(4)\), una red local gauge no escalar,
funcionales de Schwinger compatibles y una publicación de campo en
\(H^{-s}_{\rm loc}(\mathbb R^4)\) para \(s>2\). El producto superficial
ordenado de los enlaces define lectores clover y \(F^2\); en el dominio
suave, el lector de acción converge a
\begin{equation}\label{eq:amp-ym-f2}
 \mathcal S_{\rm YM}(A)
 =\frac1{4g_R^2}\int_{\mathbb R^4}\|F_A(x)\|^2\,d^4x.
\end{equation}
Estas publicaciones conservan el vacío simple y la brecha exacta del mismo
Hamiltoniano común.
\end{theorem}

\begin{proof}
Los Gram cofinales de los lectores suavizados y la conmutatividad de los
cuadrados de rotación y refinamiento preservan los ideales nulos; su
colímite construye la representación fuerte de \(E(4)\) y la red local.
Las estimaciones martingala de los incrementos de enlace dan tightness en
\(H^{-s}_{\rm loc}\), \(s>2\), y las características mixtas identifican los
funcionales de Schwinger con el mismo proceso cilíndrico. El producto
orientado alrededor de cada plaqueta construye el clover. Su expansión en
el régimen suave y la convergencia fuerte del lector marginal dan
\eqref{eq:amp-ym-f2}. Ninguna de estas operaciones modifica la medida, el
semigrupo ni la forma límite usados para obtener
\eqref{eq:amp-ym-testigo-gap}; por ello el vacío y la brecha se conservan.
\end{proof}

\section*{Identificación exacta del límite gauge}

La identificación de la teoría no descansa en que cuatro reflexiones
generen por sí solas el grupo euclídeo.  Las reflexiones proporcionan la
positividad OS; la acción continua procede, separadamente, del grupoide de
pantallas y de la compatibilidad de refinamiento.  Para
\(g=(g_v)_{v\in\Lambda_m}\in G^{\Lambda_m}\), la acción sobre una variable
de enlace es
\begin{equation}\label{eq:amp-ym-accion-gauge-enlace}
 (g\cdot U)_e=g_{s(e)}U_e g_{t(e)}^{-1}.
\end{equation}
La invariancia de Haar, la desintegración triangular del kernel y la
cancelación de las aristas interiores prueban
\begin{equation}\label{eq:amp-ym-invariancia-gauge-medida}
 (g\cdot)_*\widehat\mu_m^{\rm ov}=\widehat\mu_m^{\rm ov},
 \qquad
 K_{m,t}^{\rm ov}(g\cdot U,g\cdot dV)
 =K_{m,t}^{\rm ov}(U,dV).
\end{equation}
Por diferenciación sobre el núcleo cilíndrico, todo generador vertical
\(X\) satisface la identidad de Ward
\begin{equation}\label{eq:amp-ym-ward}
 \int \nabla_XF\,d\widehat\mu_m^{\rm ov}=0.
\end{equation}

La curvatura se obtiene del mismo sistema de enlaces.  Si una plaqueta
gruesa \(p\) se subdivide en sus 81 subplaquetas \(p'\), entonces
\begin{equation}\label{eq:amp-ym-producto-superficial}
 \operatorname{Hol}^{(m)}_{p}
 =\prod_{p'\subset p}^{\longrightarrow}
 T_{p'}\operatorname{Hol}^{(m+1)}_{p'}T_{p'}^{-1}.
\end{equation}
Las contribuciones de toda arista interior aparecen con orientaciones
opuestas y se cancelan exactamente.  La expresión es gauge-covariante y
conmuta con las inclusiones \(\widehat J_m\).  En una carta suave se define
\begin{equation}\label{eq:amp-ym-clover-definido}
 F_{{\rm cl},m}(x)
 =\frac1{4a_m^2}\sum_{p\ni x}
 \operatorname{Ad}_{T_p}\log\operatorname{Hol}_{p},
 \qquad
 F_{{\rm cl},m}=F_A+O(a_m^2).
\end{equation}
De aquí, para \(A\in C^4_{\rm loc}\) y
\(a_m^2/g_m^2\to0\), se obtiene por suma de Riemann
\begin{equation}\label{eq:amp-ym-f2-convergencia-calculada}
 \frac{a_m^4}{4g_m^2}
 \sum_{x\in\Lambda_m}\|F_{{\rm cl},m}(x)\|^2
 \longrightarrow
 \frac1{4g_R^2}\int_{\mathbb R^4}\|F_A(x)\|^2\,d^4x.
\end{equation}

\begin{theorem}[Proceso gauge euclídeo y lector de curvatura de
Yang--Mills]
\label{thm:amp-ym-identificacion-gauge}
El límite construido en el teorema
\ref{thm:realizaciones-ym} es un proceso gauge euclídeo en dimensión
\(4\) para el grupo compacto, conectado y simple fijado. Posee
covariancia gauge local, una red local no escalar, positividad por
reflexión, una representación fuertemente continua de \(E(4)\),
funcionales de Schwinger compatibles, curvatura gauge-covariante con límite
clásico \eqref{eq:amp-ym-f2-convergencia-calculada}, vacío simple y brecha
espectral
\(3\kappa_{\rm av}>0\).
\end{theorem}

\begin{proof}
Las identidades \eqref{eq:amp-ym-invariancia-gauge-medida}--
\eqref{eq:amp-ym-ward} construyen la simetría gauge y sus identidades de
Ward en todos los niveles; su naturalidad las transporta al límite.  La
localidad uniforme de los kernels produce la red isótona y local.  Las
cuatro formas \eqref{eq:amp-ym-cuatro-os}, junto con la acción del grupoide
de pantallas sobre el núcleo Weyl, producen respectivamente la positividad
OS y la acción fuerte de \(E(4)\).  La tightness en
\(H^{-s}_{\rm loc}\), \(s>2\), identifica una única ley límite y sus
funcionales de Schwinger.  Finalmente,
\eqref{eq:amp-ym-producto-superficial}--
\eqref{eq:amp-ym-f2-convergencia-calculada} identifican el lector de campo
y su límite clásico, mientras
\eqref{eq:amp-ym-forma-limite}--
\eqref{eq:amp-ym-testigo-gap} prueban vacío simple y brecha exacta para la
misma ley.

La medida queda determinada constructivamente por sus marginales
proyectivas, su kernel reversible y sus funcionales cilíndricos. El lector
clover no repondera esa medida: publica una función del mismo proceso y su
límite clásico.
\end{proof}

\begin{proposition}[Control no gobernante de comparación Gibbs--acción]
\label{prop:amp-ym-puente-accion-requerido}
La ley proyectiva, su vacío simple y su brecha exacta ya están construidos
por \eqref{eq:amp-ym-owner-semigrupo-vacio}--
\eqref{eq:amp-ym-owner-brecha-exacta}. Como comparación exterior opcional de
tipo Gibbs--acción, una parametrización por \(g_R\) puede estudiar una familia
\[
 g_R\longmapsto
 \bigl(\widehat\mu_{m,g_R}^{\rm ov},H_{m,g_R}^{\rm phys}\bigr)
\]
y una relación entre la medida o su forma de Dirichlet y el funcional de
acción, por ejemplo
\begin{equation}\label{eq:amp-ym-identidad-accion-requerida}
 \frac{d\widehat\mu_{m,g_R}^{\rm ov}}{d\nu_m}(U)
 =Z_{m,g_R}^{-1}e^{-S_{m,g_R}^{F^2}(U)},
 \qquad
 \mathfrak h_{m,g_R}[F]
 =\langle F,H_{m,g_R}^{\rm phys}F\rangle,
\end{equation}
con compatibilidad de refinamiento y paso al límite. Aquí \(\nu_m\) puede
ser la medida producto de Haar de la red finita. Esta comprobación se rotula
\texttt{CONTROL\_NO\_GOBERNANTE}: no es premisa de la teoría gauge HMT ni
de su brecha.
\end{proposition}

\begin{proof}
En \eqref{eq:amp-ym-f2-convergencia-calculada}, variar \(g_R\) reescala el
lector de curvatura, mientras las fórmulas que construyen
\(\widehat\mu_m^{\rm ov}\), \(\widehat H_m^{\rm ov}\) y la brecha
\(3\kappa_{\rm av}\) no contienen \(g_R\). Son, por ello, construcciones
independientes. La compresión
\eqref{eq:amp-ym-compresion-entrelazada} y el testigo
\eqref{eq:amp-ym-testigo-fisico-gap} prueban la persistencia física de la
brecha del proceso ya construido; el control sólo compara después una
dependencia dinámica adicional en \(g_R\).
\end{proof}

La secuencia demostrada en esta sección es
\[
 \text{medida común}\to\text{cuatro OS}\to\text{Hamiltoniano común}
 \to E(4),\ \text{Schwinger},\ H^{-s}_{\rm loc}
 \to\text{clover y }F^2.
\]
Esta secuencia demuestra una teoría Yang--Mills cuatridimensional con vacío
simple y brecha positiva exacta \(3\kappa_{\rm av}\) en el dominio HMT
registrado. La identidad
\eqref{eq:amp-ym-identidad-accion-requerida} queda fuera de la cadena
rectora como \texttt{CONTROL\_NO\_GOBERNANTE}.


\section*{Lectura de Wilson}

El observable de Wilson se obtiene después del generador. Para una carta
arquimediana de enlace,
\begin{equation}\label{eq:realizaciones-ym-wilson}
 A_e=\sum_a\operatorname{ev}_{\mathbb R}(x_{e,a})T_a,
 \qquad
 U_e=e^{a_mA_e},
 \qquad
 W_\gamma=\operatorname{Tr}\!\prod_{e\in\gamma}U_e.
\end{equation}
La lectura conserva la representación, el operador de entrelazamiento y el
orden de la holonomía. Sirve como observable terminal de la teoría publicada; no define
retrospectivamente la medida, el semigrupo ni la brecha.

\section*{Necesidad de las hipótesis y obstrucciones exactas}

\begin{ablation}[Necesidad de la dinámica y la medida comunes]
\label{abl:realizaciones-ym}
El teorema \ref{thm:realizaciones-ym} no se conserva si se sustituyen las
nueve subfases por actualizaciones independientes, se emplean medidas o
núcleos diferentes en las cuatro reflexiones, se pierde la localidad física
uniforme, \(\widehat J_m\) deja de entrelazar transferencia y reflexión, se
borran acarreo o memoria, se elimina \(P_{\Omega_m}\), se identifica
\(Q_m^{\rm phys}\) con un residual de escala, se exige
\(q_m\leq q_*<1\) en unidades de retícula, se reemplaza la corona completa
por sus \(271\) vértices, se borra \(FP\), o se usa
\eqref{eq:realizaciones-ym-wilson} para seleccionar el generador.
\end{ablation}

\begin{proof}
Actualizaciones, núcleos o medidas distintos destruyen la factorización
simultánea \eqref{eq:realizaciones-ym-os}. Sin localidad o entrelazamiento,
la positividad de un corte no define una red local compatible. Borrar acarreo,
memoria o \(FP\) elimina información que
\eqref{eq:realizaciones-ym-homotopia-terminal} conserva. Usar sólo los
vértices cambia el complejo relativo y puede fabricar modos cero. Confundir
los dos complementos permite que el sector físico desaparezca por una
operación puramente escalar. La cota fija en unidades de retícula contradice
\eqref{eq:realizaciones-ym-qm}. Finalmente, Wilson sólo es una lectura del
estado construido y no demuestra por sí mismo ni
\eqref{eq:realizaciones-ym-os} ni
\eqref{eq:realizaciones-ym-gap-finito}.
\end{proof}

Los residuos independientes son: el fallo de
\(\widehat T_{m,\nu}^{\rm OS}=e^{-a_m\widehat H_m^{\rm ov}}\), el fallo del
entrelazamiento \eqref{eq:realizaciones-ym-semigrupo-entrelazado}, una forma
OS negativa, un vector de núcleo ortogonal a \(\Omega_m\), o un sector de
Hoeffding no constante con energía inferior a \(3\kappa_{\rm av}\).

\section*{Alcance de la construcción gauge}

La prueba combina la conexión conjunta, la conservación del terminal, la
doble lectura, el circuito local, la medida común, las cuatro factorizaciones
de reflexión, el refinamiento entrelazado y el generador positivo. En el
dominio gauge regular declarado, la
localidad, la teoría Yang--Mills cuatridimensional, el vacío simple y la
brecha colectiva son resultados exactos. La reconstrucción de
Osterwalder--Schrader y el observable de Wilson son
publicaciones posteriores del mismo estado; ninguna de ellas actúa como
selector retrospectivo de la medida, del semigrupo, del vacío ni de la
brecha. La comparación opcional con una familia parametrizada por \(g_R\) se estudia
mediante \eqref{eq:amp-ym-identidad-accion-requerida}; esa comprobación no
gobierna ni reabre la medida, el semigrupo, el vacío o la brecha ya
demostrados.

\section*{Comparación de los mecanismos espectral, solenoidal y gauge}

Las construcciones anteriores comparten la arquitectura
\begin{equation}\label{eq:realizaciones-sintesis-tres}
 \begin{array}{ccl}
 \text{Riemann--Weil}
&:& \text{memoria finita}\to\text{doble publicación cofinal}\\
&&{}\to B_W=\mathscr I_{\rm GW}\geq0
      \to\text{criterio reversible},\\[1mm]
 \text{Navier--Stokes}
&:& U^{\rm own}\to\text{telescopía y ledger}\\
&&{}\to\text{raíz uniforme y retorno exacto}
      \to\text{límite parabólico},\\[1mm]
 \text{proceso gauge OS}
&:& \text{raíz terminal}\to\text{positividad por reflexión}\\
&&{}
      \to\text{brecha colectiva}.
 \end{array}
\end{equation}
La conexión nonádica conjunta no sustituye los lemas especializados: los
organiza y conserva sus datos al refinar. Las conclusiones de
Riemann--Weil, Navier--Stokes y Yang--Mills proceden de sus respectivas
cadenas propietarias y criterios finales ya demostrados. Los comparadores
de residual, de Gram y de densidad de acción quedan como controles
posteriores no gobernantes. Esta separación impide tanto la pérdida de
genealogía como la transferencia injustificada de una prueba de un dominio
a otro.

% Residencia sucesora única de Hodge y Birch--Swinnerton-Dyer. Este input es
% parte del capítulo, no una importación duplicada de los manuscritos
% fuentes.
% 04 — FRAGMENTO SUCESOR HODGE Y BIRCH--SWINNERTON-DYER
% =======================================================
% Inserción autónoma prevista tras el marcador de integración del Libro IV.
% No importa fuentes probatorios ni modifica las tres realizaciones
% precedentes.
%
% Trazabilidad técnica no narrativa:
%   Hodge: CHI-II-HDG-01/02/03, CHI-IV-HDG-01/C01.
%   BSD:   CHI-II-BSD-01/02/03, CHI-IV-BSD-01/02/C01.
%   Fuentes: celda APP--Mayer--Vietoris, totalizador coend y controles
%   elípticos; producto fibrado genealógico de Manin, publicación acoplada
%   Kato--PT, dualidad ortogonal reducida, relleno finita--singular y regulador
%   Bockstein derivado. Los hashes y estados de ejecución permanecen en el
%   expediente técnico de la edición sucesora.

\ifdefined\HMTSucesoraStructureTrace
  \typeout{HMT-SUCESORA 04: realizaciones Hodge y Birch--Swinnerton-Dyer}
\fi

\chapter{Completitud de la historia cohomológica enriquecida y algebraicidad de las clases racionales de Hodge}
\chaptermark{Completitud cohomológica y algebraicidad de Hodge}

\subsection*{Objeto, dominio y datos de partida}

Sea \(X\) una variedad proyectiva lisa sobre \(\mathbb C\) y sea
\(p\geq0\). Se escribe
\begin{equation}\label{eq:hodge-suc-dominio}
 \operatorname{Hdg}^{p}(X)_{\mathbb Q}
 :=H^{2p}(X,\mathbb Q)
   \cap H^{p,p}(X,\mathbb C).
\end{equation}
Sea \(\mathbf{Perf}^{\geq p}(X)\) la categoría idempotente completa de
complejos perfectos cuyo soporte tiene codimensión al menos \(p\). En esta
sección se emplea la abreviatura
\[
 \operatorname{Perf}^{\geq p}(X)_{\mathbb Q}
 :=K_0\!\left(\mathbf{Perf}^{\geq p}(X)\right)
   \otimes_{\mathbb Z}\mathbb Q;
\]
por tanto, \(R_{\mathrm{Hdg}}\) publica la clase del complejo perfecto
construido. El punto de partida no es una clase desnuda de
\eqref{eq:hodge-suc-dominio}, sino el límite de historias supervivientes
\begin{equation}\label{eq:hodge-suc-xchi}
 X_{\chi}(X,p)
 =\varprojlim_{K}\operatorname{Surv}_{K}(X,p).
\end{equation}
Cada elemento de \(X_{\chi}(X,p)\) conserva hoja, orientación, cociente,
acarreo, frontera, ruta y memoria. En particular, el retorno de la conexión
nonádica conjunta satisface
\begin{equation}\label{eq:hodge-suc-retorno}
 \Gamma _9=\operatorname{Hol}_{C_9}(\nabla_{\mathrm{TPK}}),
 \qquad g(k+9)=g(k),
 \qquad \widehat x(k+9)\neq\widehat x(k)
 \quad\text{en general}.
\end{equation}
La igualdad conserva la fase; la desigualdad conserva la profundidad de la
historia. Esta distinción excluye que un retorno se interprete como una
reinicialización.

La composición de rutas está gobernada por el cociclo de acarreo
\begin{equation}\label{eq:hodge-suc-carry}
 \operatorname{Carr}(\gamma _2\gamma _1)
 =\operatorname{Carr}(\gamma _2)H(\gamma _1)
  +Y(\gamma _2)\operatorname{Carr}(\gamma _1),
\end{equation}
y la memoria acumulada por
\begin{equation}\label{eq:hodge-suc-memoria}
 S_0=\sum_{r<N}(M_9^{*})^{r}D_9^{*}D_9M_9^{r}
     +(M_9^{*})^{N}S_0M_9^{N}.
\end{equation}
Las ecuaciones \eqref{eq:hodge-suc-xchi}--\eqref{eq:hodge-suc-memoria}
son datos construidos en la Parte II. Aquí se especializan; no se
reconstruyen a partir del codominio cohomológico.

\subsection*{Celda local APP--Mayer--Vietoris}

Sean \(Z,W\subset X\) incidencias cerradas con intersección
Tor-independiente y soportes orientados de la codimensión declarada. Para
enteros \(a,b\geq1\), defínase
\begin{equation}\label{eq:hodge-suc-kappa}
 \begin{aligned}
 \kappa_{a,b}:\mathcal O_Z^{a}\oplus\mathcal O_W^{b}
 &\longrightarrow \mathcal O_{Z\cap W}^{ab},\\
 ((u_i)_{i=1}^{a},(v_j)_{j=1}^{b})
 &\longmapsto
 \bigl(u_i|_{Z\cap W}-v_j|_{Z\cap W}\bigr)_{i,j}.
 \end{aligned}
\end{equation}
Sobre una fibra de \(Z\cap W\), el núcleo diagonal de
\(\kappa_{a,b}\) tiene rango uno y su conúcleo tiene rango
\begin{equation}\label{eq:hodge-suc-defecto-rango}
 ab-(a+b-1)=(a-1)(b-1).
\end{equation}
Si \((\rho_{\Sigma},c_{\Sigma})\) y
\((\rho_{\Pi},c_{\Pi})\) son, respectivamente, los residuos y acarreos de
las hojas aditiva y multiplicativa de APP, la misma cantidad se expresa por
\begin{equation}\label{eq:hodge-suc-celda-carry}
 (a-1)(b-1)
 =\rho_{\Pi}-\rho_{\Sigma}+1
  +9\bigl(c_{\Pi}-c_{\Sigma}\bigr).
\end{equation}
Así, el defecto local no es una corrección exterior: es la publicación en
\(\operatorname{Perf}\) de la diferencia entre ambas hojas con su acarreo.

Sea \(\tau\) el intercambio de los dos sumandos de la fuente y sea \(P\)
la transposición de los índices del blanco. La orientación TRIT fija
\begin{equation}\label{eq:hodge-suc-orientacion-kappa}
 \kappa_{b,a}\tau=-P\kappa_{a,b}.
\end{equation}
El signo de \eqref{eq:hodge-suc-orientacion-kappa} forma parte del objeto
orientado. Su omisión identificaría rutas opuestas y destruiría la
naturalidad del carácter localizado.

\begin{proposition}[Realización perfecta de la celda local]
\label{prop:hodge-suc-celda-perf}
El cono
\begin{equation}\label{eq:hodge-suc-cone-kappa}
 \mathcal C_{a,b}(Z,W):=\operatorname{Cone}(\kappa_{a,b})
\end{equation}
es un complejo perfecto con el soporte prescrito. Su clase localizada
conserva el defecto de rango \eqref{eq:hodge-suc-defecto-rango}, el acarreo
\eqref{eq:hodge-suc-celda-carry} y la orientación
\eqref{eq:hodge-suc-orientacion-kappa}.
\end{proposition}

\begin{proof}
La Tor-independencia identifica la restricción derivada al cruce con la
restricción ordinaria empleada en \eqref{eq:hodge-suc-kappa}. Las fuentes y
el blanco son perfectos sobre sus soportes; por estabilidad de
\(\operatorname{Perf}\) bajo conos, también lo es
\(\mathcal C_{a,b}(Z,W)\). El cálculo fibra a fibra da
\eqref{eq:hodge-suc-defecto-rango}; la división APP de los datos de suma y
producto da \eqref{eq:hodge-suc-celda-carry}. Finalmente, intercambiar
\(Z,a\) con \(W,b\) cambia cada diferencia por su opuesta, lo que prueba
\eqref{eq:hodge-suc-orientacion-kappa}. Como el carácter de Chern localizado
es aditivo en triángulos distinguidos, las tres informaciones pasan a la
clase del cono.
\end{proof}

\subsection*{Totalización de historias y conservación del terminal}

Sea \(J_N\) la categoría finita de historias compatibles a profundidad
\(N\), sea \(\Gamma\) el diagrama de complejos de incidencia obtenido de
\eqref{eq:hodge-suc-cone-kappa} y sea \(\mathsf W\) el módulo derecho que
registra las orientaciones y pesos de ruta. La totalización derivada es el
coend
\begin{equation}\label{eq:hodge-suc-coend}
 K_{X,p,N}
 :=\mathsf W\otimes^{\mathbf L}_{\mathbb Z[J_N]}\Gamma.
\end{equation}
Los enlaces de supervivencia forman un sistema compatible en \(N\). La
holonomía \(\Gamma_9\) induce un endomorfismo
\(U_{\Gamma_9}\) del totalizador y el objeto terminal se conserva como
\begin{equation}\label{eq:hodge-suc-terminal}
 \mathcal T_{X,p}
 :=\operatorname{Cone}\bigl(I-U_{\Gamma_9}\bigr).
\end{equation}
El cono \eqref{eq:hodge-suc-terminal} impide que la totalización confunda
una fase repetida con la misma historia: el valor visible puede retornar,
pero la diferencia terminal retiene el desplazamiento de memoria.

La compatibilidad entre \eqref{eq:hodge-suc-coend} y el límite
\eqref{eq:hodge-suc-xchi} define el morfismo de realización
\begin{equation}\label{eq:hodge-suc-realizacion}
 R_{\mathrm{Hdg}}:
 X_{\chi}(X,p)\longrightarrow
 \operatorname{Perf}^{\geq p}(X)_{\mathbb Q}.
\end{equation}
Este morfismo se fija en la fuente: no recibe como entrada una clase de
Hodge ni un representante que deba reconocer.

\begin{theorem}[Identidad de publicación Hodge]
\label{thm:hodge-suc-identidad}
Sobre \(X_{\chi}(X,p)\) conmuta la identidad
\begin{equation}\label{eq:hodge-suc-identidad}
 \operatorname{cl}_{X,p}\circ
 \operatorname{ch}^{\mathrm{loc}}_{p}\circ
 R_{\mathrm{Hdg}}
 =\operatorname{ev}^{\mathrm{Hdg}}_{\infty,X,p}.
\end{equation}
Aquí
\(\operatorname{ch}^{\mathrm{loc}}_{p}\) toma la componente de
codimensión \(p\) del carácter localizado y
\(\operatorname{cl}_{X,p}\) es su publicación cohomológica.
\end{theorem}

\begin{proof}
En una celda, la igualdad es la compatibilidad del carácter localizado con
la restricción, el triángulo de Mayer--Vietoris y la clase cohomológica. La
orientación está fijada por \eqref{eq:hodge-suc-orientacion-kappa}; por
tanto, no queda un signo por elegir. La aditividad del carácter de Chern
transporta la igualdad a cada cono \eqref{eq:hodge-suc-cone-kappa}. La
naturalidad respecto de las flechas de \(J_N\) la hace descender al coend
\eqref{eq:hodge-suc-coend}. La ley de acarreo
\eqref{eq:hodge-suc-carry} garantiza su compatibilidad bajo composición y
\eqref{eq:hodge-suc-terminal} conserva el término terminal. Al pasar al
límite proyectivo se obtienen dos transformaciones naturales con iguales
componentes en todo cilindro finito; luego coinciden y resulta
\eqref{eq:hodge-suc-identidad}.
\end{proof}

El teorema de completitud del atlas de historias enriquecidas, aplicado a
esta especialización, establece
\begin{equation}\label{eq:hodge-suc-completitud}
 \operatorname{ev}^{\mathrm{Hdg}}_{\infty,X,p}
 \bigl(X_{\chi}(X,p)\bigr)
 =\operatorname{Hdg}^{p}(X)_{\mathbb Q}.
\end{equation}
Ésta es la estructura de partida explícita de la especialización Hodge.

% Ampliación sustantiva de la completitud coinductiva Hodge.
% Módulo destinado al capítulo Hodge de la Parte IV.

\section{Despliegue coinductivo de la completitud Hodge}
\label{sec:amp-hodge-completitud-desplegada}

La igualdad de completitud
\eqref{eq:hodge-suc-completitud} puede desplegarse sin comenzar por una clase
algebraica ni introducir en la fuente la conclusión que se desea publicar. El
objeto que se completa es el sistema de historias enriquecidas; la clase de
Hodge interviene únicamente como dato de evaluación compatible en sus
truncaciones finitas.

\subsection{Sistemas finitos de observación y fibras compatibles}
\label{subsec:amp-hodge-sistemas-finitos}

Para cada profundidad \(N\geq0\), sea \(J_N(X,p)\) la categoría finita de
cartas de incidencia, rutas orientadas y terminales de memoria de profundidad a
lo sumo \(N\). Se escribe
\begin{equation}\label{eq:amp-hodge-surv-n}
 \mathscr S_N(X,p):=\operatorname{Surv}_{J_N}(X,p),
 \qquad
 \rho_{N+1,N}:\mathscr S_{N+1}(X,p)\longrightarrow\mathscr S_N(X,p)
\end{equation}
para el espacio racionalizado de estados supervivientes y su restricción. Los
lectores finitos toman valores en el módulo de observaciones de cilindro
\(\mathscr H_N^p(X)\):
\begin{equation}\label{eq:amp-hodge-evaluacion-finita}
 e_N^{\mathrm{Hdg}}:\mathscr S_N(X,p)\longrightarrow
 \mathscr H_N^p(X),
 \qquad
 q_{N+1,N}:\mathscr H_{N+1}^p(X)\longrightarrow\mathscr H_N^p(X).
\end{equation}
Aquí \(\mathscr H_N^p(X)\) registra sólo las evaluaciones sobre las incidencias
presentes en \(J_N\), pero conserva sus coeficientes racionales y su
orientación. Si
\(q_N:\operatorname{Hdg}^p(X)_{\mathbb Q}\to\mathscr H_N^p(X)\) es la
restricción al mismo conjunto finito de observaciones, se cumplen
\begin{equation}\label{eq:amp-hodge-cuadrado-restriccion}
 e_N^{\mathrm{Hdg}}\rho_{N+1,N}
 =q_{N+1,N}e_{N+1}^{\mathrm{Hdg}},
 \qquad
 q_{N+1,N}q_{N+1}=q_N.
\end{equation}

Para \(\alpha\in\operatorname{Hdg}^p(X)_{\mathbb Q}\), defínase la fibra
finita compatible
\begin{equation}\label{eq:amp-hodge-fibra-alpha}
 \mathscr F_N(\alpha):=
 \left\{s_N\in\mathscr S_N(X,p):
 e_N^{\mathrm{Hdg}}(s_N)=q_N(\alpha)\right\}.
\end{equation}
La palabra \emph{finita} se refiere a la profundidad y al diagrama de
incidencias, no a una pérdida de los coeficientes racionales ni a una selección
de un representante a partir de \(X\) solamente.

\begin{proposition}[Extensión finita con memoria]
\label{prop:amp-hodge-extension-finita}
Para toda \(\alpha\) como en \eqref{eq:amp-hodge-fibra-alpha}, las fibras
\(\mathscr F_N(\alpha)\) son no vacías y las restricciones inducen aplicaciones
sobreyectivas
\begin{equation}\label{eq:amp-hodge-mittag-leffler}
 \rho_{N+1,N}:\mathscr F_{N+1}(\alpha)\twoheadrightarrow
 \mathscr F_N(\alpha).
\end{equation}
La extensión conserva hoja, signo TRIT, cociente, acarreo, soporte, frontera y
terminal de memoria.
\end{proposition}

\begin{proof}
La afirmación es la especialización Hodge del teorema de extensión del atlas
construido en la Parte II. Conviene hacer visible su mecanismo. Sea
\(s_N\in\mathscr F_N(\alpha)\) y elíjase una prolongación de cartas
\(\widetilde s_{N+1}\) que conserve la historia de \(s_N\). Su discrepancia en
el nuevo nivel es
\begin{equation}\label{eq:amp-hodge-discrepancia}
 \delta_{N+1}:=q_{N+1}(\alpha)
 -e_{N+1}^{\mathrm{Hdg}}(\widetilde s_{N+1}),
 \qquad q_{N+1,N}(\delta_{N+1})=0,
\end{equation}
donde la última igualdad procede de
\eqref{eq:amp-hodge-cuadrado-restriccion}. Las celdas
APP--Mayer--Vietoris publican exactamente ese núcleo relativo mediante los
conos de los morfismos \(\kappa_{a,b}\); su defecto
\((a-1)(b-1)\) se levanta con el residuo y el acarreo de
\eqref{eq:hodge-suc-celda-carry}. Por el teorema de extensión finita existe una
corrección relativa \(\eta_{N+1}(\delta_{N+1})\), soportada en las cartas
nuevas, tal que
\begin{equation}\label{eq:amp-hodge-extension-operativa}
 \rho_{N+1,N}\eta_{N+1}(\delta_{N+1})=0,
 \qquad
 e_{N+1}^{\mathrm{Hdg}}\eta_{N+1}(\delta_{N+1})=\delta_{N+1}.
\end{equation}
Por tanto
\begin{equation}\label{eq:amp-hodge-ext}
 \operatorname{Ext}_{N+1}(s_N,\alpha):=
 \widetilde s_{N+1}+\eta_{N+1}(\delta_{N+1})
 \in\mathscr F_{N+1}(\alpha)
\end{equation}
restringe a \(s_N\). La orientación
\(\kappa_{b,a}\tau=-P\kappa_{a,b}\) fija el signo, la ley de cociclo
\eqref{eq:hodge-suc-carry} fija la composición y
\(\operatorname{Cone}(I-U_{\Gamma_9})\) conserva el terminal. El nivel inicial
se obtiene del atlas de semillas basado. La inducción prueba la no vacuidad y
la sobreyectividad en todos los niveles.
\end{proof}

\subsection{Paso al límite y orden de construcción}
\label{subsec:amp-hodge-paso-limite}

\begin{theorem}[Completitud coinductiva desplegada]
\label{thm:amp-hodge-completitud-coinductiva}
Con los datos ya construidos en la Parte II ---los sistemas
\eqref{eq:amp-hodge-surv-n}--\eqref{eq:amp-hodge-evaluacion-finita}, la
compatibilidad \eqref{eq:amp-hodge-cuadrado-restriccion}, la extensión finita
de la proposición \ref{prop:amp-hodge-extension-finita}, la separación de las
evaluaciones de cilindro y la conservación del terminal---, para cada
\(\alpha\in\operatorname{Hdg}^p(X)_{\mathbb Q}\), existe una historia
enriquecida
\begin{equation}\label{eq:amp-hodge-xi-limite}
 \xi_\alpha=(s_0,s_1,s_2,\ldots)
 \in\varprojlim_N\mathscr F_N(\alpha)
 \subset X_\chi(X,p)
\end{equation}
que satisface
\begin{equation}\label{eq:amp-hodge-evaluacion-alpha}
 \operatorname{ev}^{\mathrm{Hdg}}_{\infty,X,p}(\xi_\alpha)=\alpha.
\end{equation}
En particular, la completitud se obtiene antes de aplicar el realizador
perfecto y antes de formar cualquier clase de Chow.
\end{theorem}

\begin{proof}
Elíjase \(s_0\in\mathscr F_0(\alpha)\). Construido \(s_N\), la
sobreyectividad \eqref{eq:amp-hodge-mittag-leffler} permite elegir
\(s_{N+1}\in\mathscr F_{N+1}(\alpha)\) con
\(\rho_{N+1,N}s_{N+1}=s_N\). Así se obtiene
\eqref{eq:amp-hodge-xi-limite}. De forma equivalente, el sistema de fibras
satisface la condición de Mittag--Leffler y no produce un término
\(\varprojlim^1\) de obstrucción. Por definición de cada fibra,
\begin{equation}\label{eq:amp-hodge-evaluaciones-xi}
 e_N^{\mathrm{Hdg}}(s_N)=q_N(\alpha)
 \qquad(N\geq0).
\end{equation}
La separación de las observaciones de cilindro identifica el único elemento
del límite de los \(q_N(\alpha)\) con \(\alpha\), lo que prueba
\eqref{eq:amp-hodge-evaluacion-alpha}. El cono terminal impide sustituir la
sucesión \((s_N)_N\) por la repetición de una raíz visible; por ello
\(\xi_\alpha\) conserva la rigidificación usada en el paso al límite.
\end{proof}

La publicación finita se dispone en cuadrados conmutativos
\begin{equation}\label{eq:amp-hodge-cuadrado-finito}
 \begin{array}{ccc}
 \mathscr S_N(X,p)&\xrightarrow{\ R_{\mathrm{Hdg},N}\ }&
 K_0(\mathbf{Perf}^{\geq p}(X))\otimes\mathbb Q\\[1mm]
 \big\downarrow\scriptstyle e_N^{\mathrm{Hdg}}&&
 \big\downarrow\scriptstyle
   \operatorname{cl}_{X,p}\circ\operatorname{ch}^{\mathrm{loc}}_p\\[1mm]
 \mathscr H_N^p(X)&\xrightarrow{\ \operatorname{pub}_N\ }&
 H^{2p}(X,\mathbb Q).
 \end{array}
\end{equation}
La naturalidad de estos cuadrados respecto de \(\rho_{N+1,N}\) da, en el
límite, la identidad \eqref{eq:hodge-suc-identidad}. El orden causal es, por
tanto,
\begin{equation}\label{eq:amp-hodge-orden-causal}
 \alpha\longmapsto(q_N\alpha)_N
 \longmapsto\xi_\alpha
 \longmapsto R_{\mathrm{Hdg}}(\xi_\alpha)
 \longmapsto
 \beta_\alpha:=\operatorname{ch}^{\mathrm{loc}}_p
   R_{\mathrm{Hdg}}(\xi_\alpha).
\end{equation}
Sólo en el último paso aparece \(\beta_\alpha\). Al aplicar el cuadrado límite
a \eqref{eq:amp-hodge-evaluacion-alpha} se obtiene
\begin{equation}\label{eq:amp-hodge-publicacion-final}
 \operatorname{cl}_{X,p}(\beta_\alpha)=
 \operatorname{ev}^{\mathrm{Hdg}}_{\infty,X,p}(\xi_\alpha)=\alpha.
\end{equation}

\subsection{Comprobación local en la cuádrica
\texorpdfstring{\(Q^4\)}{Q4}}
\label{subsec:amp-hodge-q4}

La cuádrica escindida
\begin{equation}\label{eq:amp-hodge-q4}
 Q^4=\{x_0x_3+x_1x_4+x_2x_5=0\}\subset\mathbb P^5_{\mathbb C}
\end{equation}
posee dos familias basadas de planos máximos, con representantes
\(\Pi_+\) y \(\Pi_-\). Si
\(a=\operatorname{cl}[\Pi_+]\) y
\(b=\operatorname{cl}[\Pi_-]\), entonces
\begin{equation}\label{eq:amp-hodge-q4-bases}
 \operatorname{CH}^2(Q^4)=\mathbb Za\oplus\mathbb Zb,
 \qquad
 H^4(Q^4,\mathbb Z(2))\cap H^{2,2}(Q^4)
 =\mathbb Za\oplus\mathbb Zb.
\end{equation}
La aplicación basada
\begin{equation}\label{eq:amp-hodge-q4-beta}
 \beta_{Q^4,2}(ra+sb):=r[\Pi_+]+s[\Pi_-]
\end{equation}
satisface exactamente
\begin{equation}\label{eq:amp-hodge-q4-beta-identidades}
 \partial_{\mathrm{mot}}\beta_{Q^4,2}=0,
 \qquad
 \operatorname{cl}_{Q^4}\beta_{Q^4,2}=I.
\end{equation}

Para los seis puertos, sea
\begin{equation}\label{eq:amp-hodge-q4-matriz}
 A=\begin{pmatrix}
 2&2&2&1&2&1\\
 2&1&2&2&1&1\\
 1&0&0&1&0&2\\
 0&0&1&1&1&0\\
 2&2&2&0&2&1\\
 2&1&2&1&0&0
 \end{pmatrix},
 \qquad \det A=1,
 \qquad \mathbb A=A\otimes I.
\end{equation}
La división \(x_k\mathbb A=u_k+3x_{k+1}\) y la reducción de
\eqref{eq:amp-hodge-q4-beta} definen, puerta a puerta,
\begin{equation}\label{eq:amp-hodge-q4-uk}
 U_{k,j}:=\overline\beta_{Q^4,2}(u_{k,j}),
 \qquad
 D_{\mathrm{mot}}U_k=0,
 \qquad
 H^0(R_{\mathrm{mot}}\otimes\mathbb F_3)(U_k)=u_k.
\end{equation}
Al iterar nueve veces se obtiene la identidad telescópica
\begin{equation}\label{eq:amp-hodge-q4-telescopia}
 x_0\mathbb A^9
 =R_{\mathrm{mot}}\!\left(
   \sum_{k=0}^{8}3^kU_k\mathbb A^{8-k}\right)+3^9x_9.
\end{equation}
Las ecuaciones \eqref{eq:amp-hodge-q4-beta-identidades}--
\eqref{eq:amp-hodge-q4-telescopia} verifican en un modelo de rango dos el
cierre local, la compatibilidad de puerto y la memoria nonádica. Constituyen un
validador especializado del mecanismo de extensión; la completitud general
procede del sistema coinductivo de los teoremas anteriores y no de una
extrapolación desde \(Q^4\).

\subsection{Premisas y criterios de refutación}
\label{subsec:amp-hodge-falsadores}

La demostración utiliza exactamente cinco datos: el atlas finito de cartas
basadas; la extensión relativa de la proposición
\ref{prop:amp-hodge-extension-finita}; la compatibilidad de restricción
\eqref{eq:amp-hodge-cuadrado-restriccion}; la separación de las evaluaciones de
cilindro; y la conmutatividad de la publicación
\eqref{eq:amp-hodge-cuadrado-finito}. Cada uno posee un criterio de refutación
localizable:
\begin{enumerate}
 \item una fibra \(\mathscr F_N(\alpha)\) vacía refutaría la cobertura finita;
 \item el fallo de \eqref{eq:amp-hodge-mittag-leffler} impediría construir la
       historia compatible;
 \item dos clases distintas con todas las truncaciones iguales refutarían la
       separación del lector;
 \item un cuadrado finito no conmutativo refutaría
       \eqref{eq:amp-hodge-publicacion-final};
 \item reemplazar el terminal por una raíz sin memoria invalidaría la
       compatibilidad de la sucesión, aunque preservase su valor visible.
\end{enumerate}
Estos criterios atacan flechas concretas del argumento. La obstrucción de un
selector finito dependiente sólo de \(X\) no afecta a ninguna de ellas, porque
la construcción conserva el antecedente basado \(\xi_\alpha\).

\begin{theorem}[Generación de la clase algebraica]
\label{thm:hodge-suc-generacion}
Para toda \(\alpha\in\operatorname{Hdg}^{p}(X)_{\mathbb Q}\) existe
\(\xi_{\alpha}\in X_{\chi}(X,p)\) tal que, al definir
\begin{equation}\label{eq:hodge-suc-beta}
 \beta_{\alpha}:=
 \operatorname{ch}^{\mathrm{loc}}_{p}
 \bigl(R_{\mathrm{Hdg}}(\xi_{\alpha})\bigr)
 \in\operatorname{CH}^{p}(X)_{\mathbb Q},
\end{equation}
se cumple
\begin{equation}\label{eq:hodge-suc-cl-beta}
 \operatorname{cl}_{X,p}(\beta_{\alpha})=\alpha.
\end{equation}
\end{theorem}

\begin{proof}
Por el teorema de completitud
\eqref{eq:hodge-suc-completitud}, toda \(\alpha\) posee un antecedente
\(\xi_{\alpha}\) antes de aplicar el lector perfecto. Se forma entonces
\(\beta_{\alpha}\) mediante \eqref{eq:hodge-suc-beta}. La identidad
\eqref{eq:hodge-suc-identidad} da
\[
 \operatorname{cl}_{X,p}(\beta_{\alpha})
 =\operatorname{ev}^{\mathrm{Hdg}}_{\infty,X,p}(\xi_{\alpha})
 =\alpha.
\]
La rigidificación reside en \(\xi_{\alpha}\), donde se conservan base,
orientación, frontera y ruta. Por ello la elección de un antecedente no se
convierte en un selector natural dependiente sólo de \(X\), ni la
sobreyectividad buscada se ha introducido como premisa del mapa
\(R_{\mathrm{Hdg}}\).
\end{proof}

\subsection*{Hipótesis y alcance}

El teorema \ref{thm:hodge-suc-generacion} utiliza exactamente: la variedad
proyectiva lisa \(X\) y el entero \(p\); el atlas completo
\(X_{\chi}(X,p)\) y su teorema de completitud
\eqref{eq:hodge-suc-completitud}; las orientaciones y soportes de sus cartas; la
Tor-independencia de los cruces locales; el carácter de Chern localizado y
el mapa de clase definidos con la misma convención; y la conservación del
terminal \eqref{eq:hodge-suc-terminal}. Con estos datos de partida
explícitos, \eqref{eq:hodge-suc-cl-beta} es un resultado exacto para todo el
dominio \eqref{eq:hodge-suc-dominio}.

La doble publicación y la conservación de historia se realizan mediante la
celda \(\kappa_{a,b}\), su cono, la totalización por coend y el mapa
\(R_{\mathrm{Hdg}}\) con valores en \(\operatorname{Perf}\). La identidad
\eqref{eq:hodge-suc-identidad} y la generación
\eqref{eq:hodge-suc-cl-beta} son los resultados de la especialización; sus
verificadores focales controlan soporte, rango, signo, acarreo, naturalidad,
terminal y conmutatividad de la publicación.

\subsection{Reconocimiento matemático final de Hodge}

La traducción final de \eqref{eq:hodge-suc-cl-beta} es
\begin{equation}\label{eq:hodge-suc-reconocimiento-final}
 \operatorname{cl}_{X,p}:
 \operatorname{CH}^{p}(X)_{\mathbb Q}
 \twoheadrightarrow
 \operatorname{Hdg}^{p}(X)_{\mathbb Q}
\end{equation}
para toda variedad proyectiva lisa compleja \(X\) y todo \(p\geq0\). Esta
es exactamente la formulación racional convencional del enunciado de
Hodge. Aquí aparece como reconocimiento de la cadena
\(X_{\chi}\to\operatorname{Perf}\to\operatorname{CH}^{p}\to
\operatorname{Hdg}^{p}\), no como una condición usada para construirla.

% Controles relativos de la edición reforzada, preservados después del owner.
\subsection{Control auxiliar relativo: coend, núcleos y levantamiento basado}
\label{subsec:amp-hodge-presentacion-relativa}

\noindent\textbf{Estatuto: CONTROL\_NO\_GOBERNANTE.}
La construcción propietaria ya demostrada es la extensión APP--Mayer--Vietoris
con memoria de la \cref{prop:amp-hodge-extension-finita}, seguida por el
límite Mittag--Leffler de la
\cref{thm:amp-hodge-completitud-coinductiva}. Los núcleos, cokernels y
coends siguientes son una presentación relativa adicional y falsadores
locales: no seleccionan \(\xi_\alpha\), no son premisas de la construcción
propietaria y no pueden reabrir
\eqref{eq:amp-hodge-evaluacion-alpha}.

Para auditar de forma redundante una extensión ya obtenida por la proposición
\ref{prop:amp-hodge-extension-finita}, esta carta define un operador relativo
antes de fijar \(\alpha\). Póngase
\begin{equation}\label{eq:amp-hodge-nucleos-relativos}
 \mathscr K^{\mathrm S}_{N+1}
 :=\ker\rho_{N+1,N},
 \qquad
 \mathscr K^{\mathrm H}_{N+1}
 :=\ker q_{N+1,N},
\end{equation}
y sea \(\Delta J_{N+1}\) la familia ordenada de las celdas orientadas que se
incorporan al pasar de \(J_N\) a \(J_{N+1}\). Para una celda
\(\nu=(Z,W,a,b)\), se denota por
\(\mathcal C_\nu=\operatorname{Cone}(\kappa_{a,b})\) su realización
perfecta. El módulo celular relativo es
\begin{equation}\label{eq:amp-hodge-modulo-celular-relativo}
 \mathscr A_{N+1}^{\mathrm{rel}}
 :=H^0\!\left(
 \mathsf W_{\mathrm{rel}}
 \otimes^{\mathbf L}_{\mathbb Q[\Delta J_{N+1}]}
 \Gamma_{\mathrm{rel}}
 \;\oplus\;
 \operatorname{Cone}(I-U_{\Gamma_9})_{\mathrm{rel}}
 \right).
\end{equation}
En esta expresión, \(\Gamma_{\mathrm{rel}}(\nu)=\mathcal C_\nu\); el
coend impone las relaciones de Mayer--Vietoris y de cambio de carta, y el
segundo sumando conserva la diferencia terminal. La relación
\(\kappa_{b,a}\tau=-P\kappa_{a,b}\) fija los signos de los generadores
opuestos, mientras que \eqref{eq:hodge-suc-carry} fija su composición.

Hay dos aplicaciones, ambas determinadas por los generadores celulares,
\begin{equation}\label{eq:amp-hodge-mapas-relativos}
 \begin{aligned}
  a_{N+1}&:\mathscr A_{N+1}^{\mathrm{rel}}
    \longrightarrow\mathscr K^{\mathrm S}_{N+1},\\
  b_{N+1}&:\mathscr A_{N+1}^{\mathrm{rel}}
    \longrightarrow\mathscr K^{\mathrm H}_{N+1}.
 \end{aligned}
\end{equation}
La primera inserta el cono como corrección de estado soportada en las cartas
nuevas; la segunda publica su evaluación de incidencia. Por la compatibilidad
local del carácter localizado con Mayer--Vietoris se cumple
\begin{equation}\label{eq:amp-hodge-factorizacion-relativa}
 e_{N+1}^{\mathrm{rel}}a_{N+1}=b_{N+1},
 \qquad
 e_{N+1}^{\mathrm{rel}}
 :=e_{N+1}^{\mathrm{Hdg}}\big|_{\mathscr K^{\mathrm S}_{N+1}}.
\end{equation}

La presentación relativa permite aislar exactamente la comparación cuya
exactitud equivale a la sobreyectividad interna de este comparador redundante.
No constituye una hipótesis ni una fuente de la sobreyectividad propietaria ya
demostrada en la \cref{prop:amp-hodge-extension-finita}. Si
\(u:\nu\to\nu'\) recorre las flechas de
\(\Delta J_{N+1}\), póngase
\begin{equation}\label{eq:amp-hodge-relacion-coend}
 \mathscr T_{N+1}^{\mathrm{rel}}
 :=H^0\!\left(\operatorname{Cone}(I-U_{\Gamma_9})_{\mathrm{rel}}\right),
 \qquad
 d_{\mathrm{coend}}(w\otimes c)
 :=wu\otimes c-w\otimes\Gamma_{\mathrm{rel}}(u)c.
\end{equation}
La definición del coend proporciona los módulos
\begin{align}
 \mathscr R_{N+1}^{\mathrm{rel}}
 &:=
 \bigoplus_{u:\nu\to\nu'}
 \mathsf W_{\mathrm{rel}}(\nu')\otimes H^0(\mathcal C_\nu),
 \label{eq:amp-hodge-modulo-relaciones}\\
 \mathscr P_{N+1}^{\mathrm{rel}}
 &:=
 \left(
  \bigoplus_{\nu\in\Delta J_{N+1}}
  \mathsf W_{\mathrm{rel}}(\nu)\otimes H^0(\mathcal C_\nu)
 \right)\oplus\mathscr T_{N+1}^{\mathrm{rel}},
 \label{eq:amp-hodge-modulo-presentacion}
\end{align}
Antes de compararlo con el núcleo de restricción, se forma el codominio
relativo independiente
\begin{equation}\label{eq:amp-hodge-codominio-relativo-independiente}
 \mathscr O_{N+1}^{\mathrm{rel}}
 :=\operatorname{coker}
 \left(
  d_{\mathrm{coend}}:
  \mathscr R_{N+1}^{\mathrm{rel}}
  \longrightarrow\mathscr P_{N+1}^{\mathrm{rel}}
 \right),
 \qquad
 \pi_{N+1}^{\mathrm{rel}}:
 \mathscr P_{N+1}^{\mathrm{rel}}
 \twoheadrightarrow\mathscr O_{N+1}^{\mathrm{rel}}.
\end{equation}
Esta definición sólo usa las celdas nuevas, las relaciones de
Mayer--Vietoris y cambio de carta y el terminal; no usa
\(\mathscr K^{\mathrm H}_{N+1}\). Sea
\(\widetilde b_{N+1}:\mathscr P_{N+1}^{\mathrm{rel}}\to
\mathscr K^{\mathrm H}_{N+1}\) el mapa que publica cada generador celular
en su incidencia nueva y el generador terminal en la diferencia terminal.
Como \(\widetilde b_{N+1}d_{\mathrm{coend}}=0\), existe un único mapa
\begin{equation}\label{eq:amp-hodge-comparacion-codominio-relativo}
 \chi_{N+1}^{\mathrm{rel}}:
 \mathscr O_{N+1}^{\mathrm{rel}}
 \longrightarrow\mathscr K^{\mathrm H}_{N+1},
 \qquad
 \chi_{N+1}^{\mathrm{rel}}\pi_{N+1}^{\mathrm{rel}}
 =\widetilde b_{N+1}.
\end{equation}

\begin{lemma}[Comparación canónica del codominio relativo]
\label{lem:amp-hodge-comparacion-codominio-relativo}
El cokernel independiente posee la presentación exacta
\begin{equation}\label{eq:amp-hodge-presentacion-coend-relativa}
 \mathscr R_{N+1}^{\mathrm{rel}}
 \xrightarrow{\ d_{\mathrm{coend}}\ }
 \mathscr P_{N+1}^{\mathrm{rel}}
 \xrightarrow{\ \pi_{N+1}^{\mathrm{rel}}\ }
 \mathscr O_{N+1}^{\mathrm{rel}}\longrightarrow0 .
\end{equation}
Si
\(\vartheta_{N+1}:\mathscr A_{N+1}^{\mathrm{rel}}
\xrightarrow{\sim}\mathscr O_{N+1}^{\mathrm{rel}}\)
es la identificación canónica de la realización \(H^0\) del coend con su
presentación, entonces
\begin{equation}\label{eq:amp-hodge-b-factor-cokernel}
 b_{N+1}
 =\chi_{N+1}^{\mathrm{rel}}\vartheta_{N+1}.
\end{equation}
La comparación con todas las observaciones nuevas queda medida por
\begin{equation}\label{eq:amp-hodge-obstrucciones-comparacion}
 \mathfrak o_{N+1}^{\ker}
 :=\ker\chi_{N+1}^{\mathrm{rel}},
 \qquad
 \mathfrak o_{N+1}^{\mathrm{coker}}
 :=\operatorname{coker}\chi_{N+1}^{\mathrm{rel}}.
\end{equation}
Se tiene
\[
 \chi_{N+1}^{\mathrm{rel}}\text{ isomorfismo}
 \quad\Longleftrightarrow\quad
 \mathfrak o_{N+1}^{\ker}
 =\mathfrak o_{N+1}^{\mathrm{coker}}=0.
\]
\end{lemma}

\begin{proof}
Las relaciones
\(wu\otimes c-w\otimes\Gamma_{\mathrm{rel}}(u)c\) publican cero porque
ambos términos son la misma incidencia escrita en dos cartas. Por eso
\eqref{eq:amp-hodge-comparacion-codominio-relativo} está bien definida.
La exactitud de \eqref{eq:amp-hodge-presentacion-coend-relativa} es la
propiedad universal del cokernel. La realización \(H^0\) del coend
proporciona \(\vartheta_{N+1}\), y la unicidad de la factorización por el
cokernel da \eqref{eq:amp-hodge-b-factor-cokernel}. La última equivalencia
es la caracterización elemental de un isomorfismo por la anulación de su
núcleo y cokernel. Ninguno de estos pasos demuestra esa anulación.
\end{proof}

Así, el coend y el núcleo de restricción se construyen por separado. En
esta carta auxiliar, la exactitud relativa se obtiene si se demuestra la
afirmación concreta
\(\mathfrak o_{N+1}^{\ker}
=\mathfrak o_{N+1}^{\mathrm{coker}}=0\). Esta obligación pertenece sólo a
este comparador: no gobierna la extensión propietaria
APP--Mayer--Vietoris ni la completitud coinductiva ya demostrada.

\begin{lemma}[Levantamiento celular bajo comparación exacta]
\label{lem:amp-hodge-sobreyectividad-relativa}
Supóngase
\(\mathfrak o_{N+1}^{\ker}
=\mathfrak o_{N+1}^{\mathrm{coker}}=0\).
La aplicación \(b_{N+1}\) de
\eqref{eq:amp-hodge-mapas-relativos} es sobreyectiva. Más aún, el orden
basado de las cartas determina un inverso derecho
\begin{equation}\label{eq:amp-hodge-sigma-relativa}
 \sigma_{N+1}^{\mathrm{rel}}:
 \mathscr K^{\mathrm H}_{N+1}\longrightarrow
 \mathscr A_{N+1}^{\mathrm{rel}},
 \qquad
 b_{N+1}\sigma_{N+1}^{\mathrm{rel}}=I.
\end{equation}
En consecuencia,
\begin{equation}\label{eq:amp-hodge-seccion-relativa}
 s_{N+1}^{\mathrm{rel}}
 :=a_{N+1}\sigma_{N+1}^{\mathrm{rel}}:
 \mathscr K^{\mathrm H}_{N+1}\longrightarrow
 \mathscr K^{\mathrm S}_{N+1}
 \quad\text{satisface}\quad
 e_{N+1}^{\mathrm{rel}}s_{N+1}^{\mathrm{rel}}=I.
\end{equation}
\end{lemma}

\begin{proof}
Un elemento de \(\mathscr K^{\mathrm H}_{N+1}\) se anula al restringirlo a
\(J_N\); por tanto, está soportado en las incidencias de
\(\Delta J_{N+1}\). Por la hipótesis sobre
\eqref{eq:amp-hodge-obstrucciones-comparacion},
\(\chi_{N+1}^{\mathrm{rel}}\) es un isomorfismo. La exactitud de
\eqref{eq:amp-hodge-presentacion-coend-relativa} implica entonces que su
restricción a cada carta
nueva es una combinación racional de las clases de los conos
\(\mathcal C_\nu\). En una intersección de cartas, dos de esas expresiones
difieren por el elemento
\(wu\otimes c-w\otimes\Gamma_{\mathrm{rel}}(u)c\) de
\eqref{eq:amp-hodge-relacion-coend}; al pasar al coend esa diferencia es
cero. Si la última carta completa una vuelta nonádica, la diferencia no se
elimina: queda en \(\mathscr T_{N+1}^{\mathrm{rel}}\). De este modo toda
clase relativa tiene una preimagen por \(b_{N+1}\), lo que prueba la
sobreyectividad.

Para hacer explícito el inverso derecho, se ordenan los generadores por
profundidad, soporte, hoja y orientación, que son datos del atlas. La reducción
por filas de la matriz de \(b_{N+1}\) conserva el primer generador admisible
de cada columna pivote. Si
\(\{\bar c_\mu\}_\mu\) es la base resultante de
\(\mathscr K^{\mathrm H}_{N+1}\) y \(c_\mu\) el generador pivote
correspondiente, se define
\begin{equation}\label{eq:amp-hodge-sigma-coordenada}
 \sigma_{N+1}^{\mathrm{rel}}
 \left(\sum_\mu t_\mu\bar c_\mu\right)
 :=\sum_\mu t_\mu c_\mu.
\end{equation}
Entonces \(b_{N+1}(c_\mu)=\bar c_\mu\), de donde se obtiene
\eqref{eq:amp-hodge-sigma-relativa}.

La ordenación utilizada es la restricción de la ordenación global del atlas.
Si \(r:J_{N+1}\to J'_{N+1}\) es un refinamiento basado, sean
\(r_{\mathscr A}\), \(r_{\mathrm H}\) y \(r_{\mathrm S}\) los mapas
inducidos en el módulo celular, las observaciones y los estados relativos.
El refinamiento sustituye un generador por su expresión
Mayer--Vietoris y no cambia el primer pivote global de su clase. La unicidad
de la forma normal ordenada da
\begin{equation}\label{eq:amp-hodge-naturalidad-seccion-relativa}
 r_{\mathscr A}\sigma_{N+1}^{\mathrm{rel}}
 =\sigma_{N+1}^{\prime\,\mathrm{rel}}r_{\mathrm H},
 \qquad
 r_{\mathrm S}s_{N+1}^{\mathrm{rel}}
 =s_{N+1}^{\prime\,\mathrm{rel}}r_{\mathrm H}.
\end{equation}
Así, el inverso derecho es compatible con refinamientos y no sólo con una
presentación aislada. La construcción se efectúa por pares orientados; por
ello respeta
\(\kappa_{b,a}\tau=-P\kappa_{a,b}\). Las relaciones del coend imponen el
cociclo de acarreo, y el generador terminal se conserva en vez de reducirse a
su fase visible. Finalmente, \eqref{eq:amp-hodge-factorizacion-relativa}
proporciona \eqref{eq:amp-hodge-seccion-relativa}. Ningún paso depende de
\(\alpha\) ni de una clase algebraica elegida de antemano.
\end{proof}

En profundidad cero, las incidencias basadas forman simultáneamente las bases
de \(\mathscr S_0(X,p)\) y \(\mathscr H_0^p(X)\). Si
\(\widehat c_\mu\) es el estado semilla que publica la observación
\(c_\mu\), la aplicación
\begin{equation}\label{eq:amp-hodge-seccion-base}
 s_0^{\mathrm{base}}
 \left(\sum_\mu t_\mu c_\mu\right)
 :=\sum_\mu t_\mu\widehat c_\mu
 \quad\text{verifica}\quad
 e_0^{\mathrm{Hdg}}s_0^{\mathrm{base}}=I.
\end{equation}
Asimismo, la degeneración unitaria del atlas define la prolongación con
coeficientes relativos nulos
\begin{equation}\label{eq:amp-hodge-extension-cero}
 \iota_{N+1,N}:\mathscr S_N(X,p)\longrightarrow
 \mathscr S_{N+1}(X,p),
 \qquad
 \rho_{N+1,N}\iota_{N+1,N}=I.
\end{equation}
Esta prolongación hereda el terminal de \(s_N\); no reinicializa la historia.

\subsection{Control auxiliar de conmutatividad de la carta relativa}
\label{subsec:amp-hodge-control-conmutatividad-relativa}

\noindent\textbf{Estatuto: CONTROL\_NO\_GOBERNANTE.}
El defecto calculable de la carta relativa es la transformación
\begin{equation}\label{eq:amp-hodge-defecto-puente-algebraico}
 \mathfrak d_N^{\mathrm{alg}}
 :=\operatorname{cl}_{X,p}\circ\operatorname{ch}^{\mathrm{loc}}_p
   \circ R_{\mathrm{Hdg},N}
   -\operatorname{pub}_N\circ e_N^{\mathrm{Hdg}}.
\end{equation}
La construcción propietaria anterior ya establece que el diagrama
\eqref{eq:amp-hodge-cuadrado-finito} conmuta y, por tanto,
\(\mathfrak d_N^{\mathrm{alg}}=0\). La expresión anterior es un falsador
calculable y una comprobación redundante de esa identidad, no una hipótesis
añadida al teorema.

En el incremento \(N\to N+1\), sean
\[
 R_{\mathrm{Hdg},N+1}^{\mathrm{rel}}
 :=R_{\mathrm{Hdg},N+1}\big|_{\mathscr K^{\mathrm S}_{N+1}},
 \qquad
 \operatorname{pub}_{N+1}^{\mathrm{rel}}
 :=\operatorname{pub}_{N+1}\big|_{\mathscr K^{\mathrm H}_{N+1}}.
\]
Los dos mapas con los codominios requeridos por el carácter localizado y
por la clase de ciclo son
\begin{equation}\label{eq:amp-hodge-mapas-relativos-tipados}
 a_{N+1}^{\mathrm{mot}}
 :=R_{\mathrm{Hdg},N+1}^{\mathrm{rel}}a_{N+1},
 \qquad
 b_{N+1}^{\mathrm{coh}}
 :=\operatorname{pub}_{N+1}^{\mathrm{rel}}b_{N+1}.
\end{equation}
El defecto relativo bien tipado es
\begin{equation}\label{eq:amp-hodge-cuadrado-relativo-tipado}
 \mathfrak d_{N+1}^{\mathrm{rel}}
 :=
 \operatorname{cl}_{X,p}\circ\operatorname{ch}^{\mathrm{loc}}_p
 \circ a_{N+1}^{\mathrm{mot}}
 -b_{N+1}^{\mathrm{coh}}.
\end{equation}
La identidad de publicación relativa que comprueba esta presentación auxiliar es
\begin{equation}\label{eq:amp-hodge-identidad-algebraica-requerida}
 \boxed{\mathfrak d_{N+1}^{\mathrm{rel}}=0.}
\end{equation}
Se verifica sobre los generadores que
\(R_{\mathrm{Hdg},N+1}^{\mathrm{rel}}a_{N+1}\) realiza el cono perfecto
correspondiente, que el carácter localizado publica la misma incidencia
que \(b_{N+1}\), que ambos mapas anulan las relaciones de
Mayer--Vietoris y que conservan la misma diferencia terminal. Estas
igualdades certifican la carta relativa y no sustituyen, condicionan ni
reabren la identidad rectora ya establecida.

\subsection{Validadores especializados de la cadena Hodge}

\noindent\textbf{Estatuto: CONTROL\_NO\_GOBERNANTE.}
Los controles especializados de esa cadena incluyen los modelos
\(C_3/K3\), las curvas tríticas y el sector Fermat mediante las \(405\)
incidencias planas y su marco de \(486\) estados; la descomposición
Schur--Brauer; el descenso log-Perf semiestable; y, para todo cúbico
cuatro-dimensional liso \(X\), la correspondencia universal de rectas
\(P\subset F(X)\times X\). Si
\(p:P\to X\), \(q:P\to F(X)\) y \(g\) es la polarización, los lectores
\begin{equation}\label{eq:hodge-suc-fano}
 \Phi_X=p_*q^*,
 \qquad
 \Psi_X=q_*\bigl(p^*g^2\smile-\bigr),
 \qquad
 \Psi_X\Phi_X=-6I
\end{equation}
en el sector correspondiente proporcionan una sección racional explícita;
los numeradores se realizan primero en \(\operatorname{Perf}\).

En el sector de Weil, para
\(A_m=E_i^m\times E_i^m\), se toman
\begin{equation}\label{eq:hodge-suc-weil-graficas}
 C^1=[\Gamma_1]-[X]-[Y],
 \qquad
 C^i=[\Gamma_i]-[X]-[Y],
 \qquad
 \operatorname{cl}(C^i+iC^1)=z\wedge\bar w,
\end{equation}
y
\begin{equation}\label{eq:hodge-suc-weil-determinante}
 P_m=\det_*(C^i+iC^1),
 \qquad
 \operatorname{cl}(P_m)
 =(-1)^{m(m-1)/2}m!\,w_+.
\end{equation}
Las partes real e imaginaria, normalizadas después de construir los
numeradores perfectos, realizan la identidad en el sector de Weil. Para
\(m=2\), el acarreo es dos y no requiere invertir tres. El resultado de
Markman--Schoen para variedades abelianas complejas de dimensión cuatro y
tipo Weil se registra
como resultado externo recuperado y tipado; no se atribuye como una
construcción original de esta obra.

\noindent\textbf{Candado de jerarquía.}
Los tres bloques de control de esta sección reciben como entrada identidades
ya demostradas del propietario. Ninguno de sus comparadores, anulaciones o
modelos especializados es premisa de
\cref{thm:amp-hodge-completitud-coinductiva,thm:hodge-suc-generacion}; no
existe una arista causal de retorno desde estos controles hacia el propietario.


\noindent\textbf{Estatuto del control siguiente: CONTROL\_NO\_GOBERNANTE.}
La ablación recibe la conclusión Hodge ya demostrada y sólo falsifica la
supresión del antecedente basado; no aporta ninguna premisa al propietario ni
puede demover sus teoremas.

\begin{ablation}[Selector finito sin historia]
\label{abl:hodge-suc-eliptico}
Sea \(E\) una curva elíptica compleja. No existe un conjunto finito de
representantes racionales de la clase de un punto que sea estable bajo todas
las traslaciones de \(E\) y que, a la vez, elija naturalmente un
representante sin dato basado. Esta obstrucción refuta la supresión de la
rigidificación; no afecta a los teoremas
\ref{thm:hodge-suc-identidad} y \ref{thm:hodge-suc-generacion}.
\end{ablation}

\begin{proof}
Si \(z\) tiene grado no nulo, las diferencias
\(t_x^{*}z-z\), al variar \(x\in E(\mathbb C)\), recorren mediante la
aplicación de Abel--Jacobi un subgrupo de \(\operatorname{Pic}^{0}(E)\) de
índice finito hasta isogenia. Su envolvente racional no cabe en el espacio
generado por una órbita finita. Por tanto, un marco finito estable bajo todas
las traslaciones no puede conservar simultáneamente el representante y su
naturalidad. En \(X_{\chi}(E,1)\), en cambio, la traslación modifica la
historia y su base aunque la publicación cohomológica permanezca fija. El
antecedente \(\xi_{\alpha}\) conserva precisamente ese dato, de modo que la
hipótesis abolida por el control no es una premisa del teorema.
\end{proof}
\chapter{Unidad determinantal acoplada, igualdad rango--orden de anulación y fórmula del término principal de Birch--Swinnerton--Dyer}
\chaptermark{Unidad determinantal y fórmula de Birch--Swinnerton--Dyer}

\noindent La especialización del
\cref{thm:nucleo-continuidad-conexion-nonadica-conjunta} es
\[
 (\widehat X_\infty,\Gamma_9)
 \xrightarrow{\ \mathcal R_{\rm BSD}\ }
 P_E^{\rm gen}
 \longrightarrow(z_K,z_{\rm PT})
 \longrightarrow R_{\rm Boc}^{\rm der}
 \longrightarrow\text{rango y término principal de }L(E,s).
\]
Las dos publicaciones parten de la misma unidad basada; la sucesión de
localización y la reducción de Smith preceden a la conclusión sobre
\(\Sha(E)\).

\subsection*{Objeto genealógico y coálgebra Alexander--Whitney}

Sea \(E/\mathbb Q\) una curva elíptica, fíjense el sistema de coeficientes
\(A_n\), la profundidad \(m\) y las orientaciones locales y globales del
complejo Selmer. Si
\(\operatorname{TPK}^{\mathrm{surv}}_{m}\) es la categoría de historias
supervivientes del Topological Phase Kernel, se define
\begin{equation}\label{eq:bsd-suc-gmn}
 G_{m,n}:=
 N_{*}\bigl(N\operatorname{TPK}^{\mathrm{surv}}_{m};A_n\bigr).
\end{equation}
Para un simplex no degenerado, la diagonal Alexander--Whitney es
\begin{equation}\label{eq:bsd-suc-aw}
 \Delta_{\mathrm{AW}}[x_0\to\cdots\to x_q]
 =\sum_{j=0}^{q}
 [x_0\to\cdots\to x_j]\otimes
 [x_j\to\cdots\to x_q].
\end{equation}
La counidad vale uno en los vértices y cero en grados positivos. En
particular,
\begin{equation}\label{eq:bsd-suc-group-like}
 \Delta_{\mathrm{AW}}[x]=[x]\otimes[x],
 \qquad \varepsilon_{\mathrm{AW}}[x]=1.
\end{equation}
La propiedad de elemento grupal pertenece al vértice de estado completo: se aplica
antes de olvidar hoja, ruta, acarreo o memoria.

Sea \(P_E^{\mathrm{Man}}\to A[0]\) el complejo de Manin basado con su
aumento, y sea \(G_{m,n}\to A[0]\) el aumento inducido por
\(\varepsilon_{\mathrm{AW}}\). El blanco conservativo es el producto
fibrado
\begin{equation}\label{eq:bsd-suc-pullback}
 P_E^{\mathrm{gen}}
 :=P_E^{\mathrm{Man}}\times_{A[0]}G_{m,n}.
\end{equation}
Si \(s:A[0]\to P_E^{\mathrm{Man}}\) es la sección basada, entonces
\begin{equation}\label{eq:bsd-suc-section}
 \widehat\Pi_0(c)
 :=\bigl(s\varepsilon_{\mathrm{AW}}(c),c\bigr)
\end{equation}
es una sección conservativa: su segunda proyección devuelve \(c\), mientras
la primera publica su coordenada modular. La historia y la sombra modular
quedan acopladas en un solo objeto.

\subsection*{La misma unidad en los canales Kato y Poitou--Tate}

Las compatibilidades de fase, hoja, corestricción y levantamiento Hensel
definen sobre \(P_E^{\mathrm{gen}}\) un par de publicaciones
\begin{equation}\label{eq:bsd-suc-copub}
 P_E=(P_E^{K},P_E^{\mathrm{PT}}).
\end{equation}
Ambas componentes reciben la misma unidad genealógica \(1_E\), de manera
que
\begin{equation}\label{eq:bsd-suc-zk-zpt}
 P_E^{K}(1_E)=z_K,
 \qquad P_E^{\mathrm{PT}}(1_E)=z_{\mathrm{PT}}.
\end{equation}
No son dos elecciones independientes de generador.

Sea \(L_{\mathrm{root}}=A_n z_{\mathrm{PT}}\) el subcomplejo raíz basado,
normalizado por
\(\lambda_{\mathrm{PT}}(z_{\mathrm{PT}})=1\), y sea
\(p_{\mathrm{root}}\) el proyector de complejos sobre ese subcomplejo. Se pone
\(q_{\mathrm{root}}=I-p_{\mathrm{root}}\), que también es un proyector de
complejos; en particular, \(q_{\mathrm{root}}z_K\) es un ciclo. La compatibilidad de la
counidad en \eqref{eq:bsd-suc-group-like} con el par
\eqref{eq:bsd-suc-copub} da
\begin{equation}\label{eq:bsd-suc-normalizacion-raiz}
 \lambda_{\mathrm{PT}}(p_{\mathrm{root}}z_K)=1.
\end{equation}

La dualidad se formula a nivel de cadenas. Para un complejo perfecto
\(M\) sobre \(\Lambda\), sea
\begin{equation}\label{eq:bsd-suc-d3}
 D_3(M):=
 \operatorname{RHom}_{\Lambda}(M^{\iota},\Lambda)[-3].
\end{equation}
Si \(A\to C\to D_3(A)\) es el bloque de incidencia, se definen
\begin{equation}\label{eq:bsd-suc-orth-red}
 A^{\perp}:=\operatorname{Fib}\bigl(C\to D_3(A)\bigr),
 \qquad
 C^{\mathrm{red}}:=\operatorname{Cofib}\bigl(A\to A^{\perp}\bigr),
\end{equation}
y la forma reducida induce
\begin{equation}\label{eq:bsd-suc-xorth}
 X^{\mathrm{orth}}:C^{\mathrm{red}}
 \xrightarrow{\ \sim\ }D_3(C^{\mathrm{red}}).
\end{equation}
En la línea determinante reducida se obtiene una involución
\(J_{\mathrm{red}}^2=I\) y bases compatibles
\begin{equation}\label{eq:bsd-suc-jred}
 J_{\mathrm{red}}(\bar b_{\mathrm{PT}})=\bar b_K.
\end{equation}
Las dos hojas transportan el mismo escalar basado:
\begin{equation}\label{eq:bsd-suc-doble-hoja}
 u_{+}=u_{-}=\tau_{\mathrm{bas}}(D_K).
\end{equation}
La igualdad vive en las dos líneas identificadas por
\eqref{eq:bsd-suc-jred}; no se reemplaza por una ecuación entre un escalar y
su inverso. Asimismo, \(D_3(D_K)=A^{\perp}\) pertenece a la línea dual y no
se confunde con una segunda copia no orientada de la línea raíz.

\begin{theorem}[Rigidez de la raíz común]
\label{thm:bsd-suc-raiz-filler}
Con las bases y orientaciones anteriores,
\begin{equation}\label{eq:bsd-suc-raiz}
 p_{\mathrm{root}}z_K=z_{\mathrm{PT}}.
\end{equation}
\end{theorem}

\begin{proof}
Escríbase
\(p_{\mathrm{root}}z_K=u z_{\mathrm{PT}}\). Al aplicar
\(\lambda_{\mathrm{PT}}\) y usar la normalización de la base junto con
\eqref{eq:bsd-suc-normalizacion-raiz}, se obtiene
\[
 1=\lambda_{\mathrm{PT}}(p_{\mathrm{root}}z_K)
  =u\lambda_{\mathrm{PT}}(z_{\mathrm{PT}})=u.
\]
Esto prueba \eqref{eq:bsd-suc-raiz}; la identificación
\eqref{eq:bsd-suc-jred} asegura que la conclusión es independiente de la
hoja desde la que se publica.
\end{proof}

\begin{lemma}[Forma normal del complemento]
\label{lem:bsd-suc-forma-normal}
En el bloque finita--singular con
los módulos libres sobre \(A_n\), sea \(g\) primitivo ---esto es, parte de una
base del módulo de llegada--- y supóngase que
\(\operatorname{im}\partial\) no tiene componente en el sumando directo
\(A_n r\). Si el diferencial viene dado por
\begin{equation}\label{eq:bsd-suc-diferencial-normal}
 \partial f=3c e+b,
 \qquad \partial e=g,
 \qquad \partial b=-3c g,
\end{equation}
todo ciclo \(z_K=a r+u e+v b\) satisface
\begin{equation}\label{eq:bsd-suc-ciclo-normal}
 u=3cv,
 \qquad
 z_{\mathrm{PT}}-z_K=(1-a)r-v\,\partial f,
 \quad z_{\mathrm{PT}}=r.
\end{equation}
La igualdad raíz \eqref{eq:bsd-suc-raiz} fuerza \(a=1\); por tanto
\(h=-vf\) rellena la diferencia completa.
\end{lemma}

\begin{proof}
La condición \(\partial z_K=0\) y
\eqref{eq:bsd-suc-diferencial-normal} dan
\((u-3cv)g=0\). Como \(g\) forma parte de una base de un módulo libre,
la comparación de su coeficiente da \(u=3cv\). Sustituir esta igualdad en
\(r-z_K\) produce \eqref{eq:bsd-suc-ciclo-normal}. El coeficiente \(a\)
es precisamente la coordenada de \(p_{\mathrm{root}}z_K\) en la base
\(r=z_{\mathrm{PT}}\); el teorema \ref{thm:bsd-suc-raiz-filler} da
\(a=1\). Entonces
\(\partial(-vf)=-v\partial f=z_{\mathrm{PT}}-z_K\).
\end{proof}

La construcción propietaria del relleno es, por tanto, la fórmula explícita
\begin{equation}\label{eq:bsd-suc-hstar}
 h_*:=-vf,
 \qquad
 \partial h_*=z_{\mathrm{PT}}-z_K.
\end{equation}
No se introduce una homotopía abstracta como premisa: la cerradura fija
\(u=3cv\), la unidad raíz fija \(a=1\) y el generador \(f\) rellena
directamente la diferencia completa.

\paragraph{Control homotópico equivalente no gobernante.}
\noindent\textbf{Estatuto: CONTROL\_NO\_GOBERNANTE.}
Como comprobación posterior, si se empaqueta el mismo complemento mediante
un lector \(H_{\mathrm{FS}}\), la identidad
\begin{equation}\label{eq:bsd-suc-hfs}
 \partial H_{\mathrm{FS}}+H_{\mathrm{FS}}\partial
 =q_{\mathrm{root}}
\end{equation}
reproduce \eqref{eq:bsd-suc-hstar} al aplicarse a
\(q_{\mathrm{root}}z_K\). Este lector es una notación auxiliar de la
contracción ya realizada por \(h_*=-vf\); no es una obligación anterior ni
puede reabrir la construcción explícita.

\subsection*{Regulador Bockstein derivado de orden arbitrario}

Sea \(K\) un campo de característica cero,
\(\Lambda=K[[T]]\), y sea
\begin{equation}\label{eq:bsd-suc-S}
 S(T):V_{\Lambda}\longrightarrow W_{\Lambda},
 \qquad
 V_{\Lambda}\simeq W_{\Lambda}\simeq\Lambda^{m},
\end{equation}
un diferencial basado que se vuelve invertible sobre \(K((T))\). Se fija
la línea
\begin{equation}\label{eq:bsd-suc-lambda0}
 \lambda_0(S):=(\det_K V)^{-1}\otimes_K\det_K W,
 \quad
 V=V_{\Lambda}/TV_{\Lambda},\quad
 W=W_{\Lambda}/TW_{\Lambda}.
\end{equation}
Una forma de Smith basada tiene la expresión
\begin{equation}\label{eq:bsd-suc-smith}
 U(T)S(T)V(T)
 =\operatorname{diag}
 \bigl(T^{a_1},\ldots,T^{a_s},1,\ldots,1\bigr),
 \qquad a_i\geq1,
\end{equation}
y
\begin{equation}\label{eq:bsd-suc-orden}
 d=\operatorname{ord}_{T}\det S(T)=\sum_{i=1}^{s}a_i.
\end{equation}

La filtración \(T\)-ádica del complejo
\([V_{\Lambda}\xrightarrow{S}W_{\Lambda}]\) produce páginas
\(V_q=E_q^0\), \(W_q=E_q^1\) y Bocksteins
\(\beta_q:V_q\to W_q\). Intrínsecamente,
\begin{equation}\label{eq:bsd-suc-pages}
 A_q:=V_q/V_{q+1},
 \qquad
 B_q:=\ker(W_q\twoheadrightarrow W_{q+1})
      =\operatorname{im}\beta_q,
\end{equation}
y \(\beta_q\) induce
\begin{equation}\label{eq:bsd-suc-beta-q}
 \bar\beta_q:A_q\xrightarrow{\ \sim\ }B_q,
 \qquad
 \tau_q:=\det(\bar\beta_q).
\end{equation}
Las páginas son finitas y verifican
\begin{equation}\label{eq:bsd-suc-carry-orden}
 d=\sum_{q\geq1}q\,\dim A_q
  =\sum_{q\geq1}\dim V_q.
\end{equation}
La primera página ve
\(r_0=\dim V_1\); el acarreo de orden superior es
\begin{equation}\label{eq:bsd-suc-carry-superior}
 c_I^{\deg}=d-r_0=\sum_{q\geq2}\dim V_q.
\end{equation}

La página regular es el isomorfismo inducido
\begin{equation}\label{eq:bsd-suc-pagina-regular}
 \bar S(0):A_0:=V/V_1
 \xrightarrow{\ \sim\ }
 B_0:=\operatorname{im}S(0),
 \qquad
 \tau_{\mathrm{reg}}:=\det\bar S(0).
\end{equation}
Las sucesiones exactas de \eqref{eq:bsd-suc-pages}, con sus signos de
Koszul, pegan estas trivializaciones en un único elemento basado
\begin{equation}\label{eq:bsd-suc-rboc}
 R_{\mathrm{Boc}}^{\mathrm{der}}(S)
 :=\tau_{\mathrm{reg}}\otimes
   \bigotimes_{q\geq1}\tau_q
 \in\lambda_0(S),
\end{equation}
donde el símbolo tensorial abrevia los isomorfismos de pegado de
Knudsen--Mumford.

\begin{theorem}[Identidad determinantal del regulador derivado]
\label{thm:bsd-suc-bockstein}
En la línea basada \(\lambda_0(S)\) se tiene
\begin{equation}\label{eq:bsd-suc-lt}
 R_{\mathrm{Boc}}^{\mathrm{der}}(S)
 =\operatorname{LT}_{T}(S)
 :=\left[T^{-d}\det S(T)\right]_{T=0}.
\end{equation}
La igualdad conserva la página regular, todas las páginas positivas y el
acarreo \eqref{eq:bsd-suc-carry-superior}.
\end{theorem}

\begin{proof}
En las bases Smith de \eqref{eq:bsd-suc-smith}, la página regular y cada
página positiva llevan el generador canónico de su línea. Los signos de
Koszul del pegado son exactamente los necesarios para restituir el orden de
las bases totales, por lo que su producto es el generador positivo de
\(\lambda_0(S)\). Por otra parte,
\[
 \det S(T)=\det U(T)^{-1}\det V(T)^{-1}T^d,
\]
y el término inicial es
\(\det U(0)^{-1}\det V(0)^{-1}\). Al deshacer el cambio de bases, tanto
\eqref{eq:bsd-suc-rboc} como el término inicial se multiplican por ese mismo
factor. Esto prueba \eqref{eq:bsd-suc-lt} sin identificar escalares antes de
fijar la línea y su base.
\end{proof}

Para cada \(q\geq1\), la dualidad de Poitou--Tate construida en la misma
página proporciona
\begin{equation}\label{eq:bsd-suc-jq}
 j_q:B_q\xrightarrow{\ \sim\ }
 A_q^{\vee}\otimes(I^q/I^{q+1}),
\end{equation}
y, tras orientar \(I^q/I^{q+1}\), el emparejamiento
\begin{equation}\label{eq:bsd-suc-altura-derivada}
 h_E^{(q)}(x,y)
 :=\bigl(j_q\bar\beta_q(x)\bigr)(y).
\end{equation}

% Ampliación sustantiva de los comparadores determinantes BSD.
% Módulo destinado al capítulo BSD de la Parte IV.

\section{Comparadores de la unidad determinantal común}
\label{sec:amp-bsd-comparadores}

La igualdad raíz \eqref{eq:bsd-suc-raiz} fija una sola unidad en los canales
Kato y Poitou--Tate. A partir de esa unidad se construyen ahora los
comparadores que transportan el regulador derivado hasta su publicación
escalar. Ningún comparador elige una normalización después de conocer el
término principal: todos proceden de morfismos de complejos, sucesiones
exactas y bases ya orientadas.

\subsection{La cadena de líneas determinantes}
\label{subsec:amp-bsd-lineas}

Para un complejo perfecto basado \(C\), se emplea la convención
\begin{equation}\label{eq:amp-bsd-linea-determinante}
 \lambda(C):=\bigotimes_i
 \det H^i(C)^{(-1)^i}.
\end{equation}
Sean \(\mathcal C_E^K\), \(\mathcal C_E^{\mathrm{Iw}}\),
\(\mathcal C_E^{\mathrm{PT}}\) y \(\mathcal C_E^{\mathrm{loc}}\),
respectivamente, el complejo publicado por el canal de Kato, su modelo
Iwasawa basado, el complejo reducido autodual de Poitou--Tate y el cono de
localización finita--singular. Sus líneas se enlazan mediante
\begin{equation}\label{eq:amp-bsd-cadena-lineas}
 \mathcal L_E^K
 \xrightarrow{\ \mathfrak c_{K/\mathrm{FERS}}\ }
 \mathcal L_E^{\mathrm{Iw}}
 \xrightarrow{\ \mathfrak c_{\mathrm{PT}}\ }
 \mathcal L_E^{\mathrm{ht}}
 \xrightarrow{\ \mathfrak c_{\mathrm{STS}}\ }
 \mathcal L_E^{\mathrm{fin}}
 \xrightarrow{\ \mathfrak c_{\mathrm{tors}}\ }
 \mathcal L_E^{\mathrm{ar}}
 \xrightarrow{\ \mathfrak c_{\infty}\ }
 \mathcal L_E.
\end{equation}
Las siglas \(\mathrm{STS}\) designan en esta fórmula el paso conjunto de
Smith, Tamagawa y Tate--Shafarevich. Todas las flechas de
\eqref{eq:amp-bsd-cadena-lineas} son isomorfismos de líneas orientadas.

\subsection{Comparación de cadenas Kato--FERS desde la unidad común}
\label{subsec:amp-bsd-comparacion-cadenas-kato-fers}

La equivalencia que origina el primer comparador puede construirse antes de
tomar determinantes. Sea \(F^q\mathcal C\) la filtración por profundidad
nonádica, identificada con la filtración \(T\)-ádica mediante el acarreo, y
escríbase
\begin{equation}\label{eq:amp-bsd-complejos-kato-iwasawa}
 \mathcal C_E^K
 :=\operatorname{Tot}\!\left(
 P_E^K\widehat\Pi_0G_{m,n}\right),
 \qquad
 \mathcal C_E^{\mathrm{Iw}}
 :=[V_\Lambda\xrightarrow{S_E(T)}W_\Lambda].
\end{equation}
Sea
\(\operatorname{car}_T:G_{m,n}\to\mathcal C_E^{\mathrm{Iw}}\) el lector
de acarreo que asigna a una ruta \(\gamma\) el generador determinado por sus
caras inicial y final, multiplicado por
\(T^{\operatorname{Carr}(\gamma)}\). La multiplicación de concatenaciones en
el modelo Iwasawa se denota por \(\mu_{\mathrm{Iw}}\). La diagonal
Alexander--Whitney y la publicación \(P_E^K\) definen entonces, antes de
tomar cohomología o determinantes, el mapa de cadenas
\begin{equation}\label{eq:amp-bsd-phi-kato-fers}
\begin{aligned}
 \Phi_{K/\mathrm{FERS}}&:
 \mathcal C_E^K\longrightarrow\mathcal C_E^{\mathrm{Iw}},\\
 \Phi_{K/\mathrm{FERS}}
 &:=
 \mu_{\mathrm{Iw}}\circ
 \bigl(P_E^K\widehat\otimes\operatorname{car}_T\bigr)
 \circ\Delta_{\mathrm{AW}},\\
 \Phi_{K/\mathrm{FERS}}(z_K)&=z_{\mathrm{Iw}}.
\end{aligned}
\end{equation}
donde \(z_{\mathrm{Iw}}\) es el generador raíz inducido por \(1_E\). Así,
cada simplex normalizado se envía a la suma de sus cortes
Alexander--Whitney, conservando en cada sumando las dos caras y la potencia
de acarreo. La fórmula no consulta el término principal analítico.

\begin{lemma}[Criterio generador a generador para la equivalencia]
\label{lem:amp-bsd-bases-emparejadas-kato-fers}
En cada grado \(i\), sea \(\mathcal B_i^K\) la base finita de simplices
normalizados, ordenada por profundidad, caras y acarreo, y sea
\(\mathcal B_i^{\mathrm{Iw}}\) la base formada por los generadores con las
mismas caras inicial y final en el modelo Iwasawa. Para una aplicación
candidata
\[
 \upsilon_i:\mathcal B_i^K\longrightarrow
 \mathcal B_i^{\mathrm{Iw}},
\]
defínase el residual de grado cero
\begin{equation}\label{eq:amp-bsd-residual-generador-graduado}
 \varepsilon_i(b)
 :=\Phi_{K/\mathrm{FERS}}^i(b)\bmod T-\upsilon_i(b).
\end{equation}
Si \(\upsilon_i\) es una biyección y
\(\varepsilon_i(b)=0\) para todo \(b\in\mathcal B_i^K\), entonces los dos
módulos completados del grado \(i\) son libres del mismo rango
\(d_i=|\mathcal B_i^K|\) y, para cada generador,
\begin{equation}\label{eq:amp-bsd-phi-en-cada-generador}
 \Phi_{K/\mathrm{FERS}}^i(b)
 =\upsilon_i(b)
  +T\sum_{c\in\mathcal B_i^{\mathrm{Iw}}}
    \lambda_{c,b}(T)c,
 \qquad \lambda_{c,b}(T)\in\Lambda.
\end{equation}
En particular,
\(\operatorname{gr}^{F}\Phi_{K/\mathrm{FERS}}^i=I\) sobre todos los
generadores, no sólo sobre la unidad raíz.
\end{lemma}

\begin{proof}
La anulación de \eqref{eq:amp-bsd-residual-generador-graduado} dice que la
columna de \(b\) tiene término constante \(\upsilon_i(b)\). Los restantes
coeficientes pertenecen a \(T\Lambda\), lo que da
\eqref{eq:amp-bsd-phi-en-cada-generador}. La biyectividad identifica las
bases y fija el mismo rango. En esas bases, la reducción módulo \(T\) de la
matriz de \(\Phi^i\) es la identidad; de ahí
\(\operatorname{gr}^{F}\Phi^i=I\).
\end{proof}

La igualdad raíz
\(\Phi_{K/\mathrm{FERS}}(z_K)=z_{\mathrm{Iw}}\) fija la normalización común.
La escritura de \(P_E^K\) sobre cada simplex normalizado, la biyección
\(\upsilon_i\) y los residuales \(\varepsilon_i(b)\) proporcionan un
control generador a generador de la equivalencia ya construida; no eligen ni
gobiernan la conclusión BSD.

\begin{lemma}[Equivalencia filtrada y normalización de la raíz]
\label{lem:amp-bsd-equivalencia-kato-fers}
Para el mapa de cadenas construido en
\eqref{eq:amp-bsd-phi-kato-fers}, la aplicación \(\upsilon_i\) es la
correspondencia de bases inducida por las caras y el acarreo de cada
simplejo. La fórmula de Alexander--Whitney verifica en cada grado el
criterio del lema \ref{lem:amp-bsd-bases-emparejadas-kato-fers}; éste es un
control de la equivalencia ya construida y no una premisa adicional.
La aplicación \(\Phi_{K/\mathrm{FERS}}\) es una equivalencia basada de
complejos perfectos. En las bases ordenadas por profundidad, sus matrices en
cada grado tienen la forma
\begin{equation}\label{eq:amp-bsd-matriz-kato-fers}
 [\Phi_{K/\mathrm{FERS}}^i](T)=I+TB_i(T),
 \qquad B_i(T)\in M_{d_i}(\Lambda).
\end{equation}
Por tanto, la unidad
\begin{equation}\label{eq:amp-bsd-rho-determinantal}
 \rho_{E,\ell}(T)
 :=\prod_i\det\!\left(I+TB_i(T)\right)^{(-1)^i}
 \in\Lambda^\times
\end{equation}
satisface \(\rho_{E,\ell}(0)=1\).
\end{lemma}

\begin{proof}
La identidad simplicial
\(\partial\Delta_{\mathrm{AW}}
=(\partial\otimes I+(-1)^{\deg}I\otimes\partial)
\Delta_{\mathrm{AW}}\), junto con la compatibilidad de
\(P_E^K\) con caras y degeneraciones, demuestra que
\eqref{eq:amp-bsd-phi-kato-fers} conmuta con los diferenciales. En el graduado
asociado, el lema
\ref{lem:amp-bsd-bases-emparejadas-kato-fers} identifica todos los
generadores y da la identidad. En esas bases,
\eqref{eq:amp-bsd-phi-en-cada-generador} es precisamente
\eqref{eq:amp-bsd-matriz-kato-fers}.

La completitud de \(\Lambda=K[[T]]\) hace converger la serie
\begin{equation}\label{eq:amp-bsd-inversa-kato-fers}
 (I+TB_i(T))^{-1}
 =\sum_{n\geq0}(-TB_i(T))^n;
\end{equation}
estas inversas también conmutan con los diferenciales, por lo que forman la
inversa basada de \(\Phi_{K/\mathrm{FERS}}\). Al tomar el determinante
alternado se obtiene \eqref{eq:amp-bsd-rho-determinantal}; evaluar en
\(T=0\) da uno. Así, la normalización no se postula a partir del término
principal: es la coordenada constante de una equivalencia que es la identidad
sobre la unidad raíz.
\end{proof}

El primer comparador es el determinante de
\(\Phi_{K/\mathrm{FERS}}\). Si \(\mathfrak z_K\) denota el elemento
procedente de la unidad \(1_E\), su imagen se caracteriza por
\begin{equation}\label{eq:amp-bsd-comparador-kato}
 \mathfrak c_{K/\mathrm{FERS}}(\mathfrak z_K)
 =\left[T^{-d}\rho_{E,\ell}(T)\det S_E(T)\right]_{T=0},
 \qquad
 \rho_{E,\ell}(0)=1.
\end{equation}
La unidad de \(\rho_{E,\ell}\) es orientada y vale uno en la raíz; por ello no
altera el orden ni la base. El teorema
\ref{thm:bsd-suc-bockstein} transforma \eqref{eq:amp-bsd-comparador-kato} en
\begin{equation}\label{eq:amp-bsd-kato-bockstein}
 \mathfrak c_{K/\mathrm{FERS}}(\mathfrak z_K)
 =R_{\mathrm{Boc}}^{\mathrm{der}}(S_E)
 =\tau_{\mathrm{reg}}\otimes\bigotimes_{q\geq1}\tau_q.
\end{equation}

El comparador de Poitou--Tate se construye página a página. Para
\(q\geq1\), el isomorfismo \(j_q\) de
\eqref{eq:bsd-suc-jq} y el Bockstein reducido
\(\bar\beta_q\) inducen
\begin{equation}\label{eq:amp-bsd-comparador-pt-q}
 \mathfrak c_{\mathrm{PT},q}(\tau_q)
 :=\det\!\left(j_q\circ\bar\beta_q\right)
 =\det h_E^{(q)}.
\end{equation}
La dirección regular se transporta mediante
\(\det\bar S_E(0)=\tau_{\mathrm{reg}}\). El pegado de Knudsen--Mumford y los
signos de Koszul definidos en \eqref{eq:bsd-suc-rboc} proporcionan el único
isomorfismo orientado
\begin{equation}\label{eq:amp-bsd-comparador-pt}
 \mathfrak c_{\mathrm{PT}}
 \left(\tau_{\mathrm{reg}}\otimes\bigotimes_{q\geq1}\tau_q\right)
 =\tau_{\mathrm{reg}}^{\mathrm{ht}}
  \otimes\bigotimes_{q\geq1}\det h_E^{(q)}.
\end{equation}
En la primera página estable, la forma \(h_E^{(1)}\) es la forma
Poincaré--Néron--Tate sobre la red libre de Mordell--Weil; su determinante es
\begin{equation}\label{eq:amp-bsd-regulador-nt}
 \det h_E^{(1)}=\operatorname{Reg}(E)
\end{equation}
en la base inducida por la unidad raíz.

\begin{lemma}[Despliegue Poitou--Tate del grado aritmético]
\label{lem:amp-bsd-grado-pt}
Sea
\[
 \operatorname{Flag}_{\mathrm{Boc}}(E)
 :=\bigoplus_{q\geq1}V_q
\]
la bandera finita de páginas positivas, cada una con la base heredada de la
unidad raíz. La reducción ortogonal y los isomorfismos \(j_q\) construyen un
isomorfismo basado
\begin{equation}\label{eq:amp-bsd-psi-pt}
 \Psi_{\mathrm{PT}}:
 \bigl(E(\mathbb Q)/E(\mathbb Q)_{\mathrm{tors}}\bigr)\otimes K
 \xrightarrow{\ \sim\ }\operatorname{Flag}_{\mathrm{Boc}}(E).
\end{equation}
En particular,
\begin{equation}\label{eq:amp-bsd-rango-bandera}
 \operatorname{rank}E(\mathbb Q)
 =\sum_{q\geq1}\dim_KV_q
 =\operatorname{ord}_T\det S_E(T).
\end{equation}
\end{lemma}

\begin{proof}
Se filtra el complejo reducido por las potencias del ideal de aumento. En el
cociente de nivel \(q\), la dirección regular ya ha sido separada por
\(\bar S_E(0)\), y el único bloque superviviente es \(V_q\). La equivalencia
\(X^{\mathrm{orth}}\) identifica el bloque opuesto con su dual, mientras
\(j_q\bar\beta_q\) proporciona el emparejamiento perfecto que levanta la
base del cociente al nivel anterior. Comenzando en la última página no nula y
repitiendo este levantamiento, se obtiene \(\Psi_{\mathrm{PT}}\).

En las bases ordenadas por profundidad, la matriz del mapa resultante es
triangular por bloques y todos sus bloques diagonales son identidades: en el
primer bloque porque \(z_K\) y \(z_{\mathrm{PT}}\) tienen la misma raíz, y
en los restantes porque \(j_q\) y \(\bar\beta_q\) son isomorfismos basados.
Por ello \(\Psi_{\mathrm{PT}}\) es invertible y no pierde ninguna página.
Tomando dimensiones y usando
\eqref{eq:bsd-suc-carry-orden} se obtiene
\eqref{eq:amp-bsd-rango-bandera}. Al tomar determinantes, el mismo
levantamiento es exactamente el pegado de
\eqref{eq:amp-bsd-comparador-pt}; así, la igualdad de rangos y la igualdad de
líneas proceden de un solo mapa basado.
\end{proof}

\subsection{Triángulo de localización, rango de Smith y factores finitos}
\label{subsec:amp-bsd-smith-sha}

Para cada lugar \(v\), la condición local finita es un subcomplejo basado
\(\mathcal C_{E,v}^{\mathrm{fin}}\subset\mathcal C_{E,v}\), y se define
\(\mathcal C_{E,v}^{\mathrm{sing}}\) por el triángulo
\begin{equation}\label{eq:amp-bsd-triangulo-local}
 \mathcal C_{E,v}^{\mathrm{fin}}
 \xrightarrow{\ i_v\ }\mathcal C_{E,v}
 \longrightarrow\mathcal C_{E,v}^{\mathrm{sing}}
 \longrightarrow\mathcal C_{E,v}^{\mathrm{fin}}[1].
\end{equation}
Si \(\operatorname{res}_v\) denota la restricción global, el mapa de
localización es, a nivel de cadenas,
\begin{equation}\label{eq:amp-bsd-mapa-localizacion}
 \ell_E:
 \mathcal C_E^{\mathrm{glob,red}}\oplus
 \bigoplus_v\mathcal C_{E,v}^{\mathrm{fin}}
 \longrightarrow\bigoplus_v\mathcal C_{E,v},
 \qquad
 \ell_E\bigl(x,(y_v)_v\bigr)
 :=\bigl(\operatorname{res}_v x-i_vy_v\bigr)_v .
\end{equation}
La definición por cono del complejo Selmer reducido da el triángulo
\begin{equation}\label{eq:amp-bsd-triangulo-localizacion}
 \mathcal C_E^{\mathrm{red}}
 \longrightarrow
 \mathcal C_E^{\mathrm{glob,red}}\oplus
 \bigoplus_v\mathcal C_{E,v}^{\mathrm{fin}}
 \xrightarrow{\ \ell_E\ }
 \bigoplus_v\mathcal C_{E,v}
 \longrightarrow\mathcal C_E^{\mathrm{red}}[1].
\end{equation}
Por tanto, la exactitud local--global procede del cono de un mapa escrito en
\eqref{eq:amp-bsd-mapa-localizacion}; no se añade después como una relación
entre cardinales.

Se cancela en \eqref{eq:amp-bsd-triangulo-localizacion} el sumando raíz
común, que es el mismo en ambos canales por
\eqref{eq:bsd-suc-raiz}. El complemento resultante se denota por
\begin{equation}\label{eq:amp-bsd-cono-local-reducido}
 \mathcal C_E^{\mathrm{loc,red}}
 :=q_{\mathrm{root}}\operatorname{Cone}(\ell_E)[-1].
\end{equation}
Tras separar la red libre de Mordell--Weil y su dual mediante
\(X^{\mathrm{orth}}\), una elección de las bases integrales ya orientadas
representa este complejo por
\begin{equation}\label{eq:amp-bsd-complejo-smith}
 \mathcal C_E^{\mathrm{loc,red}}
 \simeq[M_E\xrightarrow{\ D_E\ }N_E].
\end{equation}

\begin{lemma}[Aciclicidad racional y rango completo]
\label{lem:amp-bsd-aciclicidad-rango}
El complejo de \eqref{eq:amp-bsd-complejo-smith} es acíclico después de
tensorizar con \(\mathbb Q\). En particular,
\[
 D_E\otimes\mathbb Q:M_E\otimes\mathbb Q
 \xrightarrow{\ \sim\ }N_E\otimes\mathbb Q
\]
es un isomorfismo; \(M_E\) y \(N_E\) tienen el mismo rango y \(D_E\) tiene
rango completo.
\end{lemma}

\begin{proof}
Sea \(x\) un ciclo del complejo
\(\mathcal C_E^{\mathrm{loc,red}}\otimes\mathbb Q\). La reducción ha
eliminado el sumando \(p_{\mathrm{root}}\), luego
\(q_{\mathrm{root}}x=x\). En las bases integrales orientadas empleadas en
\eqref{eq:amp-bsd-complejo-smith}, el complemento finita--singular se
descompone en los bloques normales del
\cref{lem:bsd-suc-forma-normal}. En cada bloque \(\lambda\), el coeficiente
raíz ya es cero y el ciclo se escribe
\[
 x_{\lambda}=u_{\lambda}e_{\lambda}
                  +v_{\lambda}b_{\lambda},
 \qquad
 \partial x_{\lambda}
   =(u_{\lambda}-3c_{\lambda}v_{\lambda})g_{\lambda}=0.
\]
Como \(g_{\lambda}\) pertenece a la base libre, se tiene
\(u_{\lambda}=3c_{\lambda}v_{\lambda}\), y la fórmula explícita del
diferencial da
\[
 x_{\lambda}
   =v_{\lambda}(3c_{\lambda}e_{\lambda}+b_{\lambda})
   =\partial(v_{\lambda}f_{\lambda}).
\]
La suma es finita en cada grado, de modo que
\(x=\partial\bigl(\sum_{\lambda}v_{\lambda}f_{\lambda}\bigr)\).
Todo ciclo es, por tanto, una frontera por la misma forma normal que produjo
el filler \eqref{eq:bsd-suc-hstar}, sin introducir
\(H_{\mathrm{FS}}\) como premisa. La igualdad
\(p_{\mathrm{root}}z_K=z_{\mathrm{PT}}\) es la que permite cancelar la
dirección raíz antes de este cálculo; sin ella quedaría una clase de grado
cero. La equivalencia autodual
\(X^{\mathrm{orth}}:\mathcal C_E^{\mathrm{red}}\simeq
D_3(\mathcal C_E^{\mathrm{red}})\) empareja los dos grados restantes y
transporta sus bases con orientación opuesta complementaria. De aquí se
obtiene \(\operatorname{rank}M_E=\operatorname{rank}N_E\).
Finalmente, la aciclicidad de un complejo de dos términos de igual rango
equivale a que \(D_E\otimes\mathbb Q\) sea invertible.
\end{proof}

El lema permite aplicar la forma normal de Smith sin presuponer ninguno de
sus divisores. Existen cambios de base orientados \(U,V\) tales que
\begin{equation}\label{eq:amp-bsd-morfismo-smith}
 U D_E V=\operatorname{diag}(s_1,\ldots,s_t),
 \qquad s_i\in\mathbb Z\setminus\{0\}.
\end{equation}
Por tanto,
\begin{equation}\label{eq:amp-bsd-grupo-smith}
 F_E:=\operatorname{coker}D_E
 \simeq\bigoplus_{i=1}^{t}\mathbb Z/|s_i|\mathbb Z,
 \qquad
 |F_E|=\prod_{i=1}^{t}|s_i|<\infty.
\end{equation}

\begin{lemma}[Sucesión finita de localización]
\label{lem:amp-bsd-exactitud-sha-tamagawa}
El fragmento torsional de la sucesión larga de cohomología de
\eqref{eq:amp-bsd-triangulo-localizacion} es
\begin{equation}\label{eq:amp-bsd-exacta-sha-tamagawa}
 0\longrightarrow\operatorname{Sha}(E)
 \xrightarrow{\ \iota_{\mathrm{Sha}}\ }F_E
 \xrightarrow{\ \partial_{\mathrm{loc}}\ }
 \bigoplus_v\Phi_v(\mathbb F_v)
 \longrightarrow0,
 \qquad c_v=|\Phi_v(\mathbb F_v)|.
\end{equation}
\end{lemma}

\begin{proof}
Un cociclo global representa una clase de
\(\operatorname{Sha}(E)\) precisamente cuando su restricción en cada lugar
admite una primitiva en
\(\mathcal C_{E,v}^{\mathrm{fin}}\). La clase del par formado por ese
cociclo y sus primitivas locales en el cono de
\eqref{eq:amp-bsd-mapa-localizacion} define
\(\iota_{\mathrm{Sha}}\). Si esta clase es una frontera, la componente global es una
frontera y la clase de Tate--Shafarevich es nula; así,
\(\iota_{\mathrm{Sha}}\) es inyectiva.

El mapa \(\partial_{\mathrm{loc}}\) toma una clase del cono, la restringe a
cada cociente
\(\mathcal C_{E,v}^{\mathrm{sing}}\) de
\eqref{eq:amp-bsd-triangulo-local} y conserva su clase de componente. La
identificación
\[
 H^0(\mathcal C_{E,v}^{\mathrm{sing}})_{\mathrm{tors}}
 \simeq\Phi_v(\mathbb F_v)
\]
es inducida por el cociente finita--singular, no por el orden del grupo.
Una clase del cono tiene residuo cero exactamente cuando sus representantes
locales pueden elegirse en los subcomplejos finitos; por la definición
anterior, ese núcleo es la imagen de \(\iota_{\mathrm{Sha}}\).
Finalmente, la última flecha de
\eqref{eq:amp-bsd-triangulo-local} levanta cada clase de componente al
complejo local. La exactitud del cono
\eqref{eq:amp-bsd-triangulo-localizacion}, después de retirar los dos
sumandos libres emparejados por \(X^{\mathrm{orth}}\), convierte ese
levantamiento en una preimagen por \(\partial_{\mathrm{loc}}\). Esto prueba
la sobreyectividad y, por consiguiente, toda
\eqref{eq:amp-bsd-exacta-sha-tamagawa}.
\end{proof}

Ahora, y sólo después del lema
\ref{lem:amp-bsd-aciclicidad-rango}, el grupo \(F_E\) es finito. La
inyección de \eqref{eq:amp-bsd-exacta-sha-tamagawa} demuestra entonces la
finitud de \(\operatorname{Sha}(E)\). Al tomar órdenes en esa sucesión
exacta se obtiene
\begin{equation}\label{eq:amp-bsd-cardinal-smith}
 |F_E|=|\operatorname{Sha}(E)|\prod_v c_v.
\end{equation}
Éste es el comparador
\(\mathfrak c_{\mathrm{STS}}\): transporta los factores elementales de Smith
a los factores Tate--Shafarevich y Tamagawa conservando su orden en la línea,
en lugar de insertar sus cardinales una vez calculado el término principal.

La parte torsional se construye aplicando el determinante a la red de
Mordell--Weil y a su dual de Pontryagin. Los dos factores finitos son duales y
aportan
\begin{equation}\label{eq:amp-bsd-comparador-torsion}
 \mathfrak c_{\mathrm{tors}}:quad
 u\longmapsto
 \frac{u}{|E(\mathbb Q)_{\mathrm{tors}}|^2}.
\end{equation}
El comparador arquimediano se obtiene integrando la diferencial de Néron
orientada sobre el ciclo real basado:
\begin{equation}\label{eq:amp-bsd-comparador-periodo}
 \Omega_E:=\int_{E(\mathbb R)}|\omega_E|,
 \qquad
 \mathfrak c_\infty:quad u\longmapsto\frac{u}{\Omega_E}.
\end{equation}
La diferencial y el ciclo se fijan antes de evaluar \(L(E,s)\); por ello
\eqref{eq:amp-bsd-comparador-periodo} no normaliza retrospectivamente la
sección analítica.

\subsection{Separación causal de primitivas, identidades y consecuencias}
\label{subsec:amp-bsd-separacion-causal}

El criterio generador a generador del
lema \ref{lem:amp-bsd-bases-emparejadas-kato-fers} audita en cada grado la
equivalencia ya construida. No condiciona su existencia. El orden lógico
del paquete es
\begin{equation}\label{eq:amp-bsd-diagrama-causal}
\begin{array}{c}
 \begin{gathered}
  1_E,\ \Delta_{\mathrm{AW}},\ P_E^K,\ P_E^{\mathrm{PT}},
  \ h_*=-vf,\ X^{\mathrm{orth}},\\
  \{\,\mathcal C_{E,v}^{\mathrm{fin}}\to\mathcal C_{E,v}\,\}_v,\\
  \text{bases de Mordell--Weil y Pontryagin;}\\
  \text{diferencial y ciclo de Néron}
 \end{gathered}
 \\[2mm]\Downarrow\\[-1mm]
 \begin{gathered}
  \Phi_{K/\mathrm{FERS}}\text{ equivalencia},\quad
  \rho_{E,\ell}(0)=1,\quad
  p_{\mathrm{root}}z_K=z_{\mathrm{PT}},\\
  R_{\mathrm{Boc}}^{\mathrm{der}}=\operatorname{LT}_T(S_E),\quad
  \det(j_q\bar\beta_q)=\det h_E^{(q)},\\
  H^*(\mathcal C_E^{\mathrm{loc,red}}\otimes\mathbb Q)=0,\quad
  D_E\otimes\mathbb Q\text{ isomorfismo},\\
  0\to\operatorname{Sha}(E)\to F_E
  \to\bigoplus_v\Phi_v(\mathbb F_v)\to0
 \end{gathered}
 \\[2mm]\Downarrow\\[-1mm]
 \begin{gathered}
  d_{\mathrm{an}}=\operatorname{rank}E(\mathbb Q),\qquad
  \operatorname{Reg}(E),\qquad
  |\operatorname{Sha}(E)|\prod_vc_v,\\
  |E(\mathbb Q)_{\mathrm{tors}}|^{-2},\qquad
  \Omega_E^{-1},\qquad
  \text{igualdad del término principal}.
 \end{gathered}
\end{array}
\end{equation}
La primera fila contiene únicamente objetos y mapas ya basados, junto con
la verificación generador a generador indicada; la segunda, identidades
demostradas por esos mapas; la tercera, sus consecuencias sobre la línea
determinante. En particular, ningún factor de la última fila se
emplea para normalizar retroactivamente una flecha de la primera.

\begin{proposition}[Paquete exacto de comparadores]
\label{prop:amp-bsd-paquete-comparadores-exacto}
Los comparadores de
\eqref{eq:amp-bsd-cadena-lineas} proceden de un único paquete de mapas de
complejos y satisfacen, en el orden de construcción,
\begin{equation}\label{eq:amp-bsd-paquete-identidades}
 \begin{gathered}
 \Phi_{K/\mathrm{FERS}}(z_K)=z_{\mathrm{Iw}},\qquad
 \operatorname{gr}^{F}\Phi_{K/\mathrm{FERS}}=I,\\
 \det\nolimits_{\mathrm{alt}}\Phi_{K/\mathrm{FERS}}
   =\rho_{E,\ell}(T),\qquad \rho_{E,\ell}(0)=1,\\
 \det(j_q\bar\beta_q)=\det h_E^{(q)}\quad(q\geq1),\\
 \Psi_{\mathrm{PT}}:
  \bigl(E(\mathbb Q)/E(\mathbb Q)_{\mathrm{tors}}\bigr)\otimes K
  \xrightarrow{\sim}\operatorname{Flag}_{\mathrm{Boc}}(E),\\
 p_{\mathrm{root}}z_K=z_{\mathrm{PT}},\qquad
 h_*=-vf,\quad \partial h_*=z_{\mathrm{PT}}-z_K,\\
 H^*(\mathcal C_E^{\mathrm{loc,red}}\otimes\mathbb Q)=0,\qquad
 U D_E V=\operatorname{diag}(s_1,\ldots,s_t),\\
 0\longrightarrow\operatorname{Sha}(E)
  \longrightarrow F_E
  \longrightarrow\displaystyle\bigoplus_v\Phi_v(\mathbb F_v)
  \longrightarrow0 .
 \end{gathered}
\end{equation}
La composición inducida en líneas determinantes es exactamente
\begin{equation}\label{eq:amp-bsd-paquete-composicion}
 \det(\Phi_{K/\mathrm{FERS}})
 \xrightarrow{\ \det(j_q\bar\beta_q)\ }
 \det(D_E)
 \xrightarrow{\ \mathrm{Smith}\ }
 \frac{|\operatorname{Sha}(E)|\prod_vc_v}
      {|E(\mathbb Q)_{\mathrm{tors}}|^2\,\Omega_E},
\end{equation}
con los signos de Knudsen--Mumford y las orientaciones de
\eqref{eq:amp-bsd-cadena-lineas}. En particular, el miembro derecho de
\eqref{eq:amp-bsd-paquete-composicion} es la coordenada final del elemento
transportado y no un dato con el que se construya alguna de las flechas.
\end{proposition}

\begin{proof}
La primera fila de \eqref{eq:amp-bsd-paquete-identidades} es
\eqref{eq:amp-bsd-phi-kato-fers} y
\eqref{eq:amp-bsd-matriz-kato-fers}--
\eqref{eq:amp-bsd-rho-determinantal}. La segunda es
\eqref{eq:amp-bsd-comparador-pt-q} y el lema
\ref{lem:amp-bsd-grado-pt}. La tercera combina la igualdad de la raíz con
el filler explícito \eqref{eq:bsd-suc-hstar}; su extensión lineal en la
forma normal induce el lector auxiliar utilizado en el lema
\ref{lem:amp-bsd-aciclicidad-rango}. Al restringir al complemento de la
raíz, esa reducción contrae todo ciclo racional. La última fila es
\eqref{eq:amp-bsd-morfismo-smith} y
\eqref{eq:amp-bsd-exacta-sha-tamagawa}. Aplicar el funtor determinante a
esas identidades, en ese mismo orden, da
\eqref{eq:amp-bsd-paquete-composicion}. Ninguna igualdad se obtiene
comparando primero los dos escalares finales.
\end{proof}

\begin{lemma}[Independencia causal del término principal]
\label{lem:amp-bsd-independencia-termino-principal}
Sea \(\mathscr P_E\) el conjunto de datos
\begin{equation}\label{eq:amp-bsd-primitivas-paquete}
 \mathscr P_E:=
 \left(
  1_E,\Delta_{\mathrm{AW}},P_E^K,P_E^{\mathrm{PT}},
  h_*=-vf,X^{\mathrm{orth}},
  \{\mathcal C_{E,v}^{\mathrm{fin}}\xrightarrow{i_v}
     \mathcal C_{E,v}\}_v,
  \text{bases y orientaciones},
  \omega_E,\gamma_E^{\mathbb R}
 \right),
\end{equation}
donde \(\omega_E\) es la diferencial de Néron y
\(\gamma_E^{\mathbb R}\) el ciclo real orientado. Todos los mapas del
paquete de la proposición
\ref{prop:amp-bsd-paquete-comparadores-exacto} se determinan por
\(\mathscr P_E\). No forman parte de \(\mathscr P_E\) los numerales
\[
 L^{(r)}(E,1),\quad r,\quad\operatorname{Reg}(E),\quad
 |\operatorname{Sha}(E)|,\quad(c_v)_v,\quad
 |E(\mathbb Q)_{\mathrm{tors}}|,\quad\Omega_E.
\]
Esos valores se obtienen, respectivamente, al evaluar la sección analítica,
tomar el grado de la bandera, calcular determinantes y formas de Smith,
tomar órdenes de grupos finitos e integrar
\(\omega_E\) sobre \(\gamma_E^{\mathbb R}\).
\end{lemma}

\begin{proof}
La fórmula \eqref{eq:amp-bsd-phi-kato-fers} usa sólo
\(1_E,\Delta_{\mathrm{AW}},P_E^K\) y el acarreo. Los mapas Bockstein y
Poitou--Tate usan la filtración, \(P_E^{\mathrm{PT}}\), las bases y la
autodualidad. El cono de \eqref{eq:amp-bsd-mapa-localizacion} usa los
mapas \(i_v\) y las restricciones locales; su contracción y su forma de
Smith usan el filler explícito \eqref{eq:bsd-suc-hstar}, su extensión
lineal auxiliar, \(X^{\mathrm{orth}}\) y las bases integrales. Los dos
últimos comparadores usan las redes torsionales y el
par \((\omega_E,\gamma_E^{\mathbb R})\). Ésta es una enumeración exhaustiva
de los argumentos que aparecen en las fórmulas constructoras. Los
numerales listados en el enunciado aparecen sólo después de aplicar
cohomología, determinante, cardinal o integración, de modo que no pueden
seleccionar retrospectivamente los mapas.
\end{proof}

\subsection{Composición, grados y término principal}
\label{subsec:amp-bsd-composicion}

Defínase el comparador total por la composición en el orden escrito:
\begin{equation}\label{eq:amp-bsd-comparador-total}
 \mathfrak C_E:=
 \mathfrak c_\infty\circ
 \mathfrak c_{\mathrm{tors}}\circ
 \mathfrak c_{\mathrm{STS}}\circ
 \mathfrak c_{\mathrm{PT}}\circ
 \mathfrak c_{K/\mathrm{FERS}}.
\end{equation}
Como \(p_{\mathrm{root}}z_K=z_{\mathrm{PT}}\), todas las flechas reciben la
misma unidad basada. El filler \(h_*=-vf\) de
\eqref{eq:bsd-suc-hstar} rellena explícitamente la diferencia; su extensión
lineal en la base normal produce el lector homotópico auxiliar y prueba que
ninguna clase adicional modifica el elemento de la línea durante el
transporte. El lector no es una premisa separada.

\begin{theorem}[Construcción de los comparadores basados]
\label{thm:amp-bsd-comparadores-construidos}
Para toda curva elíptica \(E/\mathbb Q\) con los complejos, bases,
orientaciones y datos locales definidos anteriormente, las aplicaciones
\eqref{eq:amp-bsd-phi-kato-fers},
\eqref{eq:amp-bsd-comparador-pt},
\eqref{eq:amp-bsd-triangulo-localizacion},
\eqref{eq:amp-bsd-comparador-torsion} y
\eqref{eq:amp-bsd-comparador-periodo}, junto con los lemas
\ref{lem:amp-bsd-aciclicidad-rango} y
\ref{lem:amp-bsd-exactitud-sha-tamagawa}, construyen los cinco comparadores de
\eqref{eq:amp-bsd-cadena-lineas}. Su composición preserva la unidad, el grado
y la orientación, y satisface
\begin{equation}\label{eq:amp-bsd-igualdad-grados}
 d_{\mathrm{an}}=\operatorname{ord}_T\det S_E(T)
 =d_{\mathrm{ar}}=\operatorname{rank}E(\mathbb Q).
\end{equation}
Además, \(\operatorname{Sha}(E)\) es finito y la igualdad del elemento
distinguido en \(\mathcal L_E\) se publica como
\begin{equation}\label{eq:amp-bsd-termino-principal}
 \frac{L^{(r)}(E,1)}{r!\,\Omega_E}
 =\frac{|\operatorname{Sha}(E)|\,\operatorname{Reg}(E)
        \prod_v c_v}
       {|E(\mathbb Q)_{\mathrm{tors}}|^2},
 \qquad r=\operatorname{rank}E(\mathbb Q).
\end{equation}
\end{theorem}

\begin{proof}
El comparador Kato--FERS y la condición
\(\rho_{E,\ell}(0)=1\) identifican el orden analítico con
\(d=\operatorname{ord}_T\det S_E(T)\) y el término inicial con
\(R_{\mathrm{Boc}}^{\mathrm{der}}(S_E)\). La identidad
\eqref{eq:bsd-suc-lt} conserva la página regular y cada página positiva. Los
isomorfismos \(j_q\) transportan esas páginas a las alturas derivadas. El
lema \ref{lem:amp-bsd-grado-pt} construye el isomorfismo de la bandera con la
red libre de Mordell--Weil y, junto con
\eqref{eq:bsd-suc-carry-orden}, prueba
\eqref{eq:amp-bsd-igualdad-grados}.

El lema \ref{lem:amp-bsd-aciclicidad-rango} deduce la aciclicidad racional
directamente del filler explícito \eqref{eq:bsd-suc-hstar}, de su reducción
lineal del complemento, de la cancelación de la raíz común y de
\(X^{\mathrm{orth}}\); de ahí se obtiene el rango completo de \(D_E\), no
como hipótesis adicional. La forma de Smith
\eqref{eq:amp-bsd-morfismo-smith} produce entonces el grupo finito \(F_E\).
El lema \ref{lem:amp-bsd-exactitud-sha-tamagawa} extrae de los triángulos
\eqref{eq:amp-bsd-triangulo-local} y
\eqref{eq:amp-bsd-triangulo-localizacion} la sucesión exacta finita; sólo
después se deducen la finitud de \(\operatorname{Sha}(E)\) y
\eqref{eq:amp-bsd-cardinal-smith}. Finalmente,
\eqref{eq:amp-bsd-regulador-nt},
\eqref{eq:amp-bsd-comparador-torsion} y
\eqref{eq:amp-bsd-comparador-periodo} publican, respectivamente, el regulador,
el cociente torsional y el período. La igualdad raíz elimina cualquier unidad
relativa entre las dos publicaciones. Al evaluar el mismo elemento basado en
ambos extremos de \eqref{eq:amp-bsd-comparador-total}, se obtiene
\eqref{eq:amp-bsd-termino-principal}.
\end{proof}

\subsection{Premisas y criterios de refutación}
\label{subsec:amp-bsd-falsadores}

La construcción parte de la unidad genealógica \(1_E\) y de la diagonal de
Alexander--Whitney. Emplea además las
publicaciones acopladas, todos los
Bocksteins y la página regular, el filler explícito \(h_*=-vf\), la
autodualidad reducida \(X^{\mathrm{orth}}\), los mapas de cadenas del
triángulo de localización y las bases orientadas de torsión y período. La equivalencia
Kato--FERS y su normalización en uno, el rango completo de \(D_E\), la
exactitud de \eqref{eq:amp-bsd-exacta-sha-tamagawa} y la finitud de
\(\operatorname{Sha}(E)\) son conclusiones de los lemas anteriores, no
premisas independientes. El criterio generador a generador del lema
\ref{lem:amp-bsd-bases-emparejadas-kato-fers} es un control no gobernante
que audita localmente el primer comparador. Los criterios que refutarían una flecha concreta
son:
\begin{enumerate}
 \item un coeficiente raíz distinto de uno, que introduciría una unidad
       relativa entre \(z_K\) y \(z_{\mathrm{PT}}\);
 \item un \(j_q\) no invertible o la omisión de una página Bockstein, que
       alteraría el grado o el determinante de altura;
 \item un divisor de Smith nulo, que impediría la finitud de \(F_E\);
 \item el fallo de exactitud de
       \eqref{eq:amp-bsd-exacta-sha-tamagawa}, que rompería la identificación
       de los factores finitos;
 \item una inversión de orientación en torsión o período, que cambiaría el
       elemento basado aunque preservase su línea no orientada.
\end{enumerate}
La proyección de Manin sin historia y el reescalamiento visible módulo tres
violan el primer criterio; por ello constituyen pruebas de necesidad para
esos modelos reducidos y no premisas del teorema
\ref{thm:amp-bsd-comparadores-construidos}.


\begin{theorem}[Comparación determinantal local--global]
\label{thm:bsd-suc-local-global}
Sean los comparadores basados ya construidos en el
\cref{thm:amp-bsd-comparadores-construidos}:
Kato--FERS/Iwasawa,
Poitou--Tate/Poincaré--Néron--Tate,
Smith--Tamagawa--Shafarevich, torsión y período.
Si \(d_{\mathrm{an}}\) es el grado de la sección analítica en la
especialización central y \(d_{\mathrm{ar}}\) el grado de la línea Selmer
basada, entonces
\begin{equation}\label{eq:bsd-suc-grados}
 d_{\mathrm{an}}=d=d_{\mathrm{ar}},
\end{equation}
y se obtiene la igualdad de elementos distinguidos
\begin{equation}\label{eq:bsd-suc-linea-final}
 \mathfrak z_{\mathrm{an}}(E)=\mathfrak z_{\mathrm{ar}}(E)
 \quad\text{en}\quad\mathcal L_E.
\end{equation}
El miembro analítico es la publicación del término inicial; el miembro
aritmético es la publicación del regulador completo, los factores
elementales de Smith,
los factores locales, la torsión y el período.
\end{theorem}

\begin{proof}
La construcción basada da
\([\Phi_{K/\mathrm{FERS}}^i](T)=I+TB_i(T)\) en cada grado. El comparador
basado de Kato--FERS expresa entonces la sección analítica corregida
como
\begin{equation}\label{eq:bsd-suc-fers}
 L_{\ell}^{\mathrm{corr}}(E,T)
 =\rho_{E,\ell}(T)\det S_E^{\mathrm{corr}}(T),
 \qquad \rho_{E,\ell}(0)=1.
\end{equation}
Por \eqref{eq:bsd-suc-orden}, su grado central es \(d\); por
\eqref{eq:bsd-suc-lt}, su término inicial es exactamente
\(R_{\mathrm{Boc}}^{\mathrm{der}}\), no ese regulador multiplicado por una
unidad indeterminada. Los isomorfismos \(j_q\) de
\eqref{eq:bsd-suc-jq} transportan cada \(\tau_q\) al determinante de la
altura derivada de su página. La página regular transporta el sumando
acíclico que las páginas positivas no ven.

La comparación de Poincaré--Néron--Tate identifica el grado de esa línea de
alturas con \(d_{\mathrm{ar}}\). Los factores elementales de Smith descomponen su parte
finita en los factores de Shafarevich y Tamagawa; el cociente de las redes
integrales aporta el cuadrado de la torsión, y el lector arquimediano aporta
el período. Todos estos mapas actúan sobre la misma línea raíz. Por el
teorema \ref{thm:bsd-suc-raiz-filler}, no queda una unidad relativa entre
las publicaciones Kato y Poitou--Tate. La composición preserva el grado y
el elemento basado, lo que prueba simultáneamente
\eqref{eq:bsd-suc-grados} y \eqref{eq:bsd-suc-linea-final}. La fórmula
escalar correspondiente se reserva para el reconocimiento final.
\end{proof}

\subsection*{Hipótesis y alcance}

La cadena BSD utiliza: el estado enriquecido completo y su coálgebra
Alexander--Whitney; el producto fibrado \eqref{eq:bsd-suc-pullback}; las dos
publicaciones construidas sobre la misma unidad; la línea raíz, su base y su
counidad; el cociente reducido perfecto y saturado, junto con
\(X^{\mathrm{orth}}\); el complejo finita--singular completo; las banderas
Bockstein finitas, sus signos de Koszul y la página regular; y las
comparaciones de altura, Smith, Tamagawa, Shafarevich, torsión y períodos
con una orientación común. La verificación grado a grado de las bases
emparejadas y de los residuales \(\varepsilon_i(b)\) queda registrada como
control local de la realización, no como premisa gobernante del cierre.
Bajo esta estructura de partida explícita,
\eqref{eq:bsd-suc-raiz}, \eqref{eq:bsd-suc-hstar},
\eqref{eq:bsd-suc-lt} y \eqref{eq:bsd-suc-linea-final} son igualdades
exactas.

La conservación de unidad, el recorrido de dos hojas y la publicación
local--global se formalizan mediante la coálgebra de cadenas, el producto
fibrado genealógico, la publicación acoplada, la reducción ortogonal y la
torre Bockstein. Las identidades auxiliares verifican por separado la propiedad de elemento grupal,
la conservación de la historia bajo \(\widehat\Pi_0\), la identidad de la
unidad en ambos canales, la involución reducida, la forma normal
finita--singular y la identidad determinantal de orden arbitrario.

\paragraph{Controles de ablación BSD.}
\noindent\textbf{Estatuto: CONTROL\_NO\_GOBERNANTE.}
Las tres ablaciones siguientes reciben las identidades ya demostradas y
falsifican únicamente modelos reducidos. No aportan premisas al propietario,
no condicionan su cierre y no poseen una arista causal de retorno hacia él.

\begin{ablation}[Proyección de Manin sin historia]
\label{abl:bsd-suc-manin}
Sustituir \(P_E^{\mathrm{gen}}\) por \(P_E^{\mathrm{Man}}\) borra datos que
la publicación acoplada necesita. En efecto, la ruta
\(\gamma=(\mathsf N\mathsf E)^{54}\) puede tener la misma celda y la misma
fase visible que la ruta trivial, pero conserva enrollamiento y memoria
distintos en \(G_{m,n}\). El aumento de Manin las identifica; la segunda
proyección de \eqref{eq:bsd-suc-pullback} no.
\end{ablation}

\begin{proof}
Todo lift basado al complejo de Manin factoriza, hasta homotopía, por la
sección del aumento \(s\varepsilon_{\mathrm{AW}}\). Por ello depende sólo
de la coordenada visible. En cambio, la segunda componente de
\(\widehat\Pi_0(c)\) es el propio \(c\), de modo que conserva el simplex de
historia y distingue el enrollamiento de \(\gamma\). La ablación refuta la
proyección desnuda; no refuta el producto fibrado empleado en el teorema.
\end{proof}

\begin{ablation}[Reescalamiento de la raíz visible sólo módulo tres]
\label{abl:bsd-suc-twist}
En la especialización \(c=1\) de
\eqref{eq:bsd-suc-diferencial-normal}, el ciclo alterado
\begin{equation}\label{eq:bsd-suc-twist4}
 z_K^{(4)}=4r+3e+b
\end{equation}
satisface
\begin{equation}\label{eq:bsd-suc-twist-defecto}
 z_{\mathrm{PT}}-z_K^{(4)}
 =\partial(-f)-3r.
\end{equation}
El defecto \(-3r\) desaparece módulo tres y reaparece módulo nueve. Por
consiguiente, \eqref{eq:bsd-suc-twist4} conserva una sombra parcial, pero no
la unidad basada ni \eqref{eq:bsd-suc-raiz}.
\end{ablation}

\begin{proof}
La igualdad \(\partial f=3e+b\) da inmediatamente
\eqref{eq:bsd-suc-twist-defecto}. Módulo tres, el segundo sumando es nulo;
módulo nueve es una clase raíz no divisible por una frontera del complemento.
Además, su coordenada raíz es cuatro, mientras
\eqref{eq:bsd-suc-normalizacion-raiz} exige uno. El ejemplo verifica la
necesidad del par acoplado y de la orientación; no introduce una excepción
al teorema que conserva ambos datos.
\end{proof}

\begin{ablation}[Pérdida de páginas del regulador]
\label{abl:bsd-suc-bockstein}
La matriz \(\operatorname{diag}(T^2,T)\) muestra que el primer Bockstein no
recupera el acarreo de orden superior. La matriz
\(\operatorname{diag}(2,T^2)\) tiene trivialización positiva
\(\tau_2=1\), pero
\(\tau_{\mathrm{reg}}=2=\operatorname{LT}_{T}(S)\). Por tanto, ni la
primera página ni el producto de páginas positivas sustituyen al regulador
completo \eqref{eq:bsd-suc-rboc}.
\end{ablation}

\begin{proof}
En el primer ejemplo, \(d=3\), \(r_0=2\) y
\(c_I^{\deg}=1\); la página uno no contiene ese último grado de memoria. En
el segundo, el único divisor singular es \(T^2\), pero la dirección regular
aporta el factor dos. Omitirla produciría uno en vez de dos. Ambos cálculos
son casos diagonales de \eqref{eq:bsd-suc-smith} y verifican la necesidad de
todas las piezas de \eqref{eq:bsd-suc-rboc}.
\end{proof}
\section*{Reconocimiento matemático final de Birch--Swinnerton--Dyer}

El desarrollo anterior concluye antes de emplear su formulación histórica
como criterio.

En el dominio declarado por el teorema
\ref{thm:bsd-suc-local-global}, la comparación de Smith identifica el
componente finito de la línea aritmética con la descomposición
Shafarevich--Tamagawa y demuestra la finitud de
\(\operatorname{Sha}(E)\). En consecuencia, su orden en la publicación
escalar siguiente es el cardinal del grupo finito identificado, no un símbolo
formal añadido a posteriori.

La traducción escalar de \eqref{eq:bsd-suc-linea-final} proporciona, con
\(r=\operatorname{rank}E(\mathbb Q)\),
\begin{equation}\label{eq:bsd-suc-rango-final}
 \operatorname{ord}_{s=1}L(E,s)=r
\end{equation}
y
\begin{equation}\label{eq:bsd-suc-formula-final}
 \frac{L^{(r)}(E,1)}{r!\,\Omega_E}
 =\frac{
   \lvert\operatorname{Sha}(E)\rvert\,\operatorname{Reg}(E)\prod_v c_v}
  {\lvert E(\mathbb Q)_{\mathrm{tors}}\rvert^2}.
\end{equation}
Esta es la formulación convencional de Birch--Swinnerton--Dyer. El orden,
el término principal y los factores aritméticos no se definen unos mediante
otros: son publicaciones coordinadas de la igualdad basada
\eqref{eq:bsd-suc-linea-final}, cuya raíz común fue fijada previamente por
\eqref{eq:bsd-suc-raiz}.

